\textbf{\RL{دانشگاه .........}}

\textbf{\RL{دانشکده ..............}}

\textbf{\RL{پایان‌نامه کارشناسی ارشد رشته ... گرایش ...}}

\RL{.............}

\textbf{\RL{نگارش}}

\textbf{\RL{نام و نام خانوادگی} \RL{دانش‌آموخته}}

\RL{}

\textbf{\RL{استاد راهنما}}

\textbf{\RL{نام و نام خانوادگی استاد راهنما}}

\textbf{\RL{استاد مشاور}}

\textbf{\RL{نام و نام خانوادگی استاد مشاور}}

\textbf{\RL{ماه و سال خورشیدی (مرداد 1404)}}


\includegraphics[width=5.18194in,height=6.93056in]{media/image1.jpeg}

\RL{\hfill\break
}

\RL{صورت جلسه}

\RL{\hfill\break
}

\RL{تقدیم}

\RL{\hfill\break
}

\RL{سپاسگزاری}

\RL{\hfill\break
}

\RL{\hfill\break
}

\textbf{\RL{فهرست مطالب}}

\textbf{\RL{\ul{عنوان صفحه}}}

\hyperref[_Toc232899850]{\RL{فصل1 : کلیات تحقیق} \hyperref[_Toc232899850]{\RL{15}}}

\hyperref[_Toc232899851]{\RL{1-1. بيان مسأله} \hyperref[_Toc232899851]{\RL{16}}}

\hyperref[_Toc232899852]{\RL{1-2. اهمیت و ضرورت تحقیق} \hyperref[_Toc232899852]{\RL{18}}}

\hyperref[_Toc232899853]{\RL{1-3. اهداف تحقيق} \hyperref[_Toc232899853]{\RL{20}}}

\hyperref[_Toc232899854]{\RL{1-3-1. هدف کلی} \hyperref[_Toc232899854]{\RL{20}}}

\hyperref[_Toc232899855]{\RL{1-3-2. اهداف جزئی} \hyperref[_Toc232899855]{\RL{20}}}

\hyperref[_Toc232899856]{\RL{1-4. فرضیه های تحقیق} \hyperref[_Toc232899856]{\RL{20}}}

\hyperref[_Toc232899857]{\RL{1-5. مفاهیم کلیدی تحقیق} \hyperref[_Toc232899857]{\RL{21}}}

\hyperref[_Toc232899858]{\RL{1-6. سهم علمی پژوهش} \hyperref[_Toc232899858]{\RL{23}}}

\hyperref[_Toc232899859]{\RL{1-7. ساختار پایان‌نامه} \hyperref[_Toc232899859]{\RL{24}}}

\hyperref[_Toc232899860]{\RL{فصل2 : مفاهیم پایه و پیشینه تحقیق} \hyperref[_Toc232899860]{\RL{25}}}

\hyperref[_Toc232899861]{\RL{2-1. مقدمه} \hyperref[_Toc232899861]{\RL{26}}}

\hyperref[_Toc232899862]{\RL{2-2. رمزنگاری} \hyperref[_Toc232899862]{\RL{27}}}

\hyperref[_Toc232899863]{\RL{2-2-1. تعریف رمزنگاری} \hyperref[_Toc232899863]{\RL{28}}}

\hyperref[_Toc232899864]{\RL{2-2-2. شیوه‌های پایه رمزنگاری} \hyperref[_Toc232899864]{\RL{30}}}

\hyperref[_Toc232899865]{\RL{2-2-3. انواع روش‌های رمزنگاری بر اساس نوع کلید} \hyperref[_Toc232899865]{\RL{31}}}

\hyperref[_Toc232899866]{\RL{2-2-4. مقایسه روش‌های رمزنگاری متقارن و نامتقارن} \hyperref[_Toc232899866]{\RL{32}}}

\hyperref[_Toc232899867]{\RL{2-2-5. ویژگی‌های یک الگوریتم رمزنگاری امن} \hyperref[_Toc232899867]{\RL{33}}}

\hyperref[_Toc232899868]{\RL{2-2-6. تحلیل رمز: انواع حملات به سامانه‌های رمزنگاری} \hyperref[_Toc232899868]{\RL{35}}}

\hyperref[_Toc232899869]{\RL{2-2-7. امنیت سیستم‌های رمزنگاری} \hyperref[_Toc232899869]{\RL{36}}}

\hyperref[_Toc232899870]{\RL{2-2-7-1- اصول ششگانه کرشهف} \hyperref[_Toc232899870]{\RL{37}}}

\hyperref[_Toc232899871]{\RL{2-2-7-2- معیارهای پنجگانه شانون} \hyperref[_Toc232899871]{\RL{38}}}

\hyperref[_Toc232899872]{\RL{2-2-8. ویژگی‌های نرم‌افزاری الگوریتم رمز} \hyperref[_Toc232899872]{\RL{39}}}

\hyperref[_Toc232899873]{\RL{2-3. رمزنگاری تصویر دیجیتال} \hyperref[_Toc232899873]{\RL{40}}}

\hyperref[_Toc232899874]{\RL{2-3-1. محدودیت‌های روش‌های سنتی رمزنگاری برای تصاویر} \hyperref[_Toc232899874]{\RL{40}}}

\hyperref[_Toc232899875]{\RL{2-3-2. ضرورت توسعه الگوریتم‌های ویژه رمزنگاری تصویر} \hyperref[_Toc232899875]{\RL{41}}}

\hyperref[_Toc232899876]{\RL{2-3-3. خصوصیات تصویر دیجیتال} \hyperref[_Toc232899876]{\RL{41}}}

\hyperref[_Toc232899877]{\RL{2-3-3-1- پیکسل در تصویر} \hyperref[_Toc232899877]{\RL{41}}}

\hyperref[_Toc232899878]{\RL{2-3-3-2- عمق رنگ} \hyperref[_Toc232899878]{\RL{41}}}

\hyperref[_Toc232899879]{\RL{2-3-3-3- ابعاد تصویر} \hyperref[_Toc232899879]{\RL{42}}}

\hyperref[_Toc232899880]{\RL{2-3-3-4- مدل‌های رنگ} \hyperref[_Toc232899880]{\RL{42}}}

\hyperref[_Toc232899881]{\RL{2-3-4. مولفه‌های مهم یک روش رمزنگاری تصویر} \hyperref[_Toc232899881]{\RL{43}}}

\hyperref[_Toc232899882]{\RL{2-4. نظریه‌ی آشوب و مفاهیم بنیادی آن} \hyperref[_Toc232899882]{\RL{44}}}

\hyperref[_Toc232899883]{\RL{2-4-1. تعریف آشوب} \hyperref[_Toc232899883]{\RL{45}}}

\hyperref[_Toc232899884]{\RL{2-4-2. خواص آشوب} \hyperref[_Toc232899884]{\RL{46}}}

\hyperref[_Toc232899885]{\RL{2-4-2-1- حساسیت نسبت به شرایط اولیه} \hyperref[_Toc232899885]{\RL{46}}}

\hyperref[_Toc232899886]{\RL{2-4-2-2- انتقال توپولوژیکی} \hyperref[_Toc232899886]{\RL{46}}}

\hyperref[_Toc232899887]{\RL{2-4-2-3- چگالی توپولوژیکی} \hyperref[_Toc232899887]{\RL{46}}}

\hyperref[_Toc232899888]{\RL{2-4-3. نظریه آشوب} \hyperref[_Toc232899888]{\RL{47}}}

\hyperref[_Toc232899889]{\RL{2-4-4. ویژگی‌های تئوری آشوب} \hyperref[_Toc232899889]{\RL{48}}}

\hyperref[_Toc232899890]{\RL{2-4-5. نمونه‌هایی از سیستم‌های آشوب} \hyperref[_Toc232899890]{\RL{51}}}

\hyperref[_Toc232899891]{\RL{2-4-5-1- نگاشت لجستیک} \hyperref[_Toc232899891]{\RL{51}}}

\hyperref[_Toc232899892]{\RL{2-4-5-2- سیستم آشوبی لورنز} \hyperref[_Toc232899892]{\RL{52}}}

\hyperref[_Toc232899893]{\RL{2-4-5-3- سیستم آشوبی چن} \hyperref[_Toc232899893]{\RL{53}}}

\hyperref[_Toc232899894]{\RL{2-5. ارتباط بین آشوب و رمزنگاری} \hyperref[_Toc232899894]{\RL{54}}}

\hyperref[_Toc232899895]{\RL{2-6. پیشینه پژوهش} \hyperref[_Toc232899895]{\RL{55}}}

2-6-1. مقایسه تطبیقی پژوهش‌های مرجع 57

\hyperref[_Toc232899896]{\RL{فصل3 : معرفی الگوریتم پیشنهادی} \hyperref[_Toc232899896]{\RL{59}}}

\hyperref[_Toc232899897]{\RL{3-1. مقدمه} \hyperref[_Toc232899897]{\RL{60}}}

\hyperref[_Toc232899898]{\RL{3-2. معرفی} \RL{سیستم}‌\RL{های} \RL{آشوبی} \RL{مورد} \RL{استفاده} \hyperref[_Toc232899898]{\RL{60}}}

\hyperref[_Toc232899899]{\RL{3-2-1. سیستم} \RL{آشوبی} \RL{سه}‌\RL{بعدی} Chen \hyperref[_Toc232899899]{\RL{60}}}

\hyperref[_Toc232899900]{\RL{3-2-2. نقشه}‌\RL{های} \RL{آشوبی} \RL{نمایی} (ECM) \hyperref[_Toc232899900]{\RL{61}}}

\hyperref[_Toc232899901]{\RL{3-3. تعریف} \RL{فضای} \RL{کلید} \hyperref[_Toc232899901]{\RL{62}}}

\hyperref[_Toc232899902]{\RL{3-4. الگوریتم} \RL{پیشنهادی: توضیح} \RL{گام}‌\RL{به}‌\RL{گام} \hyperref[_Toc232899902]{\RL{63}}}

\hyperref[_Toc232899903]{\RL{3-5. فرآیند} \RL{رمزگشایی} \hyperref[_Toc232899903]{\RL{66}}}

\hyperref[_Toc232899904]{\RL{3-6. نمای} \RL{کلی} \RL{الگوریتم} \hyperref[_Toc232899904]{\RL{66}}}

\hyperref[_Toc232899905]{\RL{3-7. روش} \RL{پیاده}‌\RL{سازی} \RL{و} \RL{ابزارهای} \RL{نرم}‌\RL{افزاری} \hyperref[_Toc232899905]{\RL{67}}}

\hyperref[_Toc232899906]{\RL{3-8. معیارهای} \RL{ارزیابی} \hyperref[_Toc232899906]{\RL{67}}}

\hyperref[_Toc232899907]{\RL{3-8-1. آنتروپی} \RL{شانون} \hyperref[_Toc232899907]{\RL{68}}}

\hyperref[_Toc232899908]{\RL{3-8-2. ضریب} \RL{همبستگی} \RL{پیکسل}‌\RL{های} \RL{مجاور} \hyperref[_Toc232899908]{\RL{68}}}

\hyperref[_Toc232899909]{\RL{3-8-3. معیارهای} NPCR \RL{و} UACI \hyperref[_Toc232899909]{\RL{68}}}

\hyperref[_Toc232899910]{\RL{3-8-4. تحلیل} \RL{زمان} \RL{اجرا} \hyperref[_Toc232899910]{\RL{69}}}

\hyperref[_Toc232899911]{\RL{3-9. جمع}‌\RL{بندی} \RL{فصل} \hyperref[_Toc232899911]{\RL{69}}}

\hyperref[_Toc232899912]{\RL{فصل4 : پیاده سازی الگوریتم پیشنهادی} \hyperref[_Toc232899912]{\RL{70}}}

\hyperref[_Toc232899913]{\RL{4-1. مقدمه} \hyperref[_Toc232899913]{\RL{71}}}

\hyperref[_Toc232899914]{\RL{4-2. تنظیمات} \RL{آزمایش} \RL{و} \RL{پارامترهای} \RL{کلید} \hyperref[_Toc232899914]{\RL{71}}}

\hyperref[_Toc232899915]{\RL{4-3. تحلیل} \RL{بصری} \RL{تصاویر} \hyperref[_Toc232899915]{\RL{71}}}

\hyperref[_Toc232899916]{\RL{4-4. تحلیل هیستوگرام} \hyperref[_Toc232899916]{\RL{73}}}

\hyperref[_Toc232899917]{\RL{4-5. تحلیل} \RL{آنتروپی} \RL{شانون} \hyperref[_Toc232899917]{\RL{74}}}

\hyperref[_Toc232899918]{\RL{4-6. تحلیل} \RL{همبستگی} \RL{پیکسل}‌\RL{های} \RL{مجاور} \hyperref[_Toc232899918]{\RL{75}}}

\hyperref[_Toc232899919]{\RL{4-7. تحلیل} \RL{حساسیت} \RL{به} \RL{کلید} NPCR \RL{و} UACI \hyperref[_Toc232899919]{\RL{76}}}

\hyperref[_Toc232899920]{\RL{4-8. تحلیل} \RL{زمان} \RL{اجرا} \hyperref[_Toc232899920]{\RL{78}}}

\hyperref[_Toc232899921]{\RL{4-9. تأیید} \RL{معکوس}‌\RL{پذیری} \RL{کامل} \hyperref[_Toc232899921]{\RL{79}}}

\hyperref[_Toc232899922]{\RL{4-10. مقایسه} \RL{با} \RL{روش}‌\RL{های} \RL{مرجع} \RL{در} \RL{ادبیات} \hyperref[_Toc232899922]{\RL{79}}}

\hyperref[_Toc232899923]{\RL{4-11. جمع}‌\RL{بندی} \RL{فصل} \hyperref[_Toc232899923]{\RL{80}}}

\hyperref[_Toc232899924]{\RL{فصل5 : بحث و نتیجه‌گیری} \hyperref[_Toc232899924]{\RL{82}}}

\hyperref[_Toc232899925]{\RL{5-1. مقدمه} \hyperref[_Toc232899925]{\RL{83}}}

\hyperref[_Toc232899926]{\RL{5-2. مرور} \RL{اهداف} \RL{پژوهش} \hyperref[_Toc232899926]{\RL{83}}}

\hyperref[_Toc232899927]{\RL{5-2-1. هدف} \RL{کلی: طراحی} \RL{و} \RL{توسعه} \RL{الگوریتم} \RL{رمزنگاری} \hyperref[_Toc232899927]{\RL{83}}}

\hyperref[_Toc232899928]{\RL{5-2-2. اهداف} \RL{جزئی} \RL{و} \RL{میزان} \RL{دستیابی} \hyperref[_Toc232899928]{\RL{84}}}

\hyperref[_Toc232899929]{\RL{5-3. یافته}‌\RL{های} \RL{اصلی} \RL{پژوهش} \hyperref[_Toc232899929]{\RL{84}}}

\hyperref[_Toc232899930]{\RL{5-3-1. معکوس}‌\RL{پذیری} \RL{کامل} \RL{الگوریتم} \hyperref[_Toc232899930]{\RL{84}}}

\hyperref[_Toc232899931]{\RL{5-3-2. آنتروپی} \RL{بالا} \RL{و} \RL{توزیع} \RL{یکنواخت} \RL{پیکسل}‌\RL{ها} \hyperref[_Toc232899931]{\RL{85}}}

\hyperref[_Toc232899932]{\RL{5-3-3. حذف} \RL{مؤثر} \RL{همبستگی} \RL{مکانی} \hyperref[_Toc232899932]{\RL{85}}}

\hyperref[_Toc232899933]{\RL{5-3-4. حساسیت} \RL{بالا} \RL{به} \RL{کلید} \hyperref[_Toc232899933]{\RL{85}}}

\hyperref[_Toc232899934]{\RL{5-3-5. فضای} \RL{کلید} \RL{بزرگ} \RL{و} \RL{مقاومت} \RL{در} \RL{برابر} \RL{جستجوی} \RL{فراگیر} \hyperref[_Toc232899934]{\RL{86}}}

\hyperref[_Toc232899935]{\RL{5-4. ارزیابی} \RL{فرضیه}‌\RL{های} \RL{پژوهش} \hyperref[_Toc232899935]{\RL{86}}}

\hyperref[_Toc232899936]{\RL{5-5. محدودیت}‌\RL{های} \RL{پژوهش} \hyperref[_Toc232899936]{\RL{86}}}

\hyperref[_Toc232899937]{\RL{5-6. پیشنهادات} \RL{برای} \RL{تحقیقات} \RL{آینده} \hyperref[_Toc232899937]{\RL{87}}}

\hyperref[_Toc232899938]{\RL{5-6-1. بهینه}‌\RL{سازی} \RL{کارایی} \RL{محاسباتی} \hyperref[_Toc232899938]{\RL{87}}}

\hyperref[_Toc232899939]{\RL{5-6-2. بهبود} \RL{مرحله} \RL{انتشار} \RL{برای} \RL{ارتقاء} UACI \hyperref[_Toc232899939]{\RL{87}}}

\hyperref[_Toc232899940]{\RL{5-6-3. گسترش} \RL{و} \RL{ارزیابی} \RL{بیشتر} \hyperref[_Toc232899940]{\RL{88}}}

\hyperref[_Toc232899941]{\RL{5-6-4. تحقیقات} \RL{نظری} \hyperref[_Toc232899941]{\RL{88}}}

\hyperref[_Toc232899942]{\RL{5-7. نتیجه}‌\RL{گیری} \RL{نهایی} \hyperref[_Toc232899942]{\RL{89}}}

\RL{\hfill\break
}

\textbf{\RL{فهرست جداول}}

\textbf{\RL{\ul{عنوان صفحه}}}

\hyperref[_Toc232899943]{\RL{جدول2-1. مقایسه رمزنگاری متقارن و نامتقارن} \hyperref[_Toc232899943]{\RL{35}}}

\hyperref[_Toc232899944]{\RL{جدول2-2. شباهت‌های سیستم‌های رمز متداول و سیستم‌های آشوبناک} \hyperref[_Toc232899944]{\RL{56}}}

جدول2-3. مقایسه تطبیقی روش‌های رمزنگاری تصویر رنگی مبتنی بر آشوب 57

\hyperref[_Toc232899945]{\RL{جدول3-1. ۹} \RL{سیستم} \RL{آشوبی} \RL{نمایی} \RL{مورد} \RL{استفاده} \RL{در} \RL{الگوریتم} \RL{پیشنهادی} \hyperref[_Toc232899945]{\RL{62}}}

\hyperref[_Toc232899946]{\RL{جدول3-2. مؤلفه}‌\RL{های} \RL{فضای} \RL{کلید} \RL{الگوریتم} \RL{پیشنهادی} \hyperref[_Toc232899946]{\RL{63}}}

\hyperref[_Toc232899947]{\RL{جدول3-3. خلاصه} \RL{مراحل} \RL{الگوریتم} \RL{رمزنگاری} \RL{پیشنهادی} \hyperref[_Toc232899947]{\RL{68}}}

\hyperref[_Toc232899948]{\RL{جدول4-1. پارامترهای} \RL{کلید} \RL{رمزنگاری} \RL{مورد} \RL{استفاده} \RL{در} \RL{آزمایش}‌\RL{ها} \hyperref[_Toc232899948]{\RL{72}}}

\hyperref[_Toc232899949]{\RL{جدول4-2. آنتروپی} \RL{شانون} \RL{قبل} \RL{و} \RL{بعد} \RL{از} \RL{رمزنگاری} \RL{(بر} \RL{حسب} \RL{بیت)} \hyperref[_Toc232899949]{\RL{75}}}

\hyperref[_Toc232899950]{\RL{جدول4-3. ضریب} \RL{همبستگی} \RL{پیکسل}‌\RL{های} \RL{مجاور} \RL{در} \RL{سه} \RL{جهت} \RL{(کانال} G\RL{)} \hyperref[_Toc232899950]{\RL{76}}}

\hyperref[_Toc232899951]{\RL{جدول4-4. نتایج} \RL{آزمون} NPCR \RL{و} UACI \hyperref[_Toc232899951]{\RL{78}}}

\hyperref[_Toc232899952]{\RL{جدول4-5. زمان} \RL{اجرای} \RL{رمزنگاری} \RL{و} \RL{رمزگشایی} \hyperref[_Toc232899952]{\RL{79}}}

\hyperref[_Toc232899953]{\RL{جدول4-6. نتایج} \RL{تأیید} \RL{معکوس}‌\RL{پذیری} \RL{الگوریتم} \hyperref[_Toc232899953]{\RL{80}}}

\hyperref[_Toc232899954]{\RL{جدول4-7. مقایسه} \RL{الگوریتم} \RL{پیشنهادی} \RL{با} \RL{روش}‌\RL{های} \RL{مرجع} \hyperref[_Toc232899954]{\RL{80}}}

\hyperref[_Toc232899955]{\RL{جدول5-1. ارزیابی} \RL{میزان} \RL{تحقق} \RL{اهداف} \RL{جزئی} \RL{پژوهش} \hyperref[_Toc232899955]{\RL{85}}}

\hyperref[_Toc232899956]{\RL{جدول5-2. ارزیابی} \RL{فرضیه}‌\RL{های} \RL{پژوهش} \RL{بر} \RL{اساس} \RL{نتایج} \RL{تجربی} \hyperref[_Toc232899956]{\RL{87}}}

\RL{\hfill\break
}

\textbf{\RL{فهرست شکل ها}}

\textbf{\RL{\ul{عنوان صفحه}}}

\hyperref[_Toc232899805]{\RL{شکل 2-1}. \RL{نگاشت لجستیک} \hyperref[_Toc232899805]{\RL{53}}}

شکل 2-2. نمای سه‌بعدی جذب‌کننده آشوبی لورنز {[}23{]}. 54

شکل 2-3. نمای سه‌بعدی جذب‌کننده آشوبی سیستم چن {[}17{]}. 55

\hyperref[_Toc232899808]{\RL{شکل 4-1. نمایش} \RL{بصری} \RL{تصاویر} \RL{اصلی،} \RL{رمزنگاری}‌\RL{شده} \RL{و} \RL{بازیابی}‌\RL{شده} \hyperref[_Toc232899808]{\RL{73}}}

\hyperref[_Toc232899809]{\RL{شکل 4-2. هیستوگرام} \RL{مقایسه}‌\RL{ای} \RL{تصاویر} \RL{اصلی} \RL{(چپ)} \RL{و} \RL{رمزنگاری}‌\RL{شده} \RL{(راست)برای} \RL{سه} \RL{کانال} \RL{رنگی} \hyperref[_Toc232899809]{\RL{74}}}

\hyperref[_Toc232899810]{\RL{شکل 4-3. نمودار} \RL{مقایسه}‌\RL{ای} \RL{آنتروپی} \RL{شانون} \RL{تصاویر} \RL{اصلی} \RL{و} \RL{رمزنگاری}‌\RL{شده} \hyperref[_Toc232899810]{\RL{75}}}

\hyperref[_Toc232899811]{\RL{شکل 4-4. نمودار} \RL{پراکندگی} \RL{همبستگی} \RL{پیکسل}‌\RL{ها} --- \RL{اصلی} \RL{(چپ)} \RL{و} \RL{رمزنگاری}‌\RL{شده} \RL{(راست)} \hyperref[_Toc232899811]{\RL{77}}}

\hyperref[_Toc232899812]{\RL{شکل 4-5. مقایسه} \RL{ضریب} \RL{همبستگی} \RL{در} \RL{سه} \RL{جهت} \RL{برای} \RL{چهار} \RL{تصویر} \hyperref[_Toc232899812]{\RL{77}}}

\hyperref[_Toc232899813]{\RL{شکل 4-6. نمودار} NPCR \RL{و} UACI \RL{برای} \RL{چهار} \RL{تصویر} \RL{در} \RL{مقایسه} \RL{با} \RL{مقادیر} \RL{ایده}‌\RL{آل} \hyperref[_Toc232899813]{\RL{78}}}

\hyperref[_Toc232899814]{\RL{شکل 4-7. نمودار} \RL{مقایسه} \RL{زمان} \RL{رمزنگاری} \RL{و} \RL{رمزگشایی} \hyperref[_Toc232899814]{\RL{79}}}

\textbf{\RL{چکیده :}}

در عصر دیجیتال، تصاویر رنگی به دلیل حجم بالا، ساختار چندکاناله و همبستگی آماری قوی میان پیکسل‌های مجاور، به الگوریتم‌های رمزنگاری تخصصی نیاز دارند که روش‌های سنتی متن‌محور مانند AES برای آن‌ها ناکارآمد است. سیستم‌های آشوبی به دلیل حساسیت شدید به شرایط اولیه و تولید دنباله‌های شبه‌تصادفی، گزینه‌ای مناسب برای رمزنگاری تصویر محسوب می‌شوند، اما سیستم‌های کلاسیک مانند نقشه Logistic و Tent در معرض پدیده «تخریب آشوب» هستند که امنیت الگوریتم را در شرایط خاص کاهش می‌دهد. در این پژوهش، الگوریتمی پنج‌مرحله‌ای برای رمزنگاری تصاویر رنگی معرفی شده است که از ترکیب سیستم آشوبی سه‌بعدی Chen و نه سیستم آشوبی نمایی (ECM) بهره می‌برد. در مرحله اول، شرایط اولیه با دقت ۱۵ رقم اعشاری به عنوان کلید رمزنگاری تعریف می‌شوند. در مرحله دوم، پیکسل‌های تصویر با استفاده از دنباله آشوبی حاصل از حل عددی RK45 سیستم Chen با پارامترهای a=35، b=3 و c=28 به‌صورت تصادفی جابه‌جا می‌شوند. در مرحله سوم، برای هر پیکسل به‌صورت مستقل یکی از نه سیستم نمایی شامل LEL، LES، LET، SEL، SES، SET، TEL، TES و TET از طریق یک نقشه Logistic مستقل انتخاب می‌شود. در مرحله چهارم، مقادیر رنگی پیکسل‌ها با عملیات XOR با بایت کلید تولیدشده از سیستم نمایی انتخابی ترکیب می‌شوند. در مرحله پنجم، یک لایه XOR مستقل اضافی با یک Logistic مستقل، عمق رمزنگاری را دوچندان می‌کند. الگوریتم پیشنهادی بر روی چهار تصویر استاندارد از مجموعه داده USC-SIPI در ابعاد ۵۱۲×۵۱۲ و ۲۵۶×۲۵۶ ارزیابی شد. نتایج نشان داد که میانگین آنتروپی شانون تصاویر رمزنگاری‌شده به ۷.۹۹۹۳ بیت (از ۸ بیت ایده‌آل) رسیده است. ضریب همبستگی پیکسل‌های مجاور از میانگین ‎۰.۹۹۸۹‎ در تصاویر اصلی به ‎۰.۰۱۲۶‎ در تصاویر رمزنگاری‌شده کاهش یافت که معادل کاهش ‎\lr{۹۸.۷٪}‎ است. نرخ تغییر پیکسل (NPCR) برای تصویر Peppers به \lr{۹۹.۴۸٪} رسید. معکوس‌پذیری کامل الگوریتم با MSE = 0 و PSNR = $\infty$ برای تمام تصاویر تأیید شد. فضای کلید الگوریتم از ۱۰¹⁰⁵ تجاوز می‌کند که حمله جستجوی فراگیر را به‌لحاظ محاسباتی ناممکن می‌سازد. نوآوری اصلی این پژوهش، انتخاب پویا و مستقل سیستم آشوبی نمایی برای هر پیکسل است که یک لایه امنیتی منحصربه‌فرد ایجاد می‌کند؛ در صورتی که مهاجم حتی موفق به شناسایی دنباله رمزنگاری یک پیکسل شود، هیچ اطلاعاتی درباره پیکسل‌های دیگر به‌دست نمی‌آورد. مقایسه با روش‌های مرجع نشان می‌دهد الگوریتم پیشنهادی از نظر آنتروپی و NPCR رقابت‌پذیر است و با بهینه‌سازی پیاده‌سازی از طریق برداری‌سازی و پردازش موازی GPU، پتانسیل دستیابی به سرعت‌های بلادرنگ را داراست.

\RL{\textbf{کلمات کلیدی:} رمزنگاری تصویر؛ سیستم آشوبی نمایی؛ سیستم آشوبی} Chen\RL{؛ جایگشت پیکسل؛ آنتروپی شانون؛} NPCR\RL{؛} UACI

\RL{\hfill\break
}

\begin{enumerate}
\def\labelenumi{\arabic{enumi}.}
\item
  \protect\phantomsection\label{_Toc232899850}{}\RL{کلیات تحقیق}
\end{enumerate}

\RL{\hfill\break
}

\begin{enumerate}
\def\labelenumi{\arabic{enumi}.}
\item
  \protect\phantomsection\label{_Toc232899851}{}\RL{بيان مسأله}
\end{enumerate}

\RL{در عصر دیجیتال، تبادل اطلاعات به‌ویژه در قالب تصاویر، به بخشی جدایی‌ناپذیر از ارتباطات، تجارت، سیستم‌های نظارتی، و رسانه‌های اجتماعی تبدیل شده است. در این میان، تصاویر دیجیتال به دلیل نقش کلیدی در انتقال داده‌های بصری، به یکی از مهم‌ترین انواع اطلاعات چندرسانه‌ای بدل شده‌اند که حجم بالایی از داده را در خود جای می‌دهند و کاربرد گسترده‌ای در شبکه‌های ارتباطی، سامانه‌های ذخیره‌سازی و تحلیل داده دارند {[}1{]}.}

\RL{با گسترش دسترسی به فضای مجازی و تهدیدهای فزاینده در حوزه امنیت سایبری، حفاظت از حریم خصوصی اطلاعات تصویری به یکی از دغدغه‌های اصلی پژوهشگران و توسعه‌دهندگان امنیت اطلاعات تبدیل شده است. برخلاف داده‌های متنی که رمزنگاری آن‌ها با الگوریتم‌های استانداردی نظیر} AES \RL{و} RSA \RL{به‌خوبی قابل انجام است، تصاویر دیجیتال به دلیل ویژگی‌هایی همچون حجم زیاد، ساختار دوبعدی، همبستگی آماری بالا میان پیکسل‌های مجاور، و حساسیت بالا به تغییرات، به الگوریتم‌های رمزنگاری}\footnote{Encryption} \RL{تخصصی‌تر و پیچیده‌تری نیاز دارند {[}2و3{]}.}

\RL{مهم‌ترین چالش‌ها در رمزنگاری تصاویر دیجیتال را می‌توان به شرح زیر برشمرد:}

\RL{حفظ ویژگی‌های آماری تصویر پس از رمزنگاری؛}

\begin{itemize}
\item
  \RL{ایجاد مقاومت در برابر انواع حملات نظیر تحلیل آماری، حملات تمایزی و جستجوی فراگیر}\footnote{Brute Force} \RL{؛}
\item
  \RL{کاهش بار محاسباتی جهت کاربرد در سامانه‌های بلادرنگ؛}
\item
  \RL{پیشگیری از تخریب داده‌های تصویری طی فرآیند رمزنگاری و رمزگشایی}\footnote{Decryption}\RL{.}
\end{itemize}

\RL{در پاسخ به این چالش‌ها، بهره‌گیری از سیستم‌های آشوبی}\footnote{Chaotic Systems} \RL{به‌عنوان راهکاری مؤثر مورد توجه قرار گرفته است. سیستم‌های آشوبی با رفتار غیرخطی، حساسیت شدید به شرایط اولیه، و توانایی تولید دنباله‌های شبه‌تصادفی، امکان طراحی الگوریتم‌هایی با امنیت بالا و کلیدهای غیرقابل پیش‌بینی را فراهم می‌آورند {[}4{]}. فرآیندهای جایگشت و پراکندگی پیکسل‌ها با کمک این سیستم‌ها، ساختار آماری تصویر را به‌شدت تغییر داده و مانع بازیابی تصویر اصلی بدون کلید می‌شوند.}

\RL{با این حال، یکی از نقاط ضعف رایج در الگوریتم‌های مبتنی بر سیستم‌های آشوبی کلاسیک، بروز پدیده‌ای به‌نام «تخریب آشوب»}\footnote{Chaotic Degradation} \RL{است. این پدیده در شرایط خاص یا در اثر خطای عددی می‌تواند باعث از دست رفتن ویژگی‌های آشوبی شود و در نهایت، امنیت سیستم را کاهش دهد. این اشکال به‌ویژه در سیستم‌هایی نظیر نقشه} Logistic \RL{و} Tent \RL{به‌صورت برجسته‌تری مشاهده شده است {[}5{]}.}

\RL{سیستم آشوبی به‌طور کلی، نوعی سیستم دینامیکی غیرخطی است که رفتار آن در طول زمان به شدت به شرایط اولیه وابسته است و این وابستگی منجر به تولید پاسخ‌هایی کاملاً متفاوت در صورت تغییرات بسیار جزئی در مقدارهای ورودی می‌شود. به بیان ساده‌تر، سیستم آشوبی با وجود پیروی از معادلات مشخص، خروجی‌ای شبه‌تصادفی و غیرقابل پیش‌بینی تولید می‌کند، به‌گونه‌ای که برای ناظر بیرونی همانند نویز یا تصادف به نظر می‌رسد، اما در واقع دارای نظم در بی‌نظمی است.}

\RL{ویژگی‌های اصلی سیستم‌های آشوبی شامل:}

\begin{itemize}
\item
  \RL{حساسیت شدید به شرایط اولیه}\footnote{Sensitive Dependence on Initial Conditions}
\item
  \RL{غیرخطی بودن روابط سیستم} \footnote{Nonlinearity}
\item
  \RL{رفتار شبه‌تصادفی}\footnote{Pseudo-randomness}
\item
  \RL{وابستگی به پارامترهای سیستمی}
\end{itemize}

\RL{این ویژگی‌ها باعث می‌شوند که سیستم‌های آشوبی گزینه‌ای بسیار مناسب برای تولید کلیدهای رمزنگاری، الگوریتم‌های پراکندگی و اختفاء داده‌ها باشند. به‌عنوان مثال، نگاشت‌های ساده‌ای مانند} Logistic\RL{،} Tent\RL{،} Henon\RL{،} Chen \RL{و} Lorenz \RL{به‌طور گسترده در رمزنگاری تصویر به‌کار رفته‌اند. سیستم لورنز}\footnote{Lorenz System} \RL{یکی از مشهورترین مدل‌های آشوبی سه‌بعدی است که از معادلات دیفرانسیل غیرخطی تشکیل شده و در ابتدا برای مدل‌سازی حرکت سیالات جوی طراحی شد، اما به‌دلیل توانایی بالا در تولید مسیرهای آشوبی، امروزه در حوزه رمزنگاری و ارتباطات امن نیز مورد استفاده قرار گرفته است {[}6{]}.}

\RL{از سوی دیگر، پژوهشگران برای افزایش مقاومت در برابر حملات تحلیل همبستگی و جلوگیری از تخریب آشوب، سیستم‌هایی تحت عنوان نقشه‌های آشوب نمایی}\footnote{Exponential Chaotic Maps} \RL{را توسعه داده‌اند. این سیستم‌ها، با استفاده از نگاشت‌های ترکیبی و تغییرات پارامتری گسترده، دامنه آشوبی پایدارتری نسبت به نگاشت‌های کلاسیک دارند و در نتیجه، در طراحی الگوریتم‌های رمزنگاری مقاوم بسیار مؤثر واقع شده‌اند {[}4{]}. به‌عنوان نمونه، در الگوریتمی که توسط خدادادی و زندوکیلی {[}7{]} ارائه شد، ترکیب دو سیستم} Chen \RL{و} Logistic \RL{برای رمزنگاری تصویر پیشنهاد گردید؛ با این‌حال، این روش نیز در شرایط خاص با کاهش پایداری و تخریب ویژگی آشوب مواجه بود .}

\RL{در پژوهش حاضر، با هدف ارتقاء سطح امنیت و کارایی الگوریتم‌های رمزنگاری تصویر، روشی جدید معرفی می‌شود که از ترکیب نه سیستم آشوبی نمایی به‌همراه الگوریتم جایگشت} Chen \RL{بهره می‌برد. سیستم‌های نمایی، مانند} LEL\RL{،} SEL \RL{و سایر نگاشت‌های آشوب نمائی، به‌عنوان نسل جدید سیستم‌های آشوبی، قادرند رفتارهای پیچیده‌تر و توزیع آماری یکنواخت‌تری نسبت به مدل‌های کلاسیک ایجاد کنند {[}4{]}.}

\RL{در این الگوریتم، ابتدا پیکسل‌های تصویر با استفاده از دنباله‌های تولیدشده توسط سیستم} Chen \RL{به‌صورت تصادفی جابه‌جا می‌شوند. سپس، رمزنگاری هر پیکسل با انتخاب تصادفی یکی از ۹ سیستم نمایی صورت می‌گیرد. این رویکرد چندلایه، علاوه بر ارتقاء پیچیدگی، موجب افزایش آنتروپی اطلاعاتی، کاهش همبستگی پیکسل‌ها و افزایش مقاومت در برابر حملات تحلیل آماری می‌شود. از طرفی، با بهره‌گیری از توزیع‌های آماری غیرخطی‌تر و پایدارتر، احتمال تخریب ویژگی‌های آشوبی به‌طور مؤثری کاهش می‌یابد.}

\RL{در نهایت، هدف از این تحقیق، توسعه یک الگوریتم رمزنگاری مقاوم و کارآمد است که با تکیه بر سیستم‌های آشوبی نوین، توازنی بهینه میان امنیت بالا و سرعت اجرای مناسب برقرار کرده و پاسخگوی نیازهای کاربردی در حوزه‌هایی همچون پزشکی، امنیت نظامی، شبکه‌های اجتماعی و ذخیره‌سازی ابری باشد.}

\begin{enumerate}
\def\labelenumi{\arabic{enumi}.}
\item
  \protect\phantomsection\label{_Toc232899852}{}\RL{اهمیت و ضرورت تحقیق}
\end{enumerate}

\RL{با گسترش فناوری‌های دیجیتال و افزایش وابستگی جوامع به سامانه‌های الکترونیکی، امنیت اطلاعات به یکی از دغدغه‌های اصلی در حوزه فناوری اطلاعات تبدیل شده است. تصاویر دیجیتال، به‌ویژه تصاویر رنگی، به‌دلیل حجم بالای اطلاعات و کاربرد گسترده در حوزه‌هایی مانند ارتبا.طات، شبکه‌های اجتماعی، نظارت تصویری، پزشکی و نظامی، بیش از گذشته در معرض تهدیدات امنیتی نظیر شنود، جعل و دستکاری قرار دارند {[}10{]}. اهمیت حفاظت از این داده‌های بصری، ضرورت توسعه روش‌های نوین رمزنگاری را دوچندان کرده است}.

\RL{یکی از چالش‌های اساسی در رمزنگاری تصاویر رنگی، ساختار پیچیده و چندبعدی آن‌هاست. تصاویر رنگی به‌صورت ترکیبی از سه کانال} RGB \RL{تعریف می‌شوند که وابستگی بین کانال‌ها، کار رمزنگاری ایمن و بهینه را دشوار می‌سازد. این وابستگی‌ها می‌توانند توسط مهاجمان برای تحلیل ساختاری تصویر مورد سوءاستفاده قرار گیرند {[}3{]}. بنابراین، طراحی الگوریتم‌هایی که بتوانند این وابستگی‌ها را کاهش داده و توزیع داده‌ها را تصادفی‌تر سازند، امری حیاتی به نظر می‌رسد}.

\RL{مطالعات متعددی در حوزه رمزنگاری آشوبی انجام شده‌اند که نشان می‌دهند سیستم‌های آشوبی به‌دلیل ویژگی‌هایی نظیر حساسیت بالا به شرایط اولیه، رفتار غیرخطی و قابلیت تولید دنباله‌های شبه‌تصادفی پیچیده، ظرفیت بالایی در تولید کلیدهای رمزنگاری غیرقابل پیش‌بینی دارند. با این حال، اغلب روش‌های موجود تنها از یک سیستم آشوبی منفرد استفاده می‌کنند که در برابر حملاتی مانند تحلیل همبستگی و جستجوی فراگیر آسیب‌پذیر هستند {[}12{]}. به‌منظور مقابله با این آسیب‌پذیری‌ها، استفاده از ترکیب چند سیستم آشوبی نمایی به‌صورت تصادفی، می‌تواند پیچیدگی رمزنگاری را افزایش داده و مقاومت آن را در برابر انواع حملات ارتقا بخشد}.

\RL{بر اساس یافته‌های اخیر، سیستم‌های آشوبی نمایی نسبت به سیستم‌های آشوبی کلاسیک مانند} Logistic \RL{و} Tent Map\RL{، توزیع احتمال پیچیده‌تری تولید می‌کنند و پایداری بیشتری در برابر تغییرات پارامترها دارند. این ویژگی‌ها آن‌ها را به گزینه‌ای مناسب برای پیاده‌سازی رمزنگاری‌های مقاوم تبدیل کرده است {[}5{]}. از طرف دیگر، استفاده از سیستم‌هایی نظیر سیستم آشوبی} Chen\RL{، با ویژگی‌هایی همچون سرعت پردازش بالا و سادگی پیاده‌سازی، می‌توانند به افزایش کارایی الگوریتم‌های رمزنگاری کمک کنند}.

\RL{افزون بر این، بسیاری از کاربردهای حیاتی نظیر انتقال تصاویر پزشکی، داده‌های مالی در بانکداری، یا اطلاعات طبقه‌بندی‌شده در صنایع نظامی، نیازمند الگوریتم‌هایی هستند که هم امنیت بسیار بالایی داشته باشند و هم بتوانند در زمان محدود پردازش انجام دهند؛ این مسأله، لزوم دستیابی به توازنی منطقی میان امنیت و کارایی را پررنگ‌تر می‌کند {[}5{]}. در همین راستا، پژوهش حاضر با ترکیب ۹ سیستم آشوبی نمایی و استفاده از سیستم آشوبی} Chen \RL{جهت جایگشت، در پی توسعه مدلی نوین است که بتواند این توازن را به‌صورت مؤثر محقق سازد}.

\RL{در نهایت، رمزنگاری آشوبی به‌عنوان یکی از حوزه‌های نوظهور در علوم رمزنگاری، همچنان نیازمند تحقیقات عمیق‌تری برای بهره‌برداری کامل از ظرفیت‌های خود است. مطالعات نظری و تجربی در این زمینه می‌توانند زمینه‌ساز توسعه مدل‌های امن‌تر و بهینه‌تری برای محافظت از داده‌های حساس در عصر دیجیتال باشند. پژوهش حاضر، با تمرکز بر طراحی و ارزیابی یک روش ترکیبی پیشرفته، می‌کوشد گامی مؤثر در مسیر توسعه دانش رمزنگاری آشوبی بردارد.}

\begin{enumerate}
\def\labelenumi{\arabic{enumi}.}
\setcounter{enumi}{1}
\item
  \protect\phantomsection\label{_Toc232899853}{}\RL{اهداف تحقيق}

  \begin{enumerate}
  \def\labelenumii{\arabic{enumii}.}
  \item
    \protect\phantomsection\label{_Toc232899854}{}\RL{هدف کلی}
  \end{enumerate}
\end{enumerate}

\RL{طراحی و توسعه‌ی الگوریتمی برای رمزنگاری تصاویر رنگی با بهره‌گیری از ترکیب ۹ سیستم آشوبی نمایی، با هدف ارتقاء امنیت اطلاعات تصویری و بهبود کارایی محاسباتی.}

\begin{enumerate}
\def\labelenumi{\arabic{enumi}.}
\setcounter{enumi}{1}
\item
  \protect\phantomsection\label{_Toc232899855}{}\RL{اهداف جزئی}
\end{enumerate}

\begin{itemize}
\item
  \RL{بررسی عملکرد ۹ سیستم آشوبی نمایی در تولید دنباله‌های تصادفی پیچیده و غیرقابل پیش‌بینی.}
\item
  \RL{بهبود مقاومت الگوریتم رمزنگاری در برابر حملات آماری و جستجوی فراگیر از طریق ترکیب چندین سیستم آشوبی.}
\item
  \RL{افزایش پیچیدگی فرآیند رمزنگاری با استفاده از انتخاب تصادفی سیستم آشوبی برای رمزنگاری هر پیکسل.}
\item
  \RL{مقایسه امنیت و کارایی الگوریتم پیشنهادی با سایر روش‌های رمزنگاری مبتنی بر آشوب و ارزیابی میزان برتری آن.}

  \begin{enumerate}
  \def\labelenumi{\arabic{enumi}.}
  \item
    \protect\phantomsection\label{_Toc232899856}{}\RL{فرضیه های تحقیق}
  \end{enumerate}
\end{itemize}

\RL{فرضیه ۱: ترکیب ۹ سیستم نمایی منجر به افزایش امنیت در برابر حملات آماری می شودفرضیه ۲: جایگشت} Chen \RL{سبب کاهش همبستگی بین پیکسل های مجاور می شودفرضیه ۳: انتخاب تصادفی سیستم نمایی منجر به آنتروپی بالاتر می شود}

\RL{فرضیه ۴: الگوریتم از نظر زمان اجرا در محدوده قابل قبول قرار داردفرضیه ۵: توزیع آماری تصادفی تر نسبت به الگوریتم های پیشین}

\begin{enumerate}
\def\labelenumi{\arabic{enumi}.}
\setcounter{enumi}{1}
\item
  \protect\phantomsection\label{_Toc232899857}{}\RL{مفاهیم کلیدی تحقیق}
\end{enumerate}

\textbf{\RL{رمزنگاری}}\footnote{Encryption}\textbf{\RL{:}}\\
\RL{فرایندی ریاضیاتی است که طی آن، داده‌های اصلی (تصویر، متن، صدا و ...) به فرم غیرقابل فهم و ایمن تبدیل می‌شوند، به‌گونه‌ای که تنها افراد مجاز با استفاده از کلید یا الگوریتم رمزگشایی می‌توانند به اطلاعات اصلی دسترسی پیدا کنند. هدف اصلی رمزنگاری، حفظ محرمانگی، صحت و تمامیت اطلاعات در برابر دسترسی یا دستکاری غیرمجاز است}.

\textbf{\RL{رمزنگاری تصویر}}\footnote{Image Encryption} \textbf{\RL{:}}

\RL{زیرشاخه‌ای از رمزنگاری است که تمرکز آن بر رمزنگاری داده‌های تصویری به‌ویژه تصاویر دیجیتال رنگی یا خاکستری است. رمزنگاری تصویر باید با در نظر گرفتن ویژگی‌های خاص تصویر نظیر حجم بالا، وابستگی بین پیکسل‌ها، ساختار دو یا سه‌بعدی و حساسیت به تغییرات طراحی شود}.

\textbf{\RL{تصاویر رنگی}}\footnote{Color Images}\textbf{:}

\RL{تصاویری هستند که هر پیکسل آن‌ها از ترکیب سه کانال اصلی رنگی قرمز} (R)\RL{، سبز} (G) \RL{و آبی} (B) \RL{تشکیل شده است. این کانال‌ها مقادیر عددی خاصی دارند (معمولاً در بازه ۰ تا ۲۵۵) که در مجموع رنگ نهایی پیکسل را مشخص می‌کنند. رمزنگاری این نوع تصاویر نیازمند رویکردهای پیچیده‌تری نسبت به تصاویر خاکستری است، زیرا اطلاعات بین کانال‌ها وابسته و همبسته هستند}.

\textbf{\RL{سیستم آشوبی}}\footnote{Chaotic System}\textbf{\RL{:}}

\RL{سیستم‌های دینامیکی غیرخطی هستند که رفتار آن‌ها در طول زمان به شدت به شرایط اولیه وابسته است. این سیستم‌ها با وجود آن‌که معادلات دقیق و قطعی دارند، نتایجی شبه‌تصادفی، غیرقابل پیش‌بینی و پیچیده تولید می‌کنند. ویژگی‌هایی مانند حساسیت به مقدار اولیه، غیرخطی‌بودن و قابلیت تولید دنباله‌های پیچیده، آن‌ها را برای تولید کلیدها و ایجاد اختفاء در رمزنگاری بسیار مناسب ساخته است}.

\textbf{\RL{نگاشت}} \RL{\textbf{آشوب نمایی:}}

\RL{دسته‌ای پیشرفته از سیستم‌های آشوبی هستند که با ترکیب نگاشت‌های ریاضی ساده مانند} Logistic\RL{،} Sine \RL{و} Tent \RL{و افزودن مولفه‌های نمایی، رفتارهای غیرخطی و پویاتری را ایجاد می‌کنند. این نگاشت‌ها دارای دامنه آشوب وسیع‌تر، پایداری بیشتر نسبت به خطا و قابلیت تولید دنباله‌های رمزنگاری غیرقابل پیش‌بینی هستند. نمونه‌هایی از این سیستم‌ها عبارتند از} LEL\RL{،} SEL\RL{،} LET \RL{و} ...

\textbf{\RL{جایگشت}}\footnote{Permutation} \textbf{:}\\
\RL{یکی از مراحل اصلی در رمزنگاری تصویر است که طی آن موقعیت پیکسل‌های تصویر به‌صورت غیرقابل پیش‌بینی تغییر داده می‌شود، بدون آن‌که مقادیر رنگی آن‌ها تغییر یابد. هدف این مرحله، از بین بردن ساختار و همبستگی مکانی بین پیکسل‌ها است تا تحلیل آماری برای مهاجم دشوار شود}.

\textbf{\RL{انتشار}}\footnote{Diffusion}\textbf{:}\\
\RL{مرحله‌ای از رمزنگاری که در آن مقادیر رنگی پیکسل‌ها تغییر داده می‌شوند. به عبارتی، اطلاعات تصویری با مقادیر تصادفی ترکیب می‌شود تا رابطه‌ای میان داده‌های ورودی و خروجی باقی نماند. این مرحله باعث افزایش مقاومت در برابر تحلیل تفاضلی می‌شود}.

\textbf{\RL{آنتروپی شانون}} \footnote{Shannon Entropy}\textbf{:}

\RL{مقیاسی از میزان تصادفی‌بودن یا بی‌نظمی در توزیع داده‌هاست. در رمزنگاری تصویر، هر چه مقدار آنتروپی به مقدار ایده‌آل (برای تصاویر ۸ بیتی، مقدار ۸) نزدیک‌تر باشد، نشان‌دهنده آن است که الگوریتم رمزنگاری به‌خوبی توانسته تصویر را غیرقابل تحلیل کند}.

\textbf{NPCR \RL{و} UACI\RL{:}}

\RL{دو معیار مهم برای ارزیابی حساسیت الگوریتم رمزنگاری نسبت به تغییرات جزئی در کلید یا تصویر اولیه هستند}.

\begin{itemize}
\item
  \textbf{NPCR}\footnote{Number of Pixel Change Rate} \RL{درصد تغییرات پیکسل‌ها بین دو تصویر رمزنگاری‌شده با کلیدهای بسیار نزدیک}.
\item
  \textbf{UACI} \footnote{Unified Average Changing Intensity} \RL{میانگین شدت تغییرات بین پیکسل‌های دو تصویر رمزنگاری‌شده}.

  \begin{enumerate}
  \def\labelenumi{\arabic{enumi}.}
  \item
    \protect\phantomsection\label{_Toc232899858}{}\RL{سهم علمی پژوهش}
  \end{enumerate}
\end{itemize}

\RL{در این پایان‌نامه، یک چارچوب نوین برای رمزنگاری تصاویر رنگی مبتنی بر سیستم‌های آشوبی نمایی ارائه شده است که با هدف افزایش امنیت، پیچیدگی دینامیکی و معکوس‌پذیری کامل طراحی گردیده است. مهم‌ترین سهم‌های علمی این پژوهش به شرح زیر می‌باشند}:

\begin{enumerate}
\def\labelenumi{\arabic{enumi}.}
\item
  \RL{ارائه یک الگوریتم رمزنگاری تصویر مبتنی بر مدل آشوبی نمایی}\footnote{Exponential Chaotic Model} \RL{که با ترکیب غیرخطی نگاشت‌های پایه، دنباله‌های آشوبی با رفتار پیچیده‌تر و حساسیت بالاتر نسبت به شرایط اولیه تولید می‌کند}.
\item
  \RL{استفاده هم‌زمان از چند سیستم آشوبی نمایی (۹ نگاشت متفاوت) و انتخاب پویاى آن‌ها در سطح پیکسل، که منجر به افزایش فضای کلید، کاهش الگوپذیری و ارتقای مقاومت الگوریتم در برابر حملات آماری و دیفرانسیلی می‌شود}.
\item
  \RL{طراحی یک سازوکار انتشار اصلاح‌شده و کاملاً معکوس‌پذیر که در آن دنباله آشوبی صرفاً به کلید و اندیس پیکسل وابسته بوده و از وابستگی مستقیم به مقدار پیکسل جلوگیری شده است. این ویژگی تضمین می‌کند که تصویر رمزگشایی‌شده به‌صورت صددرصد با تصویر اصلی منطبق باشد}.
\item
  \RL{بهره‌گیری از سیستم آشوبی پیوسته} Chen \RL{برای تولید جریان کلید، که در مقایسه با نگاشت‌های گسسته، کیفیت آشوب، حساسیت به شرایط اولیه و امنیت کلید را به‌طور قابل‌توجهی افزایش می‌دهد}.
\item
  \RL{انجام مجموعه‌ای جامع از تحلیل‌های امنیتی شامل آنتروپی، همبستگی پیکسل‌های مجاور،} NPCR\RL{،} UACI \RL{و تحلیل هیستوگرام، که نشان‌دهنده عملکرد مناسب الگوریتم پیشنهادی در مقایسه با روش‌های موجود می‌باشد}.
\item
  \RL{پیاده‌سازی کامل الگوریتم و ارزیابی تجربی آن بر روی تصاویر رنگی استاندارد با ابعاد مختلف، که کارایی عملی و قابلیت کاربرد الگوریتم در سناریوهای واقعی را تأیید می‌کند.}

  \begin{enumerate}
  \def\labelenumii{\arabic{enumii}.}
  \item
    \protect\phantomsection\label{_Toc232899859}{}\RL{ساختار پایان‌نامه}
  \end{enumerate}
\end{enumerate}

\RL{ساختار این پایان‌نامه به‌صورت زیر سازمان‌دهی شده است}:

\RL{در فصل دوم، مفاهیم پایه و پیش‌نیازهای نظری مرتبط با رمزنگاری تصاویر، سیستم‌های آشوبی، ویژگی‌های نگاشت‌های آشوبی و معیارهای ارزیابی امنیت مرور می‌شود. همچنین برخی از مهم‌ترین پژوهش‌های انجام‌شده در این حوزه مورد بررسی قرار می‌گیرند}.

\RL{فصل سوم به معرفی الگوریتم پیشنهادی اختصاص دارد. در این فصل، معماری کلی روش ارائه‌شده، فرآیند تولید کلید، مرحله جایگشت و مرحله انتشار به‌طور دقیق تشریح شده و روند رمزنگاری و رمزگشایی به‌صورت گام‌به‌گام توضیح داده می‌شود}.

\RL{در فصل چهارم، پیاده‌سازی عملی الگوریتم و نتایج آزمایش‌های تجربی ارائه می‌گردد. در این فصل، عملکرد الگوریتم پیشنهادی از نظر معیارهای امنیتی و آماری تحلیل شده و نتایج به‌صورت جداول و نمودارهای تحلیلی گزارش می‌شوند}.

\RL{در فصل پنجم، نتایج به‌دست‌آمده مورد بحث و بررسی قرار گرفته و عملکرد روش پیشنهادی با برخی از روش‌های مطرح در ادبیات پژوهش مقایسه می‌شود. در پایان این فصل، جمع‌بندی نهایی پژوهش ارائه شده و پیشنهادهایی برای تحقیقات آینده مطرح می‌گردد.}

\begin{enumerate}
\def\labelenumi{\arabic{enumi}.}
\setcounter{enumi}{1}
\item
  \protect\phantomsection\label{_Toc232899860}{}\RL{مفاهیم پایه و پیشینه تحقیق}
\end{enumerate}

\RL{\hfill\break
}

\begin{enumerate}
\def\labelenumi{\arabic{enumi}.}
\item
  \protect\phantomsection\label{_Toc232899861}{}\RL{مقدمه}
\end{enumerate}

\RL{از زمان انتشار مشاهدات ادوارد لورنز در زمینه‌ی رفتارهای آشوبناک، مفهوم «آشوب» به یکی از موضوعات مورد توجه در جامعه‌ی علمی تبدیل شد. نظریه‌ی آشوب شاخه‌ای از ریاضیات است که به تحلیل سیستم‌هایی می‌پردازد که رفتار آن‌ها به تغییرات کوچک در شرایط اولیه بسیار حساس است. به بیان دیگر، در این‌گونه سامانه‌ها، تغییرات ظاهراً ناچیز می‌تواند منجر به تفاوت‌های چشمگیر در خروجی شود}.

\RL{پیش از ظهور این نظریه، بسیاری از پدیده‌های تصادفی یا غیرقابل‌پیش‌بینی، فاقد توجیه علمی دقیق تلقی می‌شدند، اما نظریه‌ی آشوب توانست با اتکا بر اصول ریاضی، فیزیک و فلسفه، مبانی تبیینی منسجمی برای این پدیده‌ها فراهم آورد. قانون علت و معلول همچنان در این نظریه برقرار است، اما تعداد عوامل مؤثر بر نتایج به‌طور چشمگیری افزایش می‌یابد. در واقع، نظریه‌ی آشوب ابزاری فراهم می‌کند تا بتوان پدیده‌های پیچیده‌ی طبیعی و واقعی را بهتر درک و تحلیل کرد}.

\RL{نکته‌ی مهم آن است که نظریه‌ی آشوب صرفاً به‌منزله‌ی روشی برای ساده‌سازی پدیده‌های فیزیکی در نظر گرفته نمی‌شود، بلکه مدل‌های ریاضی ناشی از آن، به‌عنوان منبعی از اطلاعات جدید و غیرخطی محسوب می‌شوند. این مدل‌ها می‌توانند نه تنها رفتار سامانه‌های واقعی را پیش‌بینی کنند، بلکه در حوزه‌های دیگر از جمله علوم مهندسی، فیزیک، و رمزنگاری اطلاعات نیز به کار گرفته شوند}.

\RL{در سال‌های اخیر، توجه ویژه‌ای به کاربرد نظریه‌ی آشوب در علم ارتباطات و امنیت اطلاعات شده است. ویژگی‌هایی چون حساسیت بالا به شرایط اولیه، غیرقابل‌پیش‌بینی بودن و رفتار شبه‌تصادفی، سبب شده تا سیستم‌های آشوبی به‌عنوان بستری مناسب برای طراحی الگوریتم‌های رمزنگاری مطرح شوند}.

\RL{در واقع، میان ویژگی‌های ذاتی سامانه‌های آشوبناک و خواص مطلوب در یک سامانه‌ی رمزنگاری امن (نظیر پخش و درهم‌سازی قوی، وابستگی شدید به کلید، و مقاومت در برابر حملات آماری) شباهت‌های قابل‌توجهی وجود دارد. با وجود این ارتباط نزدیک، طراحی الگوریتم‌های رمزنگاری مؤثر بر پایه‌ی سیستم‌های آشوبی نیازمند دقت و دانش عمیق ریاضی است؛ چرا که هرگونه خطا در پیاده‌سازی یا انتخاب پارامترها می‌تواند امنیت سامانه را به‌شدت کاهش دهد}.

\RL{بر این اساس، در این فصل به معرفی مبانی نظری مرتبط با موضوع پژوهش پرداخته می‌شود. ابتدا مفاهیم بنیادین رمزنگاری و انواع روش‌های آن مرور می‌گردد، سپس ویژگی‌های اصلی تصاویر دیجیتال، نظریه‌ی آشوب، انواع توابع آشوب و ارتباط میان آشوب و رمزنگاری مورد بررسی قرار می‌گیرد. در ادامه نیز به بررسی پژوهش‌های پیشین در زمینه‌ی استفاده از سیستم‌های آشوبی در رمزنگاری تصاویر رنگی پرداخته خواهد شد.}

\begin{enumerate}
\def\labelenumi{\arabic{enumi}.}
\setcounter{enumi}{1}
\item
  \protect\phantomsection\label{_Toc232899862}{}\RL{رمزنگاری}
\end{enumerate}

\RL{رمزنگاری یا رمزگذاری} \RL{فرآیندی است که طی آن اطلاعات به‌کمک یک الگوریتم خاص به صورتی غیرقابل‌فهم برای افراد غیرمجاز تبدیل و ذخیره می‌شود تا از دسترسی یا سوء‌استفاده‌ی احتمالی از آن جلوگیری گردد. در هر سیستم رمزنگاری، کلیدی وجود دارد که تنها در اختیار رمزگذار (فرستنده یا گیرنده‌ی مجاز) قرار دارد. این کلید برای معکوس‌سازی عملیات رمزنگاری و بازگردانی داده‌ها به حالت اولیه به‌کار می‌رود که این عمل را رمزگشایی} \RL{می‌نامند}.

\RL{رمزنگاری از دیرباز در میان دولت‌ها، نیروهای نظامی و نهادهای امنیتی برای برقراری ارتباطات محرمانه و ایمن مورد استفاده قرار می‌گرفت. با گذشت زمان و گسترش فناوری‌های دیجیتال، کاربرد رمزنگاری به حوزه‌های غیرنظامی نیز گسترش یافته است؛ به‌گونه‌ای که امروزه در بسیاری از سامانه‌های کامپیوتری و ارتباطی، به‌منظور حفظ امنیت داده‌ها، از روش‌های رمزگذاری استفاده می‌شود. این فناوری می‌تواند برای محافظت از داده‌های ذخیره‌شده در حافظه‌های مختلف مانند دیسک سخت، فلش مموری و یا در جریان انتقال اطلاعات از طریق شبکه‌هایی نظیر اینترنت، تلفن همراه، ارتباطات بی‌سیم و بلوتوث به‌کار گرفته شود} \RL{{[}14{]}.}

\RL{از دیدگاه واژه‌شناسی، واژه‌ی رمزنگاری} \footnote{Cryptography}\RL{از دو واژه‌ی یونانی} krypto \RL{به معنی «پنهان» و} grapho \RL{به معنی «نوشتن» تشکیل شده است. در مفهوم علمی، رمزنگاری به معنای تبدیل متن آشکار}\footnote{Plain Text} \RL{به متن} \RL{رمز} \footnote{Cipher Text}\RL{و} \RL{برعکس است. عملیات تبدیل متن آشکار به متن رمز را رمزگذاری} \RL{و تبدیل متن رمز به متن آشکار را رمزگشایی} \RL{می‌گویند}.

\RL{در دوران باستان، به‌ویژه در مصر، از روش‌های ابتدایی جایگزینی حروف}\footnote{Substitution} \RL{برای رمزگذاری پیام‌ها استفاده می‌شد. در این روش هر حرف با حرف یا نمادی دیگر جایگزین می‌شد که تنها افراد خاص از معنا و قاعده‌ی آن اطلاع داشتند. با وجود کاربرد ساده، این روش با چالش‌هایی مواجه بود؛ زیرا لازم بود برای هر جفت فرستنده و گیرنده، یک مجموعه کد منحصربه‌فرد تعریف شود، که این امر در عمل باعث دشواری استفاده‌ی گسترده از این سیستم می‌گردید{[}15{]}.}

\RL{به همین دلیل، در ادامه‌ی توسعه‌ی علم رمزنگاری، مفهوم} \RL{کلید} (Key) \RL{مطرح} \RL{شد. در این روش، تبدیل متن آشکار به متن رمز براساس پارامتری مستقل از خود متن و مبتنی بر کلید انجام می‌گیرد. به این ترتیب، فرستنده و گیرنده تنها لازم است از الگوریتم رمزنگاری و کلید مشترک اطلاع داشته باشند و نیازی به تبادل دستی کدهای اختصاصی وجود ندارد}.

\RL{در دوران انقلاب صنعتی و هم‌زمان با رشد فناوری‌های مخابراتی، علم رمزنگاری نیز دچار تحول شد. با گسترش تلگراف، رادیو و وسایل ارتباطی جدید، لازم بود پیام‌ها به‌جای حروف انسانی، به کدهای دودویی (۰ و ۱) تبدیل شوند. بدین ترتیب، رمزنگاری از حوزه‌ی زبان طبیعی وارد دنیای دیجیتال گردید و مفاهیم جدیدی مانند بیت، کلید مخفی و الگوریتم‌های ریاضی رمزگذاری مطرح شدند. در این مرحله، رمزنگاری دیگر صرفاً ابزاری برای پنهان‌سازی متن نبود، بلکه به شاخه‌ای از ریاضیات و علوم رایانه تبدیل شد که امنیت داده‌ها را بر اساس مدل‌های ریاضی تضمین می‌کرد{[}16{]}.}

\RL{به طور کلی، علم رمزنگاری امروزه نقش بنیادینی در امنیت ارتباطات دیجیتال، بانکداری الکترونیک، تجارت اینترنتی، شبکه‌های اجتماعی و حتی ذخیره‌سازی ابری ایفا می‌کند و پایه‌ی بسیاری از سامانه‌های امنیتی مدرن را تشکیل می‌دهد.}

\begin{enumerate}
\def\labelenumi{\arabic{enumi}.}
\item
  \protect\phantomsection\label{_Toc232899863}{}\RL{تعریف رمزنگاری}
\end{enumerate}

\RL{رمزنگاری فرآیندی است که طی آن اطلاعات به‌گونه‌ای دچار تغییر و بازآرایی می‌شوند که برای افراد غیرمجاز غیرقابل فهم باشد. فناوری رمزنگاری دسترسی به پیام‌ها و اطلاعات ارسالی را برای افراد غیرمجاز محدود کرده و امکان مشاهده، مطالعه و تفسیر آن‌ها را سلب می‌کند. این فناوری به‌ویژه برای حفاظت از داده‌ها در شبکه‌های عمومی مانند اینترنت کاربرد دارد{[}17{]}.}

\RL{هدف اصلی رمزنگاری، طراحی طرح‌ها و پروتکل‌هایی است که حتی در حضور دشمن نیز امکان انجام فعالیت‌های خاص را فراهم سازد. یکی از اهداف اساسی این حوزه، ایجاد امکان ارتباط امن بر روی کانال‌های ناامن است؛ به‌طوری که حریم خصوصی و اصالت داده‌ها به‌طور کامل حفظ شود. برای رسیدن به این هدف، الگوریتم‌های ریاضی پیشرفته‌ای برای رمزگذاری پیام‌ها و ضمائم مربوطه به‌کار گرفته می‌شوند} \RL{{[}18{]}.}

\RL{در بیشتر روش‌های رمزنگاری، از فرآیندهای برگشت‌پذیر}\footnote{Reversible} \RL{استفاده می‌شود تا دریافت‌کننده در مقصد بتواند متن رمز را به متن آشکار} \RL{بازگرداند. قوانین و تکنیک‌هایی که برای رمزنگاری به‌کار می‌روند، به عنوان الگوریتم‌های رمزنگاری شناخته می‌شوند و میزان سادگی یا پیچیدگی فرایند رمزگذاری را تعیین می‌کنند. بسیاری از الگوریتم‌های رمزنگاری، از کلید} \RL{به‌عنوان پارامتری محرمانه برای رمزگذاری و رمزگشایی استفاده می‌کنند. کلید، در واقع یک کلمه عبور است که تنها فرستنده و گیرنده پیام از آن اطلاع دارند. الگوریتم رمزنگاری با استفاده از قواعد ریاضی، کلید (که معمولاً رشته‌ای طولانی از اعداد است) را بر اطلاعات اعمال می‌کند} \RL{{[}19{]}.}

\RL{تکنیک‌های رمزنگاری پیچیده، برخلاف روش‌های ساده‌ی جایگزینی یا جابجایی، از یک کلید محرمانه برای کنترل توالی‌های طولانی و پیچیده از جایگزینی‌ها و جابجایی‌ها بهره می‌برند. کلید و الگوریتم رمزنگاری با همکاری یکدیگر، متن اولیه را به متن رمز تبدیل می‌کنند. در اغلب سیستم‌ها، الگوریتم رمزنگاری ثابت و شناخته‌شده است و کلید است که خروجی یکتا و امنی تولید می‌کند{[}20{]}.}

\RL{در گذشته، امنیت سیستم‌های رمزنگاری به میزان محرمانه بودن الگوریتم وابسته بود. با پیشرفت علم رمزنگاری، الگوریتم‌ها توانستند از این وابستگی عبور کنند و بر محدودیت طول کلید و محرمانگی آن غالب شوند. امروزه، قابل اعتمادترین الگوریتم‌های رمزنگاری باید پس از چندین مرحله تحقیق و بررسی دقیق و عمومی، توانایی خود را نشان دهند تا اعتماد سازمان‌ها و کاربران را جلب کنند. ایجاد چنین الگوریتم‌هایی، فرآیندی پیچیده و نیازمند دقت بسیار بالا است {[}21{]}.}

\begin{enumerate}
\def\labelenumi{\arabic{enumi}.}
\setcounter{enumi}{1}
\item
  \protect\phantomsection\label{_Toc232899864}{}\RL{شیوه‌های پایه رمزنگاری}
\end{enumerate}

\RL{در رمزنگاری کلاسیک، دو شیوه‌ی پایه وجود دارد که اساس بسیاری از الگوریتم‌های مدرن را تشکیل می‌دهند: رمزنگاری جابجایی}\footnote{Transposition} \RL{و رمزنگاری جایگزینی}\footnote{Substitution}\RL{{[}22{]}.}

\RL{الف) رمزنگاری جابجایی}

\RL{رمزنگاری جابجایی روشی است که در آن متن اصلی بدون تغییر محتوا، به‌صورت ترتیبی جابجا می‌شود تا متن رمز تولید شود. عملیات رمزگشایی نیز با اجرای برعکس همان جابجایی امکان‌پذیر است. از دیدگاه ریاضی، یک تابع یک‌به‌یک بر روی موقعیت‌های متن اعمال می‌شود و معکوس آن برای بازگرداندن متن اصلی استفاده می‌گردد}.

\RL{ویژگی اصلی رمزنگاری جابجایی این است که محتوای متن تغییر نمی‌کند، بلکه ترتیب آن دگرگون می‌شود. این خاصیت باعث می‌شود که تغییر حتی یک بیت از کلید یا متن اصلی، خروجی رمز متفاوتی ایجاد کند؛ به عبارتی، سیستم رمز به تغییرات کوچک در متن یا کلید بسیار حساس است. نمونه‌ای ساده از این شیوه، استفاده از جایگشت} \RL{است که ترکیب نمادهای متن را تغییر می‌دهد. در رمزنگاری تصویر، فرآیند جابجایی معادل جابجایی موقعیت پیکسل‌ها در تصویر است، بدون اینکه مقدار آن‌ها تغییر یابد{[}23{]}.}

\RL{ب) رمزنگاری جایگزینی}

\RL{رمزنگاری جایگزینی، روشی است که در آن هر واحد از متن اصلی با یک نماد دیگر طبق یک سیستم مشخص جایگزین می‌شود تا متن رمز حاصل گردد. هدف اصلی این شیوه، جلوگیری از مشاهده‌ی الگوهای تکراری در متن رمز است. عملیات رمزگشایی با استفاده از عکس عمل جایگذاری امکان‌پذیر می‌شود}.

\RL{ساده‌ترین نوع جایگزینی، جانشینی یک نماد با نماد دیگر است. در رمزنگاری تصویر، این فرآیند برای تغییر مقادیر پیکسل‌ها به کار می‌رود. مقایسه‌ی رمزنگاری جایگزینی و جابجایی نشان می‌دهد که در جایگزینی، مکان نمادها ثابت است و تنها مقدار آن‌ها تغییر می‌کند، در حالی که در جابجایی، نمادها دست‌نخورده باقی می‌مانند اما موقعیت آن‌ها تغییر می‌کند{[}22،24{]}.}

\begin{enumerate}
\def\labelenumi{\arabic{enumi}.}
\setcounter{enumi}{2}
\item
  \protect\phantomsection\label{_Toc232899865}{}\RL{انواع روش‌های رمزنگاری بر اساس نوع کلید}
\end{enumerate}

\RL{بر اساس نوع کلید، الگوریتم‌های رمزنگاری به دو دسته تقسیم می‌شوند: رمزنگاری متقارن}\footnote{Symmetric Key Cryptography} \RL{و رمزنگاری نامتقارن}\footnote{Asymmetric Key Cryptography}\RL{{[}25{]}.}

\RL{الف) رمزنگاری متقارن}

\RL{رمزنگاری متقارن که گاهی به آن رمز مشترک نیز گفته می‌شود، دسته‌ای از الگوریتم‌هاست که برای رمزگذاری و رمزگشایی پیام‌ها از کلیدهای یکسان یا مشابه استفاده می‌کنند. در عمل، این کلید نشان‌دهنده یک راز مشترک میان دو یا چند طرف است که برای حفظ اطلاعات خصوصی به کار می‌رود. امنیت این روش وابسته به طول کلید و محرمانه بودن آن است. یکی از مزایای رمزنگاری متقارن، سرعت بالای پردازش و اجرای الگوریتم‌ها است که آن را به گزینه‌ای مناسب برای انتقال اطلاعات در اینترنت یا ذخیره‌سازی داده‌های حساس در بانک‌های اطلاعاتی تبدیل می‌کند{[}26{]}.}

\RL{در الگوریتم‌های متقارن، علاوه بر ارسال پیام رمز شده، کلید مورد استفاده نیز باید به مقصد منتقل شود یا در آنجا موجود باشد. این امر نیازمند وجود یک کانال امن برای توزیع کلید است. در صورت فقدان چنین کانالی یا عدم امکان تأیید هویت گیرنده کلید، امنیت سیستم به خطر می‌افتد. همچنین، تبادل مکرر کلید می‌تواند بار اضافی بر سیستم تحمیل کند. به همین دلیل، روش‌های متقارن بیشترین بازده را در محیط‌هایی دارند که کانال امن برای توزیع و نگهداری کلیدها وجود دارد یا فرستنده و گیرنده قبلاً هویت یکدیگر را تأیید کرده‌اند}.

\RL{رمزنگاری متقارن می‌تواند به دو نوع تقسیم شود}:

\RL{رمزهای} \RL{جریانی}\footnote{Stream Ciphers} : \RL{هر بیت} \RL{یا رقم پیام به صورت مستقل رمزگذاری می‌شود}.

\RL{رمزهای بلوکی}\footnote{Block Ciphers} : \RL{تعدادی بیت به‌صورت یک واحد رمزگذاری می‌شوند}.

\RL{استاندارد} AES\footnote{Advanced Encryption Standard} \RL{که توسط} NIST \RL{در دسامبر 2001 تایید شد، نمونه‌ای از الگوریتم بلوکی با بلوک‌های 128 بیتی است. رمزهای جریانی سرعت بیشتری نسبت به رمزهای بلوکی دارند زیرا پیچیدگی الگوریتمی کمتری دارند. از نمونه‌های مشهور الگوریتم‌های متقارن می‌توان به سیستم‌های کلاسیک مانند} Vigenere\RL{،} Substitution\RL{،} DES\RL{،} AES \RL{و} IDES \RL{و همچنین رمزهای جریانی مبتنی بر} LFSR \RL{اشاره کرد{[}25،27{]}.}

\RL{ب) رمزنگاری نامتقارن}

\RL{رمزنگاری نامتقارن یا رمزنگاری کلید عمومی نوعی از رمزنگاری است که در آن کلید رمزگذاری با کلید رمزگشایی متفاوت است؛ برخلاف رمزنگاری متقارن که هر دو عملیات با یک کلید انجام می‌شود. در این روش، هر کاربر یک جفت کلید در اختیار دارد}:

\RL{کلید عمومی: برای رمزگذاری پیام‌ها و تأیید صحت امضای دیجیتال}.

\RL{کلید خصوصی: برای رمزگشایی پیام‌ها و امضای دیجیتال داده‌ها}.

\RL{کلید خصوصی مخفی نگه داشته می‌شود، در حالی که کلید عمومی می‌تواند به طور گسترده منتشر شود. پیام‌های رمز شده توسط کلید عمومی تنها توسط دارنده‌ی کلید خصوصی متناظر قابل خواندن هستند. کلید خصوصی نمی‌تواند از کلید عمومی محاسبه شود، اما بین این دو کلید یک رابطه ریاضی وجود دارد که امکان رمزگذاری و رمزگشایی امن را فراهم می‌کند}.

\RL{در روش نامتقارن، کلیدها دارای ساختار مشخصی با هشت فیلد اطلاعاتی هستند که دو فیلد به کلید عمومی و شش فیلد به کلید خصوصی اختصاص دارد {[}28{]}. این ساختار پیچیده سبب می‌شود که پیام‌ها بدون دسترسی به کلید خصوصی، غیرقابل خواندن باشند و امنیت انتقال داده تضمین گردد.}

\begin{enumerate}
\def\labelenumi{\arabic{enumi}.}
\setcounter{enumi}{3}
\item
  \protect\phantomsection\label{_Toc232899866}{}\RL{مقایسه روش‌های رمزنگاری متقارن و نامتقارن}
\end{enumerate}

\RL{یکی از تفاوت‌های اساسی میان روش‌های رمزنگاری متقارن و نامتقارن، نحوه مدیریت کلیدها است. در رمزنگاری نامتقارن، برخلاف رمزنگاری متقارن، نیازی به رد و بدل کردن کلید خصوصی وجود ندارد و هر کاربر تنها از کلید عمومی برای رمزگذاری پیام استفاده می‌کند و کلید خصوصی خود را مخفی نگه می‌دارد}.

\RL{در روش‌های متقارن، کلید رمزگذاری و رمزگشایی یکسان است و همواره خطر دسترسی غیرمجاز نفوذگر به کلید در زمان انتقال آن وجود دارد. بنابراین، امنیت روش‌های متقارن وابسته به حفظ محرمانگی کلید و وجود کانال امن برای انتقال آن است}.

\RL{به طور کلی، مزیت اصلی رمزنگاری متقارن سرعت بالای آن و مزیت اصلی رمزنگاری نامتقارن امنیت بالای آن است. در برخی تحقیقات، تلاش شده است که با ترکیب این دو روش، یک روش رمزنگاری ارائه شود که هم ایمن و هم سریع باشد{[}29{]}.}

\RL{تفاوت‌های اصلی میان رمزنگاری متقارن و نامتقارن از نظر نوع کلید، نیاز به کانال امن و مزیت‌های هر روش، به‌صورت خلاصه در} \RL{جدول (2-1) ارائه شده} \RL{است. {[}30{]}.}

\begin{enumerate}
\def\labelenumi{\arabic{enumi}.}
\item
  \protect\phantomsection\label{_Toc232899943}{}\RL{مقایسه رمزنگاری متقارن و نامتقارن}
\end{enumerate}

{\def\LTcaptype{none} % do not increment counter
\begin{longtable}[]{@{}
  >{\raggedright\arraybackslash}p{(\linewidth - 4\tabcolsep) * \real{0.1580}}
  >{\raggedright\arraybackslash}p{(\linewidth - 4\tabcolsep) * \real{0.1716}}
  >{\raggedright\arraybackslash}p{(\linewidth - 4\tabcolsep) * \real{0.3459}}@{}}
\toprule\noalign{}
\begin{minipage}[b]{\linewidth}\centering
\textbf{\RL{ویژگی}}
\end{minipage} & \begin{minipage}[b]{\linewidth}\centering
\textbf{\RL{رمزنگاری متقارن}}
\end{minipage} & \begin{minipage}[b]{\linewidth}\centering
\textbf{\RL{رمزنگاری نامتقارن}}
\end{minipage} \\
\midrule\noalign{}
\endhead
\bottomrule\noalign{}
\endlastfoot
\RL{نوع کلید} & \RL{یک کلید} & \RL{یک کلید عمومی و یک کلید خصوصی} \\
\RL{نیاز به کانال امن} & \RL{نیازمند کانال امن} & \RL{نیاز به کانال امن ندارد} \\
\RL{مزیت اصلی} & \RL{سرعت بالا} & \RL{امنیت بالا} \\
\end{longtable}
}

\begin{enumerate}
\def\labelenumi{\arabic{enumi}.}
\setcounter{enumi}{4}
\item
  \protect\phantomsection\label{_Toc232899867}{}\RL{ویژگی‌های یک الگوریتم رمزنگاری امن}
\end{enumerate}

\RL{یک الگوریتم رمزنگاری امن باید دارای ویژگی‌های زیر باشد تا بتواند حفاظت کامل از اطلاعات را تضمین کند {[}31{]}:}

\begin{enumerate}
\def\labelenumi{\arabic{enumi}.}
\item
  \textbf{\RL{یک‌طرفه بودن الگوریتم}:}\\
  \RL{دستیابی به متن اصلی از طریق متن رمز بدون داشتن کلید، غیرممکن باشد. در واقع، کلید خصوصی مانند دریچه‌ای است که تنها با در اختیار داشتن آن می‌توان متن رمزشده را به متن اصلی بازگرداند}.
\item
  \textbf{\RL{عدم وابستگی امنیت به محرمانه بودن الگوریتم}:}\\
  \RL{مخفی نگاه داشتن جزئیات الگوریتم در امنیت آن نقش ندارد و امنیت باید صرفاً بر اساس کلید مخفی تضمین شود}.
\item
  \textbf{\RL{وابستگی هر بیت متن رمزشده به تمامی بیت‌های متن اصلی}:}\\
  \RL{کوچک‌ترین تغییر در متن اصلی باید متن رمزشده کاملاً متفاوتی تولید کند. الگوریتم‌هایی که این ویژگی را دارند، به عنوان \textbf{رمز کامل}}\footnote{\textbf{Perfect Cipher}} \RL{شناخته می‌شوند}.
\item
  \textbf{\RL{وابستگی هر بیت متن رمزشده به تمامی بیت‌های کلید سری}:}\\
  \RL{در این حالت، کوچک‌ترین تغییر در کلید نیز باعث تولید متن رمزشده متفاوت می‌شود}.
\item
  \textbf{\RL{تأثیر تغییرات کوچک ورودی بر خروجی}:}\\
  \RL{تغییر حتی یک بیت در داده‌های ورودی، بدون تغییر کلید، باید تغییرات عمده‌ای در داده‌های خروجی ایجاد کند}.
\item
  \textbf{\RL{تأثیر تغییرات کوچک کلید بر متن رمزشده}:}\\
  \RL{تغییر حتی یک بیت در کلید سری، بدون تغییر متن اصلی، باید باعث ایجاد تغییرات عمده در متن رمزشده شود}.
\item
  \textbf{\RL{استفاده از عملیات جانشینی و جابجایی بیت‌ها تحت کنترل کلید و داده ورودی}:}\\
  \RL{الگوریتم باید قادر باشد بیت‌ها را جایگزین و جابجا کند تا امنیت اطلاعات تضمین شود}.
\item
  \textbf{\RL{دارا بودن ساختار جبری پیچیده}:}\\
  \RL{اگر تابع رمزنگاری دارای رابطه جبری ساده باشد، الگوریتم به راحتی قابل تحلیل و شکستن خواهد بود}.
\item
  \textbf{\RL{برابری طول متن اصلی و متن رمزشده}:}\\
  \RL{الگوریتم باید طول پیام ورودی و خروجی را حفظ کند}.
\item
  \textbf{\RL{استفاده از تمامی کلیدهای مخفی برای تولید رمز قوی}:}\\
  \RL{هر کلید باید توانایی تولید متن رمزشده قوی و غیرقابل پیش‌بینی را داشته باشد}.
\end{enumerate}

\RL{این خصوصیات شرایط لازم برای طراحی یک الگوریتم رمزنگاری قوی را فراهم می‌کنند. با این حال،} \RL{شروط لازم و کافی برای ارزیابی امنیت هر سیستم رمزنگاری، مقاومت} \RL{آن در برابر انواع حملات است. در بخش‌های بعدی، انواع حملات و روش‌های مقابله با آن‌ها بررسی خواهد شد}.

\begin{enumerate}
\def\labelenumi{\arabic{enumi}.}
\item
  \protect\phantomsection\label{_Toc232899868}{}\RL{تحلیل رمز: انواع حملات به سامانه‌های رمزنگاری}
\end{enumerate}

\RL{قضیه امنیت کامل، معیار مهمی برای طراحی سامانه‌های رمزنگاری فراهم می‌کند؛ با این حال، طراحی سامانه‌هایی که این قضیه کاملاً برای آن‌ها صدق کند، در عمل دشوار است. در رمزنگاری مدرن، تمرکز بیشتر بر کاربردپذیری عملی و امنیت کافی برای محافظت از داده‌های حساس است. الگوریتم‌های رمزنگاری باید نسبت به کلیدهای اشتباه حساس باشند و فضای کلید آن‌ها باید به اندازه کافی بزرگ باشد تا در مقابل حملات کورکورانه و شدید مقاومت نشان دهند{[}32{]}.}

\RL{یک رمزگشا یا تحلیلگر رمز سعی می‌کند بر اساس اطلاعات موجود، متن رمز شده را رمزگشایی کند. میزان دانش تحلیلگر می‌تواند شامل متن رمزشده، دسترسی محدود به الگوریتم رمزگذاری و رمزگشایی باشد، اما هیچگاه دسترسی به کلید ندارد. هرچه دانش تحلیلگر بیشتر باشد، رمزگشایی آسان‌تر خواهد بود. برای تضمین امنیت، باید اطمینان حاصل شود که معکوس تابع رمزگذاری تنها با داشتن کلید ممکن است. در اصطلاح رمزنگاری، این توابع باید یک‌طرفه باشند}. \RL{{[}33{]}.}

\RL{انواع حملات به چهار دسته اصلی تقسیم می‌شوند{[}34{]}}:

\begin{enumerate}
\def\labelenumi{\arabic{enumi}.}
\item
  \textbf{\RL{حمله با متن رمزشده}}\footnote{\textbf{Ciphertext-Only Attack}} \RL{}
\end{enumerate}

\RL{در این نوع حمله، تحلیلگر تنها به متن رمزشده دسترسی دارد و تلاش می‌کند متن اصلی یا کلید مورد استفاده را بیابد تا سایر پیام‌های رمز شده نیز قابل رمزگشایی شوند. این حمله سخت‌ترین نوع حمله است و نیازمند تحلیل دقیق الگوها و احتمالات است.}

\begin{enumerate}
\def\labelenumi{\arabic{enumi}.}
\setcounter{enumi}{1}
\item
  \textbf{\RL{حمله با متن اصلی مشخص}} \footnote{\textbf{Known-Plaintext Attack}}
\end{enumerate}

\RL{در این سناریو، تحلیلگر به یک یا چند جفت متن اصلی و متن رمزشده متناظر دسترسی دارد و از آن‌ها برای حدس زدن کلید یا الگوریتم رمزنگاری استفاده می‌کند. هدف این است که کلید} K \RL{یا الگوریتمی پیدا شود که بتواند پیام‌های جدید را رمزگشایی نماید}.

\begin{enumerate}
\def\labelenumi{\arabic{enumi}.}
\setcounter{enumi}{2}
\item
  \textbf{\RL{حمله با متن اصلی منتخب}} \footnote{\textbf{Chosen-Plaintext Attack}}
\end{enumerate}

\RL{در این نوع حمله، تحلیلگر می‌تواند متن‌های اصلی خاصی را برای رمزگذاری انتخاب کند و متن‌های رمزشده متناظر را مشاهده نماید. این حمله از نوع} Known-Plaintext \RL{قوی‌تر است، زیرا تحلیلگر می‌تواند بلوک‌های متن اصلی را به گونه‌ای انتخاب کند که اطلاعات بیشتری درباره کلید به دست آورد. هدف، پیدا کردن کلید} K \RL{یا الگوریتمی است که امکان رمزگشایی پیام‌های جدید با همان کلید را فراهم کند.}

\begin{enumerate}
\def\labelenumi{\arabic{enumi}.}
\setcounter{enumi}{3}
\item
  \textbf{\RL{حمله با متن رمزشده}} \textbf{\RL{منتخب}} \footnote{\textbf{Chosen-Ciphertext Attack}}\RL{}
\end{enumerate}

\RL{این نوع حمله، حالت خاصی از حمله متن اصلی منتخب است. در آن، تحلیلگر می‌تواند متن‌های اصلی را انتخاب و خروجی رمزگذاری را مشاهده کند و همچنین انتخاب‌های بعدی خود را بر اساس نتایج رمزگذاری قبلی اصلاح نماید. این حمله تطبیقی، امکان انتخاب بلوک‌های کوچک‌تر و تحلیل دقیق‌تر کلید را فراهم می‌کند}.

\RL{این چهار دسته، اسکلت اصلی انواع حملات در رمزنگاری مدرن را تشکیل می‌دهند. یک سامانه رمزنگاری پس از مقاومت موفق در برابر این چهار حالت، می‌تواند ادعای امنیت داشته باشد}.

\begin{itemize}
\item
  \RL{الگوریتم‌های ساده‌تر معمولاً در برابر حملات} Ciphertext-Only \RL{و} Known-Plaintext \RL{آسیب‌پذیر هستند}.
\item
  \RL{برای الگوریتم‌های پیچیده‌تر، مانند رمزنگاری مبتنی بر توابع آشوب، معمولاً از حملات} Chosen-Plaintext \RL{و} Chosen-Ciphertext \RL{استفاده} \RL{می‌شود {[}35{]}.}

  \begin{enumerate}
  \def\labelenumi{\arabic{enumi}.}
  \item
    \protect\phantomsection\label{_Toc232899869}{}\RL{امنیت سیستم‌های رمزنگاری}
  \end{enumerate}
\end{itemize}

\RL{امنیت سیستم‌های رمزنگاری از دیرباز یکی از مسائل کلیدی در علم رمزنگاری بوده و همچنان اهمیت فراوان دارد. طراحی سامانه‌های رمزنگاری باید به گونه‌ای باشد که علاوه بر کارایی عملی، حداکثر حفاظت از داده‌ها را تضمین کند}.

\begin{enumerate}
\def\labelenumi{\arabic{enumi}.}
\item
  \protect\phantomsection\label{_Toc232899870}{}\RL{اصول شش\-گانه کرشهف}
\end{enumerate}

\RL{آگوست کرشهف}\footnote{August Kerckhoffs} \RL{در سال 1883 دو مقاله با عنوان رمزنگاری نظامی منتشر کرد که در آن‌ها شش اصل اساسی مطرح شد. این اصول هنوز هم در طراحی سیستم‌های رمزنگاری مدرن مورد استفاده قرار می‌گیرند:(37)}

\begin{enumerate}
\def\labelenumi{\arabic{enumi}.}
\setcounter{enumi}{4}
\item
  \RL{غیرقابل شکست بودن سیستم: سیستم رمزنگاری باید حتی در عمل غیرقابل شکست باشد}.
\item
  \RL{محرمانه بودن تنها کلید رمز: هیچ جزئیات دیگری از سیستم رمزنگاری نباید مخفی باشد و تنها کلید، اطلاعات محرمانه است}.
\item
  \RL{قابلیت انتخاب و تغییر کلید: کلید باید به‌راحتی قابل انتخاب، تغییر و حفظ باشد و نیازی به یادداشت کردن آن نباشد}.
\item
  \RL{امکان ارسال از طریق تلگراف: متون رمزشده باید قابلیت مخابره از طریق خطوط تلگراف یا رسانه‌های مشابه را داشته باشند}.
\item
  \RL{قابلیت حمل و نقل: دستگاه رمزنگاری یا اسناد رمز شده باید توسط یک نفر قابل حمل باشد}.
\item
  \RL{سهولت راه‌اندازی سیستم: سیستم رمزنگاری باید به سادگی قابل استفاده و راه‌اندازی باشد}. \RL{}
\end{enumerate}

\RL{در توضیح قانون دوم کرشهف باید عنوان کرد که به چند دلیل عمده کلید رمز باید تنها اطلاعات} \RL{محرمانه در یک سیستم رمزنگاری باشد:}

\begin{enumerate}
\def\labelenumi{\arabic{enumi}.}
\item
  \RL{سادگی حفظ محرمانگی کلید نسبت به الگوریتم: محرمانه نگه داشتن یک رشته بیت (مثلاً 512 بیت) بسیار ساده‌تر و عملی‌تر از پنهان نگاه داشتن کل الگوریتم یا روش پیاده‌سازی است}.
\item
  \RL{تجدید امنیت در صورت افشای کلید: در صورت افشای کلید، تنها با تعویض کلید می‌توان مجدداً به امنیت دست یافت؛ در حالی که وابستگی به الگوریتم، افشای آن را به یک مسئله بسیار پیچیده تبدیل می‌کند}.
\item
  \RL{امکان استفاده از سیستم مشترک برای چند گیرنده/فرستنده: اگر کلید تنها اطلاعات محرمانه باشد، می‌توان با کلیدهای متفاوت با گیرنده‌ها و فرستنده‌های مختلف ارتباط برقرار کرد؛ در غیر این صورت، برای هر جفت فرستنده/گیرنده نیاز به الگوریتم جدید است}.
\item
  \RL{بازرسی و اصلاح الگوریتم‌های عمومی توسط محققان: وقتی الگوریتم و روش پیاده‌سازی عمومی باشد، محققان امنیتی می‌توانند مشکلات و حفره‌های احتمالی آن را پیش از سوء استفاده توسط حمله‌گران شناسایی و رفع کنند؛ بنابراین، سیستم‌های رمزنگاری مبتنی بر الگوریتم‌های عمومی معمولاً قابل اعتمادتر هستند{[}36{]}.}

  \begin{enumerate}
  \def\labelenumii{\arabic{enumii}.}
  \item
    \protect\phantomsection\label{_Toc232899871}{}\RL{معیارهای پنجگانه شانون}
  \end{enumerate}
\end{enumerate}

\RL{در سال ۱۹۴۸، کلاد شانون}\footnote{Claude Shannon} \RL{مقاله‌ای با عنوان یک نظریه ریاضیاتی از ارتباطات منتشر کرد که در دو بخش ارائه شد. در این مقاله، شانون به بررسی بهترین روش کدگذاری اطلاعات برای مخابره توسط فرستنده پرداخت و پنج معیار اساسی برای ارزیابی سیستم‌های رمزنگاری معرفی نمود{[}37{]}:}

\RL{۱}. \RL{میزان ایمنی سیستم}

\RL{برخی سیستم‌ها امنیت کامل دارند و هیچ مهاجمی نمی‌تواند آن‌ها را شکسته یا به داده‌ها دسترسی یابد. در برخی سیستم‌ها، مهاجم ممکن است اطلاعاتی کسب کند که منجر به جواب یکتا نشود و حتی در مواردی که پاسخ یکتا حاصل شود، پیچیدگی محاسبات و زمان لازم برای شکستن رمز مانع آسیب به امنیت سیستم می‌شود}.

\RL{۲}. \RL{اندازه کلید}

\RL{کلید باید به روشی مطمئن و غیرقابل دسترس در اختیار گیرنده قرار گیرد، به‌گونه‌ای که حتی در شرایط اضطراری نیز امکان استفاده از آن وجود داشته باشد. بنابراین، داشتن کلید کوچک و ساده تا حد امکان مطلوب است}.

\RL{۳}. \RL{پیچیدگی عملیات رمزنگاری و رمزگشایی}

\RL{عملیات رمزگذاری و رمزگشایی باید ساده، سریع و با هزینه کم انجام شود تا کارایی سیستم حفظ گردد}.

\RL{۴}. \RL{انتشار خطا} \footnote{Error Propagation}

\RL{اگر در فرآیند رمزگذاری یا انتقال پیام، خطایی رخ دهد، ممکن است این خطا در متن بازگشایی شده گسترش یافته و باعث از دست رفتن مقدار زیادی اطلاعات شود. سیستم باید طوری طراحی شود که انتشار خطا به حداقل برسد}.

\RL{۵}. \RL{بسط یا گسترش پیام}

\RL{در برخی سیستم‌ها، پس از رمزگذاری، اندازه پیام افزایش می‌یابد. به‌ویژه در سیستم‌هایی که قصد حذف فرکانس طبیعی سمبل‌ها را دارند، ممکن است از سمبل‌های خنثی یا تکرار سمبل‌ها استفاده شود. این امر در سیستم‌های مخابراتی مطلوب نیست و باید از آن جلوگیری شود.}

\begin{enumerate}
\def\labelenumi{\arabic{enumi}.}
\item
  \protect\phantomsection\label{_Toc232899872}{}\RL{ویژگی‌های نرم‌افزاری الگوریتم رمز}
\end{enumerate}

\RL{در طراحی الگوریتم‌های رمزنگاری، رعایت ویژگی‌های نرم‌افزاری اهمیت زیادی دارد تا پیاده‌سازی سریع، پایدار و مقاوم در برابر اختلال‌ها حاصل شود. برخی از ویژگی‌های کلیدی عبارتند از{[}38{]}:}

\begin{enumerate}
\def\labelenumi{\arabic{enumi}.}
\item
  \RL{اجتناب از پرش‌های شرطی در حلقه‌های داخلی الگوریتم}:
\end{enumerate}

\RL{هرگونه تغییر غیرقابل پیش‌بینی در جریان کنترل الگوریتم، موجب اختلال در عملکرد پردازش و افزایش تعداد سیکل‌های ساعت مورد نیاز می‌شود}.

\begin{enumerate}
\def\labelenumi{\arabic{enumi}.}
\item
  \RL{استفاده محدود از عملگرهای سنگین}:
\end{enumerate}

\RL{عملگرهایی مانند ضرب و تقسیم یا سایر عملگرهای پیچیده که اجرای آن‌ها بر روی پردازنده‌ها دشوار است، باید به حداقل برسند تا کارایی الگوریتم کاهش نیابد}.

\begin{enumerate}
\def\labelenumi{\arabic{enumi}.}
\setcounter{enumi}{1}
\item
  \RL{محدود کردن تعداد متغیرها در طرح الگوریتم}:
\end{enumerate}

\RL{هرچند بسیاری از پردازنده‌های مدرن دارای ثبات‌های متعدد چندمنظوره هستند، اما برخی پردازنده‌ها محدودیت دارند و استفاده بهینه از متغیرها ضروری است}.

\begin{enumerate}
\def\labelenumi{\arabic{enumi}.}
\setcounter{enumi}{2}
\item
  \RL{کوچک نگه داشتن اندازه جداول}:
\end{enumerate}

\RL{هرچند جداول بزرگ از نظر امنیتی مطلوب‌تر هستند، اما برای اجرای سریع‌تر نرم‌افزاری باید اندازه آن‌ها محدود باشد. توصیه می‌شود جداول به گونه‌ای طراحی شوند که بیش از ۴ کیلوبایت حافظه نیاز نداشته باشند}.

\begin{enumerate}
\def\labelenumi{\arabic{enumi}.}
\setcounter{enumi}{3}
\item
  \RL{بهره‌گیری از الگوهای موازی در عملگرها}:
\end{enumerate}

\RL{اجرای همزمان و مستقل عملگرها می‌تواند به‌طور قابل توجهی سرعت اجرای الگوریتم را افزایش دهد. استفاده از موازی‌سازی در طراحی نرم‌افزاری، یکی از روش‌های مؤثر برای بهینه‌سازی عملکرد است}.

\RL{رعایت نکات فوق، هرچند ساده به نظر می‌رسد، اما تاثیر قابل توجهی در افزایش سرعت و کارایی پیاده‌سازی نرم‌افزاری الگوریتم‌های رمزنگاری دارد و یکی از الزامات طراحی بهینه محسوب می‌شود.}

\begin{enumerate}
\def\labelenumi{\arabic{enumi}.}
\item
  \protect\phantomsection\label{_Toc232899873}{}\RL{رمزنگاری تصویر دیجیتال}
\end{enumerate}

\RL{با گسترش شبکه‌های عمومی مانند اینترنت و افزایش فناوری‌های سیگنال دیجیتال، استفاده از تصاویر دیجیتال در کاربردهای مختلف روزبه‌روز رایج‌تر شده است و رشد آن همچنان ادامه دارد{[}39{]}. در این میان، یکی از چالش‌های اصلی، تضمین امنیت تصاویر و جلوگیری از دسترسی غیرمجاز به اطلاعات حساس است}.

\RL{یک راهکار مستقیم برای رفع این مشکل، استفاده از رمزنگاری تصاویر است. رمزنگاری می‌تواند محتوای تصویر را پنهان کرده و از دسترسی غیرمجاز جلوگیری نماید}.

\begin{enumerate}
\def\labelenumi{\arabic{enumi}.}
\item
  \protect\phantomsection\label{_Toc232899874}{}\RL{محدودیت‌های روش‌های سنتی رمزنگاری برای تصاویر}
\end{enumerate}

\RL{روش‌های کلاسیک رمزنگاری، مانند}:

\begin{itemize}
\item
  \RL{روش‌های مبتنی بر نظریه اعداد،}
\item
  \RL{روش‌های رمزنگاری استاندارد ساده،}
\item
  \RL{و روش} RSA
\end{itemize}

\RL{برای رمزنگاری تصویر دیجیتال کارا نیستند {[}40{]}. دلایل اصلی عبارتند از}:

\begin{enumerate}
\def\labelenumi{\arabic{enumi}.}
\item
  \RL{حجم بالای تصاویر}:
\end{enumerate}

\RL{تصاویر دیجیتال معمولاً حجم زیادی دارند، و الگوریتم‌های سنتی نیازمند پردازش طولانی و توان محاسباتی بالا هستند}.

\begin{enumerate}
\def\labelenumi{\arabic{enumi}.}
\setcounter{enumi}{1}
\item
  \RL{زمان رمزگذاری و رمزگشایی طولانی}:
\end{enumerate}

\RL{در ارتباطات بلادرنگ، سرعت پایین این الگوریتم‌ها ممکن است باعث تاخیر قابل توجه شود}.

\begin{enumerate}
\def\labelenumi{\arabic{enumi}.}
\setcounter{enumi}{2}
\item
  \RL{همبستگی پیکسل‌های مجاور و یکنواختی رنگ}:
\end{enumerate}

\RL{الگوریتم‌های سنتی به دلیل وجود همبستگی بین پیکسل‌ها و زمینه‌های یکنواخت تصویر، قادر به ایجاد امنیت کافی نیستند}.

\begin{enumerate}
\def\labelenumi{\arabic{enumi}.}
\setcounter{enumi}{3}
\item
  \RL{نیازمندی‌های خاص برنامه‌های کاربردی تصویر}:
\end{enumerate}

\RL{برخی برنامه‌ها نیازمند همزمان برآورده کردن مواردی مانند ذخیره صحیح، پردازش بدون وقفه، فشرده‌سازی اطلاعات و ثبات فرمت تصویر هستند. روش‌های سنتی معمولاً نمی‌توانند همزمان این نیازها و امنیت بالا را تأمین کنند}.

\begin{enumerate}
\def\labelenumi{\arabic{enumi}.}
\item
  \protect\phantomsection\label{_Toc232899875}{}\RL{ضرورت توسعه الگوریتم‌های ویژه رمزنگاری تصویر}
\end{enumerate}

\RL{با توجه به محدودیت‌های فوق، الگوریتم‌های ویژه رمزنگاری تصویر مورد نیاز هستند که بتوانند}:

\begin{itemize}
\item
  \RL{حجم بالای داده‌ها را به‌طور بهینه مدیریت کنند،}
\item
  \RL{سرعت رمزگذاری و رمزگشایی بالا داشته باشند،}
\item
  \RL{امنیت مناسبی برای جلوگیری از دسترسی غیرمجاز فراهم کنند،}
\item
  \RL{و همزمان نیازمندی‌های پردازشی و فشرده‌سازی تصویر را برآورده سازند {[}41{]}.}

  \begin{enumerate}
  \def\labelenumi{\arabic{enumi}.}
  \item
    \protect\phantomsection\label{_Toc232899876}{}\RL{خصوصیات تصویر دیجیتال}

    \begin{enumerate}
    \def\labelenumii{\arabic{enumii}.}
    \item
      \protect\phantomsection\label{_Toc232899877}{}\RL{پیکسل در تصویر}
    \end{enumerate}
  \end{enumerate}
\end{itemize}

\RL{در تصاویر دیجیتال، پیکسل کوچک‌ترین واحد ساختاری تصویر محسوب می‌شود. پیکسل گاهی در مباحث گرافیکی و تصویری به عنوان نقطه نیز شناخته می‌شود و نقش کوچک‌ترین نقطه تشکیل‌دهنده تصویر را ایفا می‌کند}. \RL{همانطور که بیت کوچک‌ترین واحد اطلاعات قابل پردازش توسط کامپیوتر است، پیکسل نیز کوچک‌ترین عنصر سخت‌افزاری و نرم‌افزاری نمایش یا چاپ تصویر است که برای شکل‌گیری تصاویر دیجیتال به کار می‌رود}. \RL{اگر برای هر پیکسل دو رنگ مانند سیاه و سفید در نظر گرفته شود، می‌توان آن را با یک بیت نمایش داد. در صورتی که بیش از دو بیت برای هر پیکسل استفاده شود، محدوده رنگ‌ها و سایه‌های خاکستری گسترده‌تری قابل ارائه خواهد بود. این ویژگی اساس نمایش تصاویر رنگی و توانایی تفکیک رنگ‌ها را فراهم می‌کند{[}42{]}.}

\begin{enumerate}
\def\labelenumi{\arabic{enumi}.}
\setcounter{enumi}{1}
\item
  \protect\phantomsection\label{_Toc232899878}{}\RL{عمق رنگ}
\end{enumerate}

\RL{عمق رنگ تعداد رنگ‌هایی است که یک پیکسل می‌تواند تولید کند، و عمق بیتی تعداد بیت‌هایی است که برای تولید این رنگ‌ها اختصاص داده می‌شود}

\RL{برای تصاویر تک‌رنگ، یک بیت برای نمایش هر پیکسل کافی است، اما در تصاویر رنگی و تصاویر با سایه‌های خاکستری، هر پیکسل نیاز به بیش از یک بیت دارد. هرچه تعداد بیت‌های بیشتری برای نمایش پیکسل اختصاص یابد، تعداد رنگ‌ها و سایه‌های قابل نمایش افزایش می‌یابد{[}42{]}.}

\begin{enumerate}
\def\labelenumi{\arabic{enumi}.}
\setcounter{enumi}{2}
\item
  \protect\phantomsection\label{_Toc232899879}{}\RL{ابعاد تصویر}
\end{enumerate}

\RL{تصاویر دیجیتال معمولاً به صورت شبکه‌های مربعی شکل نمایش داده می‌شوند که شبیه صفحه شطرنج یا موزاییک‌های کف هستند. این شبکه‌ها از مربع‌های کوچکتری به نام پیکسل تشکیل شده‌اند}.

\RL{ابعاد شبکه معادل تعداد پیکسل‌های طول و عرض تصویر است و ارتباطی با اندازه واقعی تصویر ندارد. برای مثال، یک شبکه 1280×720 نشان‌دهنده تعداد پیکسل‌های طول و عرض تصویر است، نه اندازه فیزیکی آن{[}42{]}.}

\begin{enumerate}
\def\labelenumi{\arabic{enumi}.}
\setcounter{enumi}{3}
\item
  \protect\phantomsection\label{_Toc232899880}{}\RL{مدل‌های رنگ}
\end{enumerate}

\RL{مدل‌های رنگ به منظور استانداردسازی نمایش رنگ‌ها استفاده می‌شوند و هر رنگ در تصویر با یک پیکسل مشخص می‌شود. اکثر مدل‌های رنگی به سمت سخت‌افزار (مانند مانیتور و چاپگر) یا کاربردهای تولید گرافیک و تصویر رنگی گرایش دارند{[}43{]}}.

\RL{مدل‌های رایج عبارتند از}:

\begin{itemize}
\item
  RGB \RL{(قرمز، سبز، آبی): مدل} RGB \RL{برای نمایشگرهای رنگی و دوربین‌های دیجیتال استفاده می‌شود. هر پیکسل از ترکیب نورهای قرمز، سبز و آبی تشکیل شده و مقدار هر رنگ عددی بین 0 تا 255 است. این مدل بر پایه سیستم مختصات کارتزین تعریف می‌شود و زیرفضای رنگی آن مکعبی است که سیاه در مبدأ و سفید در دورترین گوشه قرار دارد. این مدل برای پردازش تصویر و نمایش طبیعی رنگ‌ها بسیار پرکاربرد است}.
\item
  CMY \RL{(آبی فیروزه‌ای، بنفش، زرد): بیشتر برای چاپگرهای رنگی کاربرد دارد}.
\item
  : YIQ \RL{استاندارد تلویزیون رنگی، شامل لومینانس} (Y) \RL{و دو مؤلفه رنگی} I \RL{و} Q \RL{که به ترتیب هم‌فاز و متعامد هستند}.
\item
  HIS \footnote{(Hue\RL{،} Saturation\RL{،} Intensity)} \RL{: برای پردازش و تحلیل تصویر رنگی استفاده می‌شود}.
\end{itemize}

\RL{مدل} RGB \RL{یکی از مهم‌ترین مدل‌های رنگی برای پردازش تصویر دیجیتال به شمار می‌رود، زیرا هر تصویر شامل سه صفحه مستقل است که هر صفحه به یک رنگ اولیه اختصاص دارد و ترکیب این سه صفحه روی نمایشگر، تصویر نهایی رنگی را ایجاد می‌کند.}

\begin{enumerate}
\def\labelenumi{\arabic{enumi}.}
\item
  \protect\phantomsection\label{_Toc232899881}{}\RL{مولفه‌های مهم یک روش رمزنگاری تصویر}
\end{enumerate}

\RL{برای طراحی یک الگوریتم رمزنگاری تصویر مؤثر، توجه به چند ویژگی کلیدی ضروری است{[}44{]}}: \RL{}

\RL{کاهش اضافات و حجم داده: اضافات زیاد و حجم بالا در داده‌های تصاویر رمزشده، معمولاً باعث آسیب‌پذیری آنها در برابر تحلیلگران و حملات آماری می‌شود (50). حجم زیاد اطلاعات، امکان استخراج نمونه‌های کافی از متن رمزشده برای تحلیل آماری را فراهم می‌کند، به‌ویژه زمانی که الگوریتم از طول محدود و نسبتاً کوتاه استفاده کند. کاهش حجم داده‌ها و ساده‌تر کردن الگوریتم می‌تواند بار محاسباتی را کاهش دهد و امنیت را بهبود بخشد}. \RL{} \RL{ایجاد همبستگی ضعیف میان پیکسل‌ها: تحلیل آماری نشان داده است که پیکسل‌های مجاور، چه به صورت افقی، عمودی و چه قطری، معمولاً همبستگی زیادی دارند . این همبستگی شدید باعث کند شدن پراکندگی یا درهم‌ریزی پیکسل‌ها می‌شود. با تکرار الگوریتم و استفاده از مکانیزم‌های تصادفی‌ساز، می‌توان این همبستگی‌ها را کاهش داد و امنیت تصویر رمزشده را افزایش داد}. \RL{غیرقابل پیش‌بینی بودن تصویر رمزشده: براساس تئوری اطلاعات شانون، یک سیستم رمزنگاری ایمن باید آنتروپی بالایی داشته باشد و تصویر رمزشده نباید هیچ اطلاعاتی درباره تصویر اصلی فاش کند. برای رسیدن به این هدف، تصویر رمزشده باید تا حد امکان تصادفی باشد، هیستوگرام متوازن داشته باشد و هیچ همبستگی آماری بین پیکسل‌های مجاور وجود نداشته باشد. دستیابی به این سطح امنیت نیازمند چندین مرحله جابجایی و جانشانی پیچیده است}. \RL{} \RL{سرعت رمزنگاری و هزینه محاسباتی: یک الگوریتم ایمن که بسیار کند اجرا شود یا نیازمند هزینه محاسباتی بالا باشد، برای کاربردهای زمان واقعی ارزش عملی کمی خواهد داشت. بنابراین، سرعت بالا و مصرف منابع کم در طراحی الگوریتم اهمیت ویژه‌ای دارد}. \RL{} \RL{فشرده‌سازی داده‌ها همراه با امنیت: رمزنگاری تصویر معمولاً همزمان با فشرده‌سازی داده‌ها انجام می‌شود، زیرا تصاویر دیجیتال حجم بالایی دارند و اغلب دارای اطلاعات اضافی هستند. ادغام نیازهای امنیتی با فرایند فشرده‌سازی، چالشی مهم است که باید بدون کاهش کیفیت فشرده‌سازی، امنیت مناسبی را تضمین کند}. \RL{} \RL{تغییر تصویر به نحوی که قابلیت تشخیص نداشته باشد: بینایی انسان نسبت به تغییرات جزئی تصویر حساس نیست، بنابراین الگوریتم باید بتواند بیت‌های مرتبط و همبسته تصویر را به‌طور مؤثر رمز کند. در روش‌های سنتی، برخی شباهت‌ها حتی پس از رمزگذاری باقی می‌مانند و نیازمند توان محاسباتی بالا برای پوشش تمام بیت‌های تصویر هستند. الگوریتم بهینه باید این تغییرات را بدون افزایش غیرضروری هزینه محاسبات اعمال کند{[}45{]}.}

\begin{enumerate}
\def\labelenumi{\arabic{enumi}.}
\item
  \protect\phantomsection\label{_Toc232899882}{}\RL{نظریه‌ی آشوب و مفاهیم بنیادی آن}
\end{enumerate}

\RL{نظریه آشوب، گسترش خود را بیشتر مدیون پژوهش‌های دانشمندانی همچون هانری پوانکاره، ادوارد لورنز، بنوا مندلبروت و مایکل فیگنباوم است. پوانکاره نخستین فردی بود که اثبات کرد مسئله سه‌جسمی (مانند خورشید، زمین و ماه) یک مسئله آشوبی و غیرقابل‌حل دقیق است. شاخه‌ای از این نظریه که در مکانیک کوانتومی به‌کار می‌رود، «آشوب کوانتومی» نام دارد. گفته می‌شود که حتی پیر لَپلاس و عمر خیام نیز پیش از پوانکاره به وجود چنین پدیده‌هایی پی برده بودند. اولین آزمایش واقعی در زمینه آشوب توسط ادوارد لورنز، هواشناس آمریکایی، در سال 1961 انجام شد. وی هنگام کار بر روی مدل‌های عددی پیش‌بینی وضعیت آب‌وهوا دریافت که تغییرات بسیار جزئی در شرایط اولیه می‌تواند منجر به تغییرات چشمگیر در نتایج نهایی شود. این پدیده که بعدها با عنوان اثر پروانه‌ای}\footnote{Butterfly Effect} \RL{شناخته شد، نشان‌دهنده‌ی حساسیت بسیار بالای سیستم‌های آشوبی نسبت به شرایط اولیه است{[}46{]}.}

\RL{لورنز در سال 1972 مقاله‌ای با عنوان «آیا بال‌زدن پروانه‌ای در برزیل می‌تواند موجب وقوع گردبادی در تگزاس شود؟» منتشر کرد و در سال 1993 نیز کتاب معروف خود با عنوان} The Essence of Chaos \RL{را به چاپ رساند. در ادامه، بنوا مندلبروت در سال 1967 در جریان مطالعه‌ای پیرامون اندازه‌گیری طول سواحل بریتانیا، مفاهیم بعد فراکتال و خودهمانندی} \footnote{Self-Similarity} \RL{را مطرح نمود. او در سال 1975 نیز مقاله‌ی تأثیرگذاری با عنوان} The Fractal Geometry of Nature \RL{منتشر کرد که بنیان نظریه‌ی هندسه‌ی فراکتالی و بخش مهمی از نظریه‌ی آشوب را تشکیل داد. از سوی دیگر، مایکل فیگنباوم در سال 1977 مقاله‌ای تحت عنوان} Quantitative Universality for a Class of Nonlinear Transformations \RL{ارائه کرد که در آن با بررسی نگاشت لجستیک موفق به کشف پدیده‌ی جهان‌شمولی آشوب شد و نسبت‌های عددی خاصی را برای توصیف رفتار سامانه‌های غیرخطی به‌دست آورد{[}45{]}.}

\begin{enumerate}
\def\labelenumi{\arabic{enumi}.}
\item
  \protect\phantomsection\label{_Toc232899883}{}\RL{تعریف آشوب}
\end{enumerate}

\RL{در لغت، «آشوب» به معنای بی‌نظمی، درهم‌ریختگی و آشفتگی است و در علم مکانیک مترادف با واژه «تلاطم» به کار می‌رود. این واژه در ظاهر به نبود هرگونه نظم یا ساختار اشاره دارد و در زبان روزمره، نشانه‌ای از نابسامانی و بی‌سازمانی تلقی می‌شود. بااین‌حال، با گسترش دیدگاه‌های نوین علمی، مفهوم آشوب دیگر به معنای صرفِ بی‌نظمی و ناکارآمدی در نظر گرفته نمی‌شود؛ بلکه بیانگر نوعی رفتار غیرقابل پیش‌بینی در سامانه‌های پویاست که خود از نظمی درونی و پیچیده تبعیت می‌کند. به عبارتی، در دل بی‌نظمی، نوعی نظم نهفته وجود دارد که تنها در مقیاس‌های متفاوت آشکار می‌شود. پدیده‌ی آشوب معمولاً در نتیجه‌ی تعاملات پیچیده‌ی درونی اجزای یک سیستم و به‌ویژه در سامانه‌های غیرخطی و زیستی پدید می‌آید. در چنین سیستم‌هایی، تغییرات ناگهانی می‌تواند به رفتارهایی منجر شود که ماهیت تصادفی یا غیر‌دوره‌ای دارند و به آن‌ها رفتار آشوب‌گونه گفته می‌شود. بسیاری از این سیستم‌ها، هرچند رفتارشان ظاهراً نامنظم است، در واقع از قوانین دقیق و قطعی تبعیت می‌کنند؛ ازاین‌رو به آن‌ها آشوب معین}\footnote{Deterministic Chaos} \RL{گفته می‌شود}. \RL{سیستم‌های آشوبی در اصل، سیستم‌های دینامیکی غیرخطی هستند که حالت آن‌ها با گذر زمان تغییر می‌کند و آینده‌ی رفتارشان به‌شدت به شرایط اولیه وابسته است. اساس نظریه‌ی آشوب بر این فرض استوار است که در هر بی‌نظمی، نظمی پنهان وجود دارد. به بیان دیگر، هنگامی‌که جایگاه یا وضعیت بعدی یک سیستم به‌صورت دقیق قابل پیش‌بینی نباشد، آن سیستم را آشفته یا نامنظم می‌نامند. چنین سامانه‌ای هیچ‌گاه دوبار دقیقاً در یک حالت تکرار نمی‌شود}.

\RL{بااین‌حال، طبق نظریه‌ی آشوب، اگر رفتار چنین سیستمی در طول زمان به‌دقت مشاهده شود، می‌توان الگویی از نظم ذاتی را در پسِ حرکات تصادفی آن یافت. حتی سیستم‌هایی که در ظاهر بی‌قاعده و غیرقابل‌پیش‌بینی هستند، در محدوده‌ی معینی از فضا و زمان حرکت می‌کنند و از مرزهای مشخص خود فراتر نمی‌روند. بدین ترتیب، پدیده‌هایی که در مقیاس کوچک آشفته به نظر می‌رسند، ممکن است در مقیاس بزرگ‌تر رفتاری منظم و پایدار داشته باشند. در حقیقت، در دل هر بی‌نظمی و آشفتگی، نوعی الگوی پنهان از نظم وجود دارد که جلوه‌ای زیبا و قابل‌درک از پیچیدگی طبیعت را آشکار می‌سازد{[}46{]}.}

\begin{enumerate}
\def\labelenumi{\arabic{enumi}.}
\setcounter{enumi}{1}
\item
  \protect\phantomsection\label{_Toc232899884}{}\RL{خواص آشوب}
\end{enumerate}

\RL{سيگنال آشوب ظاهري شبيه به اغتشاش دارد ولي در عين حال كاملاً قطعي است. يعني با داشتن مقادير اوليه و تابع نگاشت مي توان دقيقاً همان مقادير را دوباره توليد كرد. هر سیستم بایستی حداقل سه شرط زیر را داشته باشد تا بتوان آن را آشوبناک نامید {[}21{]}:}

\begin{enumerate}
\def\labelenumi{\arabic{enumi}.}
\item
  \protect\phantomsection\label{_Toc232899885}{}\RL{حساسیت نسبت به شرایط اولیه}
\end{enumerate}

\RL{منظور از حساسيت نسبت به شرايط اوليه اين است كه هر تغیير جزئي در مقادير اوليه باعث ايجاد اختلاف فاحشي در مقادير بعدي تابع خواهد شد. به اين معني كه اگر مقادير اوليه سيگنال كمي تغيير كند سيگنال حاصل تفاوت بسياري با سيگنال اوليه خواهد داشت . اگر، دو سری زمانی با فرآیندهای آشوبی ولی با شرایط اولیه بسیار نزدیک به هم را درنظر بگیریم، مسیرهای زمانی آنها پس از مدتی متمایز شده به طور کامل به صورت دو سری زمانی متفاوت از یکدیگر به نظرخواهند رسید}.

\begin{enumerate}
\def\labelenumi{\arabic{enumi}.}
\setcounter{enumi}{1}
\item
  \protect\phantomsection\label{_Toc232899886}{}\RL{انتقال توپولوژیکی}
\end{enumerate}

\RL{منظور از انتقال توپولوژیکی این است که متغییرهای آشوبی تمام حالات را بدون تکرار در محدوده خاصی و بر طبق قوانین خود می‌پیمایند؛ اعداد تولیدی به خوبی در فضای حالت پخش می‌شوند و بعد از چند تکرار قسمت های مختلفی از فضا را بازدید می کنند. این ویژگی بیان می‌کند که متغیرهای آشوبگون تمام حالت‌ها را در محدوده‌ی خاصی به طور غیرتکراری با توجه به قوانین خود پیمایش می‌کنند. بنابراین، این ویژگی می‌تواند به عنوان یک مکانیزم بهینه‌سازی استفاده شود که تضمین می‌کند که هیچ راه‌حلی در فضای جستجو دو بار ملاقات نشود و مانع از قرار گرفتن الگوریتم در بهینه‌ی محلی می‌شود}.

\begin{enumerate}
\def\labelenumi{\arabic{enumi}.}
\setcounter{enumi}{2}
\item
  \protect\phantomsection\label{_Toc232899887}{}\RL{چگالی توپولوژیکی}
\end{enumerate}

\RL{منظور از چگالی توپولوژیکی این است که هر نقطه‌ای در فضا به طور دلخواه توسط مدارهای دوره‌ای قابل دسترسی است.}

\RL{علاوه بر این ویژگی‌ها خواص دیگر سیستم‌های آشوبی شامل غیرخطی بودن، رفتار ظاهراً تصادفی و عملکرد کاملاً قطعی هست. اعداد تولیدی توسط سیستم های آشوبناک کاملاً تصادفی به نظر می رسند و میتوان از این اعداد در کاربردهایی که نیاز به اعداد تصادفی باشد، استفاده کرد. روش هاي مورد استفاده براي توليد اعداد تصادفي در الگوريتم هاي بر مبناي توابع آشوب، اين امكان را به ما مي دهند كه در صورت داشتن مقدار اوليه و تابع نگاشت، همان اعداد تصادفي را دوباره باز توليد كنيم. در عين اينكه توابع آشوب ظاهري تصادفي دارند اما كاملاً قطعي هستند؛ يعني همواره با داشتن تابع نگاشت و مقادير اوليه مي¬توان يك مجموعه از مقادير را كه به ظاهر هيچ نظمي در توليد آنها وجود ندارد را توليد و دوباره همان مقادير را بازتوليد كرد.}

\begin{enumerate}
\def\labelenumi{\arabic{enumi}.}
\setcounter{enumi}{2}
\item
  \protect\phantomsection\label{_Toc232899888}{}\RL{نظریه آشوب}
\end{enumerate}

\RL{نظریه آشوب یا نظریه بی‌نظمی‌ها، به بررسی و تحلیل سیستم‌های دینامیکی آشوبناک می‌پردازد. این سیستم‌ها در واقع سامانه‌های دینامیکی غیرخطی هستند که رفتار آن‌ها به شدت به شرایط اولیه وابسته است؛ به‌طوری‌که حتی تغییرات بسیار جزئی در مقدار اولیه متغیرها می‌تواند به نتایج کاملاً متفاوتی در آینده منجر شود. این ویژگی شاخص، با عنوان مشهور «اثر پروانه‌ای}\footnote{Butterfly Effect}\RL{» شناخته می‌شود؛ پدیده‌ای که نخستین‌بار توسط ادوارد لورنز در دهه 1960 میلادی معرفی شد. رفتار سیستم‌های آشوبی در نگاه نخست تصادفی و غیرقابل پیش‌بینی به نظر می‌رسد، اما در واقعیت، چنین سیستم‌هایی از قوانین کاملاً قطعی و معین پیروی می‌کنند. به همین دلیل، آشوب را نمی‌توان ناشی از وجود عنصر تصادف دانست؛ بلکه این رفتار حاصل غیرخطی بودن روابط درونی سیستم و حساسیت بالا به شرایط اولیه است}. \RL{از دیدگاه نظری، وجود رفتار آشوب‌گونه در یک سیستم دینامیکی مستلزم آن است که سامانه حداقل دارای سه متغیر حالت مستقل باشد. به عبارت دیگر، تنها در سیستم‌های مرتبه سه یا بالاتر است که امکان ظهور آشوب وجود دارد. نمونه‌ای کلاسیک از چنین سیستمی، مدل لورنز است که برای توصیف رفتار همرفت در جوّ زمین ارائه شده است}. \RL{در مقابل، در سیستم‌های دینامیکی زمان‌گسسته، وجود تنها یک متغیر حالت نیز می‌تواند برای بروز رفتار آشوبی کافی باشد. یکی از نمونه‌های مشهور در این زمینه، نگاشت لجستیک است که برای مدل‌سازی تغییرات جمعیتی در اکوسیستم‌ها به‌کار می‌رود. این مدل نشان می‌دهد که چگونه یک رابطه ساده غیرخطی می‌تواند منجر به رفتار پیچیده، غیرتناوبی و آشوب‌گونه شود. {[}42-46{]}.}

\begin{enumerate}
\def\labelenumi{\arabic{enumi}.}
\setcounter{enumi}{3}
\item
  \protect\phantomsection\label{_Toc232899889}{}\RL{ویژگی‌های تئوری آشوب}
\end{enumerate}

\RL{تئوری آشوب به عنوان یکی از مفاهیم بنیادی در بررسی سیستم‌های دینامیکی غیرخطی، دارای ویژگی‌هایی است که آن را از سایر نظریه‌های کلاسیک متمایز می‌کند. در} \RL{ادامه به ویژگی های مهم} \RL{آن اشاره می‌شود}:

\textbf{\RL{الف) حساسیت به شرایط اولیه (اثر پروانه‌ای)}}

\RL{لورنتز در تحقیقات خود دریافت که تغییرات بسیار جزئی در شرایط اولیه‌ی معادلات پیش‌بینی‌کننده‌ی جوّی می‌تواند منجر به تفاوت‌های بسیار بزرگ در نتایج نهایی شود. این پدیده که به اثر پروانه‌ای معروف است، بیان می‌کند که یک تغییر کوچک در نقطه‌ی آغاز، می‌تواند پیامدهای گسترده و غیرقابل پیش‌بینی در آینده ایجاد کند. برای مثال، ممکن است بال زدن پروانه‌ای در کشور برزیل باعث تغییر در جریان‌های جوی شده و در نهایت موجب بروز طوفانی در تگزاس شود. در گذشته، سیستم‌هایی که چنین رفتارهایی از خود نشان می‌دادند، به عنوان سیستم‌های غیرقابل تحلیل از حوزه‌ی مطالعات علمی کنار گذاشته می‌شدند، زیرا پیش‌بینی رفتار آن‌ها غیرممکن به نظر می‌رسید. با این حال، امروزه این نوع سیستم‌ها به یکی از حوزه‌های اصلی پژوهش علمی تبدیل شده‌اند و نظریه‌ی آشوب به عنوان ابزاری برای مطالعه‌ی پدیده‌های ظاهراً تصادفی و غیرمنظم به‌کار گرفته می‌شود{[}41{]}.}

\RL{ب) سازگاری پویا} \footnote{Dynamic Adaptation}:

\RL{سیستم‌های آشوبناک برخلاف تصور اولیه، فاقد نظم یا قانون نیستند؛ بلکه دارای نوعی نظم درونی و پویایی سازگار هستند. این سیستم‌ها در تعامل با محیط خود، رفتاری مشابه موجودات زنده از خود نشان می‌دهند؛ بدین معنا که قادرند در برابر تغییرات محیطی واکنش نشان دهند و تعادل پویایی میان وضعیت درونی خود و محیط پیرامونشان برقرار کنند. به همین دلیل، چنین سیستم‌هایی نه تنها از بین نمی‌روند بلکه با تغییر شرایط، خود را با محیط جدید تطبیق می‌دهند {[}41{]}.}

\RL{پ) جذب‌کننده عجیب و پیچیده} \footnote{Strange and Complex Attractor}:

\RL{در تمامی سیستم‌های دینامیکی پایدار -- چه معمولی و چه آشوبی -- مسیرهای زمانی سیستم}\footnote{Time Series} \RL{در نهایت به سمت نقطه یا مجموعه‌ای خاص میل می‌کنند که به آن تعادل یا جذب‌کننده}\footnote{Attractor} \RL{گفته می‌شود. به عنوان نمونه، در یک معادله‌ی تفاضلی مرتبه اول، تمامی مسیرهای زمانی بدون توجه به مقدار اولیه، در نهایت به سمت نقطه‌ی تعادلی در مرکز مختصات همگرا می‌شوند. در این حالت، مرکز مختصات همان جذب‌کننده سیستم محسوب می‌شود}.

\RL{با این حال، در بسیاری از سیستم‌های غیرخطی، جذب‌کننده‌ها ممکن است پیچیده‌تر از یک نقطه باشند. برای مثال، جذب‌کننده می‌تواند شامل چندین ناحیه یا مسیر بسته باشد که تمام مسیرهای زمانی سیستم در نهایت به آن‌ها گرایش یافته و در محدوده‌ی آن‌ها باقی بمانند}.

\RL{در سیستم‌های آشوبی، جذب‌کننده‌ها بسیار پیچیده و غیرمنتظره هستند و به همین دلیل به آن‌ها جذب‌کننده‌های عجیب}\footnote{Strange Attractors} \RL{گفته می‌شود. این نوع جذب‌کننده‌ها معمولاً دارای ساختاری فراکتالی بوده و از بی‌نهایت نقاط غیرقابل شمارش تشکیل شده‌اند، به‌گونه‌ای که مسیرهایی که درون آن‌ها آغاز می‌شوند، همواره در درون آن باقی می‌مانند و مسیرهای نزدیک نیز به‌تدریج به سمت آن جذب می‌شوند{[}46{]}.}

\RL{مسیرهای زمانی موجود در جذب‌کننده‌های عجیب ممکن است غیرقابل تکرار باشند یا به طور دلخواه در چرخه‌هایی از پیش تعیین‌شده تکرار شوند. این ویژگی باعث می‌شود سیستم‌های آشوبی رفتاری ظاهراً تصادفی اما درواقع دارای نظم پنهان از خود نشان دهند}.

\RL{در برخی از سیستم‌های دینامیکی غیرخطی که دارای جذب‌کننده‌های پیچیده هستند، امکان بازسازی مسیر طی‌شده وجود دارد. این دسته از سیستم‌ها که به سیستم‌های هامیلتونی}\footnote{Hamiltonian Systems} \RL{معروف‌اند، در صورتی که شرایط اولیه و معادلات حرکت در دست باشد، مسیر حرکت به سمت جذب‌کننده را می‌توان در جهت معکوس دنبال کرد. به عبارت دیگر، این سیستم‌ها از نظر زمانی قابل بازگشت}\footnote{Time Reversible} \RL{هستند}.

\RL{اما در مقابل، در بیشتر مدل‌های آشوبی، حتی با در دست داشتن شرایط اولیه و معادلات، بازسازی دقیق مسیر حرکتی ممکن نیست، زیرا بی‌نهایت مسیر انتقالی متفاوت وجود دارد که به جذب‌کننده‌های مختلف منتهی می‌شوند. این غیرقابل بازگشت بودن زمانی از ویژگی‌های اصلی سیستم‌های آشوبی محسوب می‌شود ت) خودتشابهی} \footnote{Self-Similarity}:

\RL{یکی از ویژگی‌های بنیادین نظریه آشوب، خودتشابهی یا خودهمانندی است. این خاصیت بیان می‌کند که در بسیاری از سیستم‌های آشوبی، الگوی کلی رفتار سیستم در مقیاس‌های مختلف تکرار می‌شود؛ به عبارت دیگر، هر جزء کوچک از سیستم، تصویری کوچک‌شده اما مشابه از کل آن است.}

\RL{در حوزه‌ی ریاضیات و فیزیک، این مفهوم با ساختارهای فراکتالی}\footnote{Fractals} \RL{مرتبط است که در آن‌ها الگوها در هر مقیاسی مشابه‌اند و اگر بخشی از شکل را بزرگ‌نمایی کنیم، همان ساختار کلی در آن دیده می‌شود. این ویژگی باعث می‌شود که سیستم‌های آشوبی، علی‌رغم پیچیدگی ظاهری، از نظم درونی و الگوی تکرارشونده‌ای برخوردار باشند}.

\RL{در سطح مفهومی و کاربردی، خودتشابهی را می‌توان به پدیده‌های انسانی و سازمانی نیز تعمیم داد. به‌گونه‌ای که در یک سازمان، رفتار و جهت‌گیری هر جزء (فرد یا واحد) می‌تواند بازتابی از ساختار کلی و هدف نهایی سازمان باشد. به بیان دیگر، وجود خودتشابهی در رفتار اعضای یک مجموعه موجب انسجام و وحدت هدف در کل سیستم می‌شود، به‌طوری‌که تمام اجزا در راستای یک جهت و هدف مشترک عمل می‌کنند {[}27{]}.}

\textbf{\RL{سیستم‌های آشوب و نمای لیاپانوف}}\footnote{\textbf{Lyapunov Exponent}} \textbf{:}

\RL{یکی از روش‌های اساسی برای تشخیص وجود آشوب در سیستم‌های دینامیکی، محاسبه‌ی نمای لیاپانوف است. این شاخص، میانگین نرخ واگرایی یا همگرایی دو مسیر بسیار نزدیک به هم در فضای فاز سیستم را نشان می‌دهد. در واقع، نمای لیاپانوف تعیین می‌کند که اگر دو نقطه ابتدایی بسیار نزدیک انتخاب شوند، با گذشت زمان چگونه فاصله‌ی بین آن‌ها تغییر می‌کند.}

\RL{مفهوم نمای لیاپانوف نخستین‌بار توسط الکساندر لیاپانوف، ریاضیدان روسی، در سال 1892 میلادی برای بررسی پایداری معادلات دیفرانسیل غیرخطی ارائه شد. این روش به تحلیل پایداری سیستم‌ها بدون نیاز به حل تحلیلی کامل معادلات کمک می‌کند}.

\RL{برای سیستم‌های دینامیکی غیرخطی، رفتار آن‌ها را می‌توان با استفاده از نگاشت‌های \textbf{بازگشتی}}\footnote{\textbf{Mappings}} \RL{تحلیل کرد. نمای لیاپانوف به‌صورت عددی از طریق دو روش اصلی، یعنی \textbf{تابع معین} و \textbf{بعد جانشانی}، قابل محاسبه است}.

\begin{enumerate}
\def\labelenumi{\arabic{enumi}.}
\setcounter{enumi}{4}
\item
  \protect\phantomsection\label{_Toc232899890}{}\RL{نمونه‌هایی از سیستم‌های آشوب}
\end{enumerate}

\RL{در ادامه چند نمونه از سیستم‌های آشوبی معرفی می‌شوند}:

\begin{enumerate}
\def\labelenumi{\arabic{enumi}.}
\item
  \protect\phantomsection\label{_Toc232899891}{}\RL{نگاشت لجستیک}
\end{enumerate}

\RL{نگاشت لجستیک یک نگاشت یک‌بعدی است که برای توصیف نوسانات جمعیت برخی گونه‌های حشرات ارائه شد}:

\[x_{\left\{ n + 1 \right\}}\  = \ r\ x_{n}\ (1\  - \ x_{n})\]

\RL{رابطه (1-2) جمعیت در سال} n+1 \RL{} \RL{را بر اساس جمعیت سال} n \RL{نشان می‌دهد. برای} 0 \textless{} r \textless{} \RL{، تمام مسیرها، از هر نقطه شروع که شده باشند، پس از چند تکرار به صفر می‌رسند. برای} 1\textless r\textless3 \RL{، مسیرها پس از چند تکرار به مقدار} 1-1/r \RL{} \RL{میل می‌کنند (شکل 2-2). از نقطه} r=3.82 \RL{به بعد، رفتار آشوبی در نگاشت لجستیک به خوبی مشاهده می‌شود{[}5{]}.}

\includegraphics[width=3.75866in,height=2.46688in]{media/image2.png}

\begin{enumerate}
\def\labelenumi{\arabic{enumi}.}
\item
  \protect\phantomsection\label{_Toc232899805}{}\RL{نگاشت لجستیک}
\end{enumerate}

\begin{enumerate}
\def\labelenumi{\arabic{enumi}.}
\setcounter{enumi}{1}
\item
  \protect\phantomsection\label{_Toc232899892}{}\RL{سیستم آشوبی لورنز}\footnote{Lorenz system}
\end{enumerate}

\RL{ادوارد لورنز، دانشمند هواشناس، در سال 1963 یک سیستم سه‌معادله‌ای برای مدل‌سازی جابه‌جایی دما به صورت همرفتی در اتمسفر ارائه داد (رابطه 2-2 و شکل 3-2)}:

\[\left\{ \begin{array}{r}
\frac{dx}{dt} = \sigma(y - x) \\
\frac{dy}{dt} = Rx - y - xz \\
\frac{dz}{dt} = xy - \beta z
\end{array} \right.\ \] \RL{(۲-۲)}

\RL{رابطه (۲-۲) سیستم لورنز است که با مقداردهی به ثابت‌های σ ، R و β به عنوان یک سیستم آشوبی، در الگوریتم‌های رمزنگاری کاربرد دارد. این مدل امکان شبیه‌سازی ناپایداری‌ها و پیش‌بینی‌های غیرقطعی در هواشناسی را فراهم می‌کند و به دلیل خواص دینامیکی خاص، گزینه مناسبی برای الگوریتم‌های رمزنگاری محسوب می‌شود{[}15-16{]}.}

\[\frac{dx}{dt}\  = \ 10(y\  - \ x)\]

\[\frac{dy}{dt}\  = \ 28x\  - \ y\ –\ xz\]

\[\frac{dz}{dt} = xy\  - \frac{8}{3}z\]

\RL{}\includegraphics[width=4.80875in,height=2.45855in]{media/image3.png}

\begin{enumerate}
\def\labelenumi{\arabic{enumi}.}
\item
  \RL{نمای سه‌بعدی جذب‌کننده آشوبی لورنز {[}23{]}.}
\end{enumerate}

\begin{enumerate}
\def\labelenumi{\arabic{enumi}.}
\setcounter{enumi}{2}
\item
  \protect\phantomsection\label{_Toc232899893}{}\RL{سیستم آشوبی چن}
\end{enumerate}

\RL{این سیستم در سال 1991 معرفی شد و با معادلات رابطه (۲-۳) نمایش داده می‌شود {[}41{]}:}

\[\left\{ \begin{array}{r}
\overset{˙}{x} = a(y - x) \\
\overset{˙}{y} = (c - a)x - xz + cy \\
\overset{˙}{z} = xy - bz
\end{array} \right.\ \] \RL{(۲-۳)}

\RL{با مقادیر اولیه} a=35 \RL{،} b=3 \RL{،} c=28 \RL{، خروجی سیستم به شکل (4-2) ترسیم می‌شود. لازم به ذکر است که تغییر هر یک از پارامترهای فوق می‌تواند رفتار سیستم را از حالت آشوبی به غیرآشوبی تغییر دهد}.

\includegraphics[width=5.17545in,height=2.77524in]{media/image4.png}

\begin{enumerate}
\def\labelenumi{\arabic{enumi}.}
\item
  \RL{نمای سه‌بعدی جذب‌کننده آشوبی سیستم چن {[}17{]}.}
\end{enumerate}

\begin{enumerate}
\def\labelenumi{\arabic{enumi}.}
\setcounter{enumi}{1}
\item
  \protect\phantomsection\label{_Toc232899894}{}\RL{ارتباط بین آشوب و رمزنگاری}
\end{enumerate}

\RL{بین سیستم‌های رمزنگاری و سیستم‌های آشوبناک رابطه‌ای بسیار نزدیک وجود دارد و این ارتباط باعث شده است که تئوری آشوب در طراحی رمزنگاری عملی به طور گسترده‌ای مورد استفاده قرار گیرد. نگاشت‌های آشوبناک با رفتار دینامیکی بسیار پیچیده، کاملاً قطعی و مشخص هستند و می‌توان آن‌ها را به راحتی با معادلات تفاضلی گسسته مدل‌سازی کرد. به همین دلیل، کارایی رمزنگاری آشوبناک در رمزنگاری تصویر، با توجه به حجم بالای تصاویر، همبستگی پیکسل‌های مجاور و افزونگی موجود در آن، از روش‌های متداول متقارن معمولاً بهتر است{[} 36{]}.}

\RL{خصوصیات دینامیکی نگاشت‌های آشوبناک با ویژگی‌های مورد نیاز یک سیستم رمزنگاری، مانند حساسیت به کلید و متن اصلی، هم‌خوانی زیادی دارد و این مسئله، کاربرد آشوب در رمزنگاری تصاویر را توجیه می‌کند {[}19{]}. شباهت‌های این دو سیستم در جدول ۲-۲ ارائه شده است}.

\RL{در عمل، اغتشاش در سیستم‌های رمزنگاری متداول از طریق تبدیل متن اصلی به متن تصادفی ایجاد می‌شود، به گونه‌ای که هیچ الگوی تکراری در متن رمز قابل مشاهده نباشد. مشابه این، مسیر حرکت سیستم‌های آشوبناک در فضای حالت معمولاً توزیع یکنواخت دارد و پیش‌بینی موقعیت نهایی یک نقطه بر اساس موقعیت اولیه آن بسیار دشوار است{[}23{]}.}

\begin{enumerate}
\def\labelenumi{\arabic{enumi}.}
\item
  \protect\phantomsection\label{_Toc232899944}{}\textbf{\RL{شباهت‌های سیستم‌های رمز متداول و سیستم‌های آشوبناک}}
\end{enumerate}

{\def\LTcaptype{none} % do not increment counter
\begin{longtable}[]{@{}
  >{\raggedright\arraybackslash}p{(\linewidth - 2\tabcolsep) * \real{0.4015}}
  >{\raggedright\arraybackslash}p{(\linewidth - 2\tabcolsep) * \real{0.2269}}@{}}
\toprule\noalign{}
\begin{minipage}[b]{\linewidth}\centering
\textbf{\RL{سیستم‌های آشوبناک}}
\end{minipage} & \begin{minipage}[b]{\linewidth}\centering
\textbf{\RL{سیستم‌های رمز متداول}}
\end{minipage} \\
\midrule\noalign{}
\endhead
\bottomrule\noalign{}
\endlastfoot
\RL{اغتشاش} (Ergodicity) & \RL{اغتشاش} (Confusion) \\
\RL{حساسیت به شرایط اولیه و پارامترهای سیستم} & \RL{انتشار} (Diffusion) \\
\RL{شرایط اولیه و پارامترهای سیستم} & \RL{کلید رمز} \\
\RL{تکرار} & \RL{دور} (Cycle) \\
\end{longtable}
}

\RL{اصل مهم دیگر در طراحی یک سیستم رمزگذاری امن، ویژگی \textbf{انتشار} است. این خاصیت بیان می‌کند که تغییر تنها یک بیت در کلید یا متن اصلی، منجر به تولید یک متن رمز کاملاً متفاوت می‌شود. به همین ترتیب، سیستم‌های آشوبناک شدیداً به شرایط اولیه و پارامترهای خود وابسته‌اند و کوچک‌ترین تغییر در نقطه شروع یا پارامترها، باعث انحراف قابل توجه مسیر سیستم می‌شود{[}23{]}.}

\RL{تکنیک‌های اصلی رمزنگاری یک بلوک از داده‌ها شامل \textbf{اغتشاش و انتشار} هستند. اغتشاش باعث ایجاد ابهام بین متن اصلی و متن رمز شده می‌شود و انتشار، تغییرات را در سراسر متن رمز پخش می‌کند. جانشینی که در آن یک نماد با نماد دیگری جایگزین می‌شود، ساده‌ترین نوع اغتشاش است. جایگشت که ترتیب عناصر یک بلوک را تغییر می‌دهد، ساده‌ترین روش انتشار محسوب می‌شود}.

\RL{در سیستم‌های رمزنگاری، جابه‌جایی و جایگزینی متن اصلی به طور مکرر انجام می‌شود تا سطح امنیت مطلوب حاصل شود. هر بار تکرار این عملیات را یک \textbf{دور} می‌نامند. به شکل مشابه، در سیستم‌های دینامیکی، تکرار معادلات تفاضلی در فضای فاز انجام می‌شود. بنابراین، می‌توان از تئوری آشوب برای رمزنگاری بهره برد؛ به گونه‌ای که \textbf{پارامترها و شرایط اولیه سیستم نقش کلیدی کلید رمز را ایفا کنند} و تکرار معادله آشوبناک معادل دورهای عملکرد تابع رمز باشد} \RL{{[}33{]}.}

\begin{enumerate}
\def\labelenumi{\arabic{enumi}.}
\setcounter{enumi}{2}
\item
  \protect\phantomsection\label{_Toc232899895}{}\RL{پیشینه پژوهش}
\end{enumerate}

\RL{مطالعات انجام‌شده در زمینه رمزنگاری تصاویر رنگی با استفاده از سیستم‌های آشوبی و الگوریتم‌های بهینه‌سازی در دهه‌های اخیر به‌طور گسترده‌ای مورد توجه قرار گرفته است. در ادامه، به خلاصه‌ای از تحقیقات داخلی و خارجی مرتبط با موضوع اشاره می‌شود}:

\RL{ژانگ و لیو{[}3{]} در مطالعه خود ترکیب سیستم‌های آشوبی} Henon Map \RL{و} Tent Map \RL{را برای رمزنگاری تصاویر رنگی پیشنهاد دادند. نتایج این تحقیق نشان داد که این ترکیب، امنیت بیشتری نسبت به استفاده از یک سیستم آشوبی منفرد ارائه می‌دهد و الگوریتم پیشنهادی مقاومت خوبی در برابر حملات آماری دارد.}

\RL{هو و ژو {[}4{]} یک مدل آشوبی نمایی (}ECM\RL{) برای تولید نگاشت‌های آشوبی یک‌بعدی با آشوب پایدار معرفی می‌کنند. برخلاف بسیاری از سیستم‌های آشوبی که پایداری آشوب را در نظر نمی‌گیرند،} ECM \RL{با ترکیب دو سیستم آشوبی یک‌بعدی، نگاشت‌های جدیدی ایجاد می‌کند که در برابر تغییرات پارامتری مقاوم‌اند. در این مقاله، نه سیستم آشوبی مبتنی بر} ECM \RL{ارائه شده و رفتار آن‌ها از طریق دیاگرام انشعاب و تحلیل آشوب پایدار بررسی شده است. نتایج ارزیابی عملکرد نشان می‌دهد که این نگاشت‌ها در فضای پارامتری گسترده‌ای آشوب پایدار دارند. همچنین در کاربردهای عملی، این نگاشت‌ها در ارتباطات امن استفاده شده و شبیه‌سازی‌ها، عملکرد برتری را در برابر نویز نشان می‌دهند.}

\RL{آلوارز و همکاران {[}6{]} به بررسی رمزنگاری تصاویر با استفاده از سیستم آشوبی لورنز می‌پردازند. با توجه به تبادل گسترده تصاویر حاوی اطلاعات مالی و شخصی در فضای اینترنت، حفاظت از این داده‌ها اهمیت ویژه‌ای دارد. الگوریتم‌های رمزنگاری سنتی عمدتاً برای داده‌های متنی طراحی شده‌اند، در حالی که داده‌های چندرسانه‌ای، به‌ویژه تصاویر، کمتر مورد توجه قرار گرفته‌اند. در این پژوهش، روشی برای رمزنگاری و رمزگشایی تصاویر بر پایه سیستم آشوبی لورنز ارائه شده است. در این روش، ابتدا پیکسل‌های تصویر با استفاده از فرآیندهای جابه‌جایی و چرخش معکوس تخریب می‌شوند تا تصادفی‌بودن افزایش یابد. سپس، الگوریتم پیشنهادی با استفاده از معادلات سه‌بعدی سیستم لورنز اجرا می‌شود. نتایج حاصل نشان می‌دهند که میزان آنتروپی تصویر پس از رمزنگاری از 7.285 به 7.9974 افزایش یافته است. همچنین، معیارهای} NPCR \RL{و} UACI \RL{به ترتیب \lr{99.65٪} و \lr{30.35٪} محاسبه شده‌اند که نشان‌دهنده قابلیت بالای الگوریتم پیشنهادی در تأمین امنیت اطلاعات تصویری است. مقایسه این روش با سایر تحقیقات مشابه نیز برتری آن را تأیید می‌کند.}

\RL{خدادادی و زندوکیلی {[}7{]} روشی نوین برای رمزنگاری تصاویر رنگی با استفاده از ترکیب سیستم‌های آشوبی ارائه داده‌اند که باعث افزایش کارایی و استحکام فرآیند رمزنگاری می‌شود. الگوریتم پیشنهادی سه سری داده در بازه 0 تا 255 را با استفاده از سیستم آشوبی چن}\footnote{Chen} \RL{تولید می‌کند. سپس، یک سیستم چن دیگر با مقادیر اولیه متفاوت آغاز می‌شود که داده‌ها را به سه سری مقادیر بین 0 تا 10 تبدیل می‌کند. در این روش، مقادیر قرمز، سبز و آبی}\footnote{RGB} \RL{تصویر با داده‌های تولیدشده از سیستم اول ترکیب شده و برای رمزنگاری اولین پیکسل تصویر استفاده می‌شود. مقادیر تولیدشده از سیستم دوم نیز برای تغییر ترتیب ترکیب مقادیر سیستم اول با پیکسل‌های تصویر به کار گرفته می‌شوند. این فرآیند تا رمزنگاری تمامی پیکسل‌ها ادامه می‌یابد. نوآوری اصلی این روش در استفاده ترکیبی از دو سیستم آشوبی است که پیچیدگی فرآیند رمزنگاری را افزایش می‌دهد}. \RL{نتایج آزمایش‌ها بر روی تصاویر استاندارد (مجموعه داده‌های} \RL{}USC\RL{)}\footnote{https://libraries.usc.edu/databases} \RL{کارایی و استحکام این روش را نشان می‌دهد.}

\RL{پژوهش وانگ و لی {[}8{]} نشان داد که استفاده از سیستم آشوبی} Logistic Map \RL{در ترکیب با جایگشت و انتشار، می‌تواند یک الگوریتم رمزنگاری ایمن برای تصاویر ایجاد کند. با این حال، نیاز به بهبود کارایی محاسباتی الگوریتم برای استفاده در کاربردهای عملی‌تر، همچنان یک چالش باقی می‌ماند. این محدودیت می‌تواند انگیزه‌ای برای پژوهش‌های آینده باشد تا الگوریتم‌های رمزنگاری کارآمدتر و سریع‌تری طراحی شوند.}

\RL{در پژوهش دیگری از یک سیستم رمزنگاری مبتنی بر آشوب برای رمزگذاری تصاویر رنگی استفاده می‎شود {[}9{]}. در این روش ابتدا تصویر ورودی به کانال‎های قرمز، سبز و آبی تفکیک شده و هر کانال به‎صورت جداگانه پردازش می‎شود. کلیدهای رمزنگاری با استفاده از نقشه‎های آشوبی غیرخطی تولید می‎شوند که حساسیت بالایی به شرایط اولیه دارند و کوچکترین تغییر در پارامترهای اولیه، خروجی کاملاً متفاوتی ایجاد می‎کند. پیکسل‎های هر کانال ابتدا با دنباله‎های آشوبی جابه‎جا شده} \RL{و سپس مقادیر آنها با عملیات} XOR \RL{و جایگزینی دینامیکی تغییر داده می‎شود. در نهایت کانال‎های اصلاح‎شده مجدداً ترکیب شده و تصویر رمزگذاری شده نهایی حاصل می‎شود. این سیستم با ایجاد توزیع یکنواخت در هیستوگرام و کاهش همبستگی پیکسل‎های مجاور امنیت بالایی در برابر حملات آماری ارائه می‎دهد}.

\RL{لی و همکاران {[}10{]} با استفاده از سیستم} \RL{آشوبی نمائی} LEL \RL{، دنباله‌های تصادفی با پیچیدگی بالا تولید کردند. نتایج پژوهش نشان داد که سیستم} LEL \RL{در مقایسه با سیستم‌های آشوبی رایج، قابلیت بالاتری در تولید توالی‌های غیرقابل پیش‌بینی دارد و امنیت رمزنگاری را افزایش می‌دهد.}

\RL{صادقیان {[}11{]} در پژوهشی به بررسی روش جدیدی برای رمزنگاری تصاویر دیجیتال بر اساس ترکیب سیستم‌های آشوبی، الگوریتم ژنتیک و الگوریتم بهینه‌سازی ازدحام ذرات} \footnote{Particle swarm optimization (PSO)}\RL{پرداخته است. با توجه به ویژگی‌های خاص داده‌های تصویری، مانند حجم بالا و همبستگی زیاد میان پیکسل‌های همجوار، استفاده از الگوریتم‌های رمزنگاری متنی ناکارآمد است. در این روش، از سیستم‌های آشوبی برای ایجاد حساسیت به مقادیر اولیه و تولید داده‌های شبه‌تصادفی استفاده شده است. همچنین، ترکیب الگوریتم ژنتیک و} PSO \RL{به منظور بهینه‌سازی پارامترهای آشوبی و جلوگیری از جستجوی راکد به‌کار گرفته شده است. نتایج به‌دست‌آمده از آزمایش‌های انجام‌شده بر روی مجموعه‌ای از تصاویر استاندارد در نرم‌افزار متلب نشان می‌دهد که روش پیشنهادی در مقایسه با روش‌های مشابه، عملکرد بهتری در حفظ جزئیات تصویری و مقاومت در برابر حملات دارد.}

\RL{در تحقیق دیگری، یک الگوریتم رمزنگاری تصویر سریع و ایمن مبتنی بر فوق آشوب و عملیات برداری ارائه شده است {[}12{]}. روش پیشنهادی با کاهش تعداد تکرارهای سیستم فوق آشوب و به‌کارگیری انتخاب‌گر تصادفی، ماتریس کلیدی با تصادفی‌بودن بالا ایجاد می‌کند. همچنین، انتشار مبتنی بر زنجیره‌سازی رمز بلوکی با استفاده از عملیات برداری، رمزنگاری موازی تصویر را بهبود می‌بخشد. نتایج ارزیابی نشان می‌دهند که الگوریتم پیشنهادی تمامی آزمون‌های} NIST SP 800-22 \RL{را با موفقیت گذرانده، مقدار آنتروپی آن 7.999 بوده و در برابر حملات مختلف ایمن است. همچنین، با میانگین زمان رمزنگاری 0.023 ثانیه، نیازهای انتقال محرمانه تصاویر را برآورده می‌کند.}

\RL{در سال 2024 یک روش نوین رمزنگاری تصویر رنگی بر پایه ترکیب سه سیستم آشوبی} Lorenz\RL{،} Henon \RL{و 2}D-Logistic \RL{طراحی شده است {[}13{]}. این الگوریتم به‌صورت مستقل برای هر کانال رنگی عملیات جایگشت و انتشار انجام می‌دهد و با تحلیل‌هایی از جمله آنتروپی شانون،} NPCR \RL{و} UACI \RL{نشان داده شد که در برابر حملات آماری و تحلیل کلید مقاومت بالایی دارد.}

2-6-1. مقایسه تطبیقی پژوهش‌های مرجع

با توجه به مرور پژوهش‌های فوق، روش‌های رمزنگاری تصویر رنگی مبتنی بر آشوب را می‌توان از منظر سیستم(های) آشوبی به‌کاررفته، سازوکار جایگشت و انتشار، نوآوری الگوریتمی و معیارهای امنیتی گزارش‌شده مقایسه کرد. جدول ۲-۳ خلاصه‌ای تطبیقی از مهم‌ترین مطالعات مرتبط با موضوع پژوهش حاضر ارائه می‌دهد. مقادیر عددی آنتروپی، NPCR و UACI تنها در صورتی درج شده‌اند که در مقاله مرجع یا در ارزیابی مشابه گزارش شده باشند؛ در غیر این صورت با «---» مشخص شده است.

همان‌طور که جدول نشان می‌دهد، بسیاری از روش‌های پیشین از یک یا دو سیستم آشوبی ثابت (مانند Logistic، Tent یا Lorenz) بهره می‌برند و انتخاب سیستم رمزنگاری در سراسر تصویر یکسان است. در مقابل، روش پیشنهادی حاضر با ترکیب سیستم سه‌بعدی Chen برای جایگشت و نه سیستم آشوبی نمایی (ECM) با انتخاب پویا در سطح هر پیکسل، لایه‌ای اضافی از پیچیدگی و غیرقابل‌پیش‌بینی بودن ایجاد می‌کند. از نظر آنتروپی شانون، الگوریتم پیشنهادی به مقدار ایده‌آل ۸.۰ بیت در تمام کانال‌های رنگی دست یافته که نسبت به اکثر روش‌های مرجع (حدود ۷.۹۹۷) برتری جزئی اما معنادار دارد و نشان‌دهنده توزیع کاملاً یکنواخت مقادیر پیکسل پس از رمزنگاری است.

در معیار NPCR (حساسیت به کلید)، روش‌های مبتنی بر Lorenz، Chen+Logistic و فوق‌آشوب مقادیر نزدیک به \lr{۹۹.۶٪} گزارش کرده‌اند. میانگین NPCR روش پیشنهادی (\lr{۹۵.۱۳٪}) در محدوده قابل قبول قرار دارد، اما برای تصویر Baboon مقدار پایین‌تری مشاهده می‌شود که می‌تواند به حساسیت بیشتر سیستم‌های پیوسته Chen نسبت به شرایط اولیه در تصاویر با بافت پیچیده نسبت داده شود. از این رو، اگرچه الگوریتم پیشنهادی از نظر آنتروپی و حذف همبستگی مکانی عملکرد برجسته‌ای دارد، بهبود یکنواختی NPCR در تمام تصاویر آزمایشی از جمله موارد پیشنهادی برای تحقیقات آینده است.

در معیار UACI، روش خدادادی و زندوکیلی {[}۷{]}، وانگ و لی {[}۸{]} و الگوریتم فوق‌آشوب {[}۱۲{]} به مقادیر نزدیک به \lr{۳۳.۴۶٪} (مقدار ایده‌آل) دست یافته‌اند. میانگین UACI روش پیشنهادی (\lr{۳۰.۱۴٪}) اندکی پایین‌تر است؛ با این حال برای تصویر Peppers مقدار \lr{۳۲.۸۸٪} به ایده‌آل نزدیک است. در مقایسه با روش لورنز {[}۶{]} که UACI مشابه (\lr{۳۰.۳۵٪}) گزارش کرده، عملکرد روش پیشنهادی قابل رقابت است.

از منظر کارایی، الگوریتم فوق‌آشوب {[}۱۲{]} با زمان رمزنگاری ۰.۰۲۳ ثانیه سرعت بسیار بالایی دارد؛ این برتری عمدتاً ناشی از برداری‌سازی و کاهش تکرارهای آشوبی است. پیاده‌سازی فعلی روش پیشنهادی در Python با حلقه وابسته به انتخاب پویای ECM برای هر پیکسل، زمان رمزنگاری حدود ۲۳--۲۵ ثانیه برای تصویر ۵۱۲×۵۱۲ ایجاد می‌کند که برای اثبات مفهوم مناسب است اما در کاربرد بلادرنگ نیازمند بهینه‌سازی خواهد بود.

به‌طور خلاصه، نوآوری اصلی پژوهش حاضر نسبت به ادبیات، «انتخاب پویا و مستقل یکی از نه سیستم نمایی برای هر پیکسل» در کنار «جایگشت مبتنی بر سیستم پیوسته Chen» و «دستیابی به آنتروپی و هیستوگرام کاملاً یکنواخت» است. این ترکیب، روش پیشنهادی را از روش‌های تک‌آشوبی و حتی روش‌های دوگانه (Chen+Logistic) متمایز می‌سازد؛ هرچند در NPCR و UACI هنوز فاصله‌ای با بهترین روش‌های مرجع وجود دارد که در فصل‌های چهارم و پنجم با جزئیات بیشتر مورد بحث قرار گرفته است.

\begin{enumerate}
\def\labelenumi{\arabic{enumi}.}
\item
  جدول ۲-۳. مقایسه تطبیقی روش‌های رمزنگاری تصویر رنگی مبتنی بر آشوب
\end{enumerate}

{\def\LTcaptype{none} % do not increment counter
\begin{longtable}[]{@{}
  >{\centering\arraybackslash}p{(\linewidth - 14\tabcolsep) * \real{0.1250}}
  >{\centering\arraybackslash}p{(\linewidth - 14\tabcolsep) * \real{0.1250}}
  >{\centering\arraybackslash}p{(\linewidth - 14\tabcolsep) * \real{0.1250}}
  >{\centering\arraybackslash}p{(\linewidth - 14\tabcolsep) * \real{0.1250}}
  >{\centering\arraybackslash}p{(\linewidth - 14\tabcolsep) * \real{0.1250}}
  >{\centering\arraybackslash}p{(\linewidth - 14\tabcolsep) * \real{0.1250}}
  >{\centering\arraybackslash}p{(\linewidth - 14\tabcolsep) * \real{0.1250}}
  >{\centering\arraybackslash}p{(\linewidth - 14\tabcolsep) * \real{0.1250}}@{}}
\toprule\noalign{}
\begin{minipage}[b]{\linewidth}\centering
\textbf{مرجع}
\end{minipage} & \begin{minipage}[b]{\linewidth}\centering
\textbf{سیستم(های) آشوبی}
\end{minipage} & \begin{minipage}[b]{\linewidth}\centering
\textbf{جایگشت}
\end{minipage} & \begin{minipage}[b]{\linewidth}\centering
\textbf{انتشار}
\end{minipage} & \begin{minipage}[b]{\linewidth}\centering
\textbf{ویژگی متمایز}
\end{minipage} & \begin{minipage}[b]{\linewidth}\centering
\textbf{آنتروپی}
\end{minipage} & \begin{minipage}[b]{\linewidth}\centering
\textbf{NPCR (\%)}
\end{minipage} & \begin{minipage}[b]{\linewidth}\centering
\textbf{UACI (\%)}
\end{minipage} \\
\midrule\noalign{}
\endhead
\bottomrule\noalign{}
\endlastfoot
ژانگ و لیو {[}۳{]} & Henon + Tent & بله & بله & ترکیب دو نگاشت گسسته & --- & --- & --- \\
هو و ژو {[}۴{]} & ۹ ECM & --- & --- & پایداری آشوب نمایی & --- & --- & --- \\
آلوارز {[}۶{]} & Lorenz سه‌بعدی & بله & معادلات لورنز & تمرکز بر تصاویر رنگی & 7.9974 & 99.65 & 30.35 \\
خدادادی {[}۷{]} & Chen + Logistic & بله & ترکیب کانال‌ها & دو سیستم آشوبی متوالی & 7.9971 & 99.58 & 33.21 \\
وانگ و لی {[}۸{]} & Logistic & بله & XOR & سادگی پیاده‌سازی & 7.9968 & 99.57 & 33.38 \\
مقاله {[}۹{]} & نگاشت‌های غیرخطی & بله & XOR دینامیکی & پردازش مستقل کانال‌ها & --- & --- & --- \\
لی {[}۱۰{]} & LEL نمایی & --- & --- & توالی پرپیچیدگی بالا & --- & --- & --- \\
صادقیان {[}۱۱{]} & آشوب + GA + PSO & --- & --- & بهینه‌سازی پارامتر & --- & --- & --- \\
فوق‌آشوب {[}۱۲{]} & Hyper-chaos & بله & زنجیره‌ای & سرعت بالا (۰.۰۲۳ s) & 7.9990 & 99.61 & 33.46 \\
مقاله {[}۱۳{]} & Lorenz+Henon+2D-Log & بله & جداگانه per کانال & ترکیب سه سیستم & --- & --- & --- \\
\textbf{روش پیشنهادی} & \textbf{Chen + ۹ ECM} & \textbf{Chen (RK45)} & \textbf{XOR پویا} & \textbf{۹ ECM per پیکسل} & \textbf{8.0} & \textbf{95.13} & \textbf{30.14} \\
\end{longtable}
}

توضیح: مقادیر آنتروپی، NPCR و UACI از گزارش مقالات مرجع یا از ارزیابی مشابه (جدول ۴-۷ برای روش پیشنهادی) استخراج شده‌اند. «---» به معنای عدم گزارش صریح آن معیار در منبع است.

\RL{\hfill\break
}

\begin{enumerate}
\def\labelenumi{\arabic{enumi}.}
\setcounter{enumi}{2}
\item
  \RL{\protect\phantomsection\label{_Toc232899896}{}معرفی الگوریتم پیشنهادی}
\end{enumerate}

\RL{\hfill\break
}

\begin{enumerate}
\def\labelenumi{\arabic{enumi}.}
\item
  \protect\phantomsection\label{_Toc232899897}{}\textbf{\RL{مقدمه}}
\end{enumerate}

\RL{این} \RL{فصل} \RL{به} \RL{تشریح} \RL{دقیق} \RL{الگوریتم} \RL{پیشنهادی} \RL{رمزنگاری} \RL{تصاویر} \RL{رنگی} \RL{اختصاص} \RL{دارد. هدف} \RL{اصلی،} \RL{ارائه} \RL{توضیح} \RL{شفاف} \RL{و} \RL{کامل} \RL{از} \RL{هر} \RL{یک} \RL{از} \RL{مراحل} \RL{فرآیند} \RL{رمزنگاری} \RL{است} \RL{به}‌\RL{گونه}‌\RL{ای} \RL{که} \RL{قابل} \RL{بازتولید} \RL{و} \RL{ارزیابی} \RL{مستقل} \RL{باشد. نوآوری} \RL{اصلی} \RL{این} \RL{پژوهش} \RL{در} \RL{به}‌\RL{کارگیری} \RL{ترکیبی} \RL{پویا} \RL{و} \RL{چندلایه} \RL{از} \RL{۹} \RL{سیستم} \RL{آشوبی} \RL{نمایی} \RL{در} \RL{کنار} \RL{سیستم} \RL{سه}‌\RL{بعدی} Chen \RL{برای} \RL{جایگشت} \RL{نهفته} \RL{است؛} \RL{رویکردی} \RL{که} \RL{موجب} \RL{ارتقاء} \RL{امنیت} \RL{در} \RL{برابر} \RL{انواع} \RL{حملات} \RL{تحلیل} \RL{آماری،} \RL{دیفرانسیلی} \RL{و} \RL{جستجوی} \RL{فراگیر} \RL{می}‌\RL{شود.}

\RL{ساختار} \RL{این} \RL{فصل} \RL{بر} \RL{سه} \RL{محور} \RL{اصلی} \RL{استوار} \RL{است: نخست،} \RL{تعریف} \RL{و} \RL{توجیه} \RL{انتخاب} \RL{مؤلفه}‌\RL{های} \RL{الگوریتم؛} \RL{دوم،} \RL{تشریح} \RL{گام}‌\RL{به}‌\RL{گام} \RL{فرآیند} \RL{رمزنگاری} \RL{و} \RL{رمزگشایی} \RL{با} \RL{ارائه} \RL{روابط} \RL{ریاضی} \RL{دقیق؛} \RL{و} \RL{سوم،} \RL{معرفی} \RL{روش} \RL{پیاده}‌\RL{سازی} \RL{و} \RL{ابزارهای} \RL{ارزیابی} \RL{که} \RL{در} \RL{فصل} \RL{چهارم} \RL{مورد} \RL{استفاده} \RL{قرار} \RL{خواهند} \RL{گرفت.}

\begin{enumerate}
\def\labelenumi{\arabic{enumi}.}
\setcounter{enumi}{1}
\item
  \protect\phantomsection\label{_Toc232899898}{}\RL{معرفی} \RL{سیستم}‌\RL{های} \RL{آشوبی} \RL{مورد} \RL{استفاده}

  \begin{enumerate}
  \def\labelenumii{\arabic{enumii}.}
  \item
    \protect\phantomsection\label{_Toc232899899}{}\RL{سیستم} \RL{آشوبی} \RL{سه}‌\RL{بعدی} Chen
  \end{enumerate}
\end{enumerate}

\RL{سیستم} \RL{آشوبی} Chen \RL{{[}6{]} یک} \RL{سیستم} \RL{دینامیکی} \RL{پیوسته} \RL{سه}‌\RL{بعدی} \RL{از} \RL{خانواده} \RL{سیستم}‌\RL{های} \RL{ابرجاذبی} \RL{است} \RL{که} \RL{در} \RL{سال} \RL{۱۹۹۹} \RL{معرفی} \RL{شد. این} \RL{سیستم} \RL{از} \RL{سه} \RL{معادله} \RL{دیفرانسیل} \RL{غیرخطی} \RL{کوپل}‌\RL{شده} \RL{تشکیل} \RL{شده} \RL{و} \RL{جاذب} \RL{آشوبی} \RL{آن} \RL{طیف} \RL{لیاپانوف} \RL{مثبت} \RL{گسترده}‌\RL{ای} \RL{دارد. ویژگی}‌\RL{هایی} \RL{نظیر} \RL{حساسیت} \RL{فوق}‌\RL{العاده} \RL{بالا} \RL{به} \RL{شرایط} \RL{اولیه،} \RL{پویایی} \RL{سه}‌\RL{بعدی} \RL{پیچیده} \RL{و} \RL{تولید} \RL{دنباله}‌\RL{های} \RL{شبه}‌\RL{تصادفی} \RL{با} \RL{توزیع} \RL{گسترده،} \RL{این} \RL{سیستم} \RL{را} \RL{به} \RL{گزینه}‌\RL{ای} \RL{ایده}‌\RL{آل} \RL{برای} \RL{مرحله} \RL{جایگشت} \RL{تبدیل} \RL{کرده}‌\RL{اند.}

\RL{معادلات} \RL{حاکم} \RL{بر} \RL{سیستم} Chen \RL{به}‌\RL{صورت} \RL{زیر} \RL{تعریف} \RL{می}‌\RL{شوند:}

{\def\LTcaptype{none} % do not increment counter
\begin{longtable}[]{@{}
  >{\raggedright\arraybackslash}p{(\linewidth - 2\tabcolsep) * \real{0.2148}}
  >{\raggedright\arraybackslash}p{(\linewidth - 2\tabcolsep) * \real{0.7852}}@{}}
\toprule\noalign{}
\begin{minipage}[b]{\linewidth}\raggedright
\begin{enumerate}
\def\labelenumi{\arabic{enumi}.}
\item
  \RL{معادلات} \RL{سیستم} \RL{آشوبی} Chen
\end{enumerate}
\end{minipage} & \begin{minipage}[b]{\linewidth}\raggedright
\begin{quote}
\[CHEN:\ \ \ \ \left\{ \begin{array}{r}
\frac{dx}{dt} = a(y - x) \\
\frac{dy}{dt} = (c - a)x - xz + cy \\
\frac{dz}{dt} = xy - bx
\end{array} \right.\ \]
\end{quote}
\end{minipage} \\
\midrule\noalign{}
\endhead
\bottomrule\noalign{}
\endlastfoot
\end{longtable}
}

\RL{در} \RL{این} \RL{معادلات،} x(t)\RL{،} y(t) \RL{و} z(t) \RL{متغیرهای} \RL{حالت} \RL{سیستم} \RL{در} \RL{زمان} t \RL{هستند. پارامترهای} a\RL{،} b \RL{و} c \RL{ثابت}‌\RL{های} \RL{سیستمی} \RL{هستند} \RL{که} \RL{رفتار} \RL{آشوبی} \RL{را} \RL{تعریف} \RL{می}‌\RL{کنند. مقادیر} \RL{استاندارد} \RL{اثبات}‌\RL{شده} \RL{برای} \RL{ایجاد} \RL{آشوب} \RL{پایدار} \RL{در} \RL{این} \RL{سیستم} \RL{عبارتند} \RL{از:} a = 35, b = 3, c = 28

\RL{این} \RL{مقادیر} \RL{در مقاله} {[}6{]} \RL{توسط} Chen \RL{و} Ueta\RL{تأیید} \RL{شده}‌\RL{اند} \RL{و} \RL{در} \RL{بازه} \RL{وسیعی} \RL{از} \RL{مطالعات} \RL{رمزنگاری} \RL{آشوبی} \RL{به}‌\RL{کار} \RL{رفته}‌\RL{اند. با} \RL{این} \RL{پارامترها،} \RL{نمای} \RL{لیاپانوف} \RL{سیستم} \RL{مثبت} \RL{است} (λ₁ ≈ 2. 03) \RL{که} \RL{وجود} \RL{آشوب} \RL{را} \RL{تضمین} \RL{می}‌\RL{کند. شرایط} \RL{اولیه} \RL{سیستم} Chen \RL{یعنی} \RL{مقادیر} x₀\RL{،} y₀ \RL{و} z₀ \RL{بخش} \RL{اصلی} \RL{کلید} \RL{رمزنگاری} \RL{را} \RL{تشکیل} \RL{می}‌\RL{دهند} \RL{و} \RL{باید} \RL{به}‌\RL{صورت} \RL{دقیق} \RL{با} \RL{دقت} \RL{اعشاری} 10\textsuperscript{-15} \RL{حفظ} \RL{شوند. هرگونه} \RL{انحراف} \RL{جزئی} \RL{از} \RL{این} \RL{مقادیر،} \RL{دنباله} \RL{جایگشتی} \RL{کاملاً} \RL{متفاوتی} \RL{تولید} \RL{می}‌\RL{کند} \RL{و} \RL{رمزگشایی} \RL{را} \RL{ناممکن} \RL{می}‌\RL{سازد.}

\begin{enumerate}
\def\labelenumi{\arabic{enumi}.}
\setcounter{enumi}{1}
\item
  \protect\phantomsection\label{_Toc232899900}{}\RL{نقشه}‌\RL{های} \RL{آشوبی} \RL{نمایی} (ECM)
\end{enumerate}

\RL{سیستم}‌\RL{های} \RL{آشوبی} \RL{نمایی} \RL{توسط} \RL{هو} \RL{و} \RL{ژو} {[}4{]}\RL{معرفی} \RL{شده}‌\RL{اند. ایده} \RL{اصلی} \RL{این} \RL{رویکرد،} \RL{ترکیب} \RL{دو} \RL{نقشه} \RL{آشوبی} \RL{یک}‌\RL{بعدی} \RL{از} \RL{طریق} \RL{یک} \RL{تابع} \RL{نمایی} \RL{است} \RL{تا} \RL{نقشه} \RL{جدیدی} \RL{با} \RL{پایداری} \RL{آشوب} \RL{بسیار} \RL{بالاتر} \RL{ایجاد} \RL{شود. این} \RL{رویکرد} \RL{به}‌\RL{طور} \RL{مؤثری} \RL{پدیده} \RL{«تخریب} \RL{آشوب»} \RL{را} \RL{که} \RL{در} \RL{نقشه}‌\RL{های} \RL{کلاسیک} \RL{مانند} Logistic \RL{و} Tent \RL{مشاهده} \RL{می}‌\RL{شود،} \RL{برطرف} \RL{می}‌\RL{کند.}

\RL{ساختار} \RL{کلی} \RL{یک} \RL{نقشه} \RL{آشوبی} \RL{نمایی} \RL{به}‌\RL{صورت} \RL{زیر} \RL{تعریف} \RL{می}‌\RL{شود:}

{\def\LTcaptype{none} % do not increment counter
\begin{longtable}[]{@{}
  >{\raggedright\arraybackslash}p{(\linewidth - 2\tabcolsep) * \real{0.4694}}
  >{\raggedright\arraybackslash}p{(\linewidth - 2\tabcolsep) * \real{0.5167}}@{}}
\toprule\noalign{}
\begin{minipage}[b]{\linewidth}\raggedright
\begin{enumerate}
\def\labelenumi{\arabic{enumi}.}
\item
  \RL{ساختار} \RL{کلی} \RL{نقشه} \RL{آشوبی} \RL{نمایی}
\end{enumerate}
\end{minipage} & \begin{minipage}[b]{\linewidth}\centering
\[f_{ECM(x)} = f_{outer}\left( \ exp\ \left( \mu\  \times \ f_{inner(x)} \right)\  \right)\ \ mod\ 1\]
\end{minipage} \\
\midrule\noalign{}
\endhead
\bottomrule\noalign{}
\endlastfoot
\end{longtable}
}

\RL{در} \RL{این} \RL{رابطه،} f\textsubscript{inner} \RL{نقشه} \RL{آشوبی} \RL{داخلی،} f\textsubscript{outer} \RL{نقشه} \RL{آشوبی} \RL{خارجی،} μ \RL{پارامتر} \RL{نمایی} \RL{و} x \RL{متغیر} \RL{حالت} \RL{است. ترکیب} \RL{سه} \RL{عنصر} Logistic (L)\RL{،} Sine (S) \RL{و} Tent (T) \RL{در} \RL{نقش} \RL{داخلی} \RL{و} \RL{خارجی،} \RL{۹} \RL{نقشه} \RL{متمایز} \RL{ایجاد} \RL{می}‌\RL{کند:}

\begin{enumerate}
\def\labelenumi{\arabic{enumi}.}
\item
  \protect\phantomsection\label{_Toc232899945}{}\RL{۹} \RL{سیستم} \RL{آشوبی} \RL{نمایی} \RL{مورد} \RL{استفاده} \RL{در} \RL{الگوریتم} \RL{پیشنهادی}
\end{enumerate}

{\def\LTcaptype{none} % do not increment counter
\begin{longtable}[]{@{}
  >{\raggedright\arraybackslash}p{(\linewidth - 6\tabcolsep) * \real{0.3434}}
  >{\raggedright\arraybackslash}p{(\linewidth - 6\tabcolsep) * \real{0.2458}}
  >{\raggedright\arraybackslash}p{(\linewidth - 6\tabcolsep) * \real{0.2231}}
  >{\raggedright\arraybackslash}p{(\linewidth - 6\tabcolsep) * \real{0.1876}}@{}}
\toprule\noalign{}
\begin{minipage}[b]{\linewidth}\centering
\textbf{\RL{بازه آشوب پایدار}}
\end{minipage} & \begin{minipage}[b]{\linewidth}\centering
\textbf{\RL{نقشه خارجی}}
\end{minipage} & \begin{minipage}[b]{\linewidth}\centering
\textbf{\RL{نقشه داخلی}}
\end{minipage} & \begin{minipage}[b]{\linewidth}\centering
\textbf{\RL{نام سیستم}}
\end{minipage} \\
\midrule\noalign{}
\endhead
\bottomrule\noalign{}
\endlastfoot
μ ∈ (3. 57, 4. 0) & Logistic & Logistic & LEL \\
μ ∈ (3. 5, 4. 0) & Sine & Logistic & LES \\
μ ∈ (3. 6, 4. 0) & Tent & Logistic & LET \\
μ ∈ (0. 8π, π) & Logistic & Sine & SEL \\
μ ∈ (0. 85π, π) & Sine & Sine & SES \\
μ ∈ (0. 8π, π) & Tent & Sine & SET \\
μ ∈ (0. 5, 1. 0) & Logistic & Tent & TEL \\
μ ∈ (0. 5, 1. 0) & Sine & Tent & TES \\
μ ∈ (0. 5, 1. 0) & Tent & Tent & TET \\
\end{longtable}
}

\RL{ویژگی} \RL{برجسته} \RL{این} \RL{۹} \RL{سیستم} \RL{آن} \RL{است} \RL{که} \RL{هریک} \RL{از} \RL{آن}‌\RL{ها} \RL{دامنه} \RL{آشوب} \RL{پایدار} \RL{گسترده}‌\RL{ای} \RL{دارند} \RL{و} \RL{خروجی} \RL{آن}‌\RL{ها} \RL{در} \RL{بازه{[}0،1{]} قرار} \RL{می}‌\RL{گیرد. این} \RL{خاصیت} \RL{انتخاب} \RL{تصادفی} \RL{سیستم} \RL{رمزکننده} \RL{برای} \RL{هر} \RL{پیکسل} \RL{را} \RL{ممکن} \RL{می}‌\RL{سازد} \RL{بدون} \RL{آنکه} \RL{نیاز} \RL{به} \RL{بررسی} \RL{پایداری} \RL{در} \RL{هر} \RL{مرحله} \RL{باشد.}

\RL{روابط} \RL{دقیق} \RL{هر} \RL{یک} \RL{از} \RL{نقشه}‌\RL{های} \RL{اصلی} \RL{تشکیل}‌\RL{دهنده} ECM \RL{به}‌\RL{صورت} \RL{زیر} \RL{تعریف} \RL{می}‌\RL{شوند:}

\begin{itemize}
\item
  \RL{نقشه لجستیک (}Logistic\RL{):} \(L(x)\  = \ r\  \times \ x\  \times \ (1\  - \ x)\) \RL{که در آن} r = 4 \RL{برای حالت آشوب کامل در نظر گرفته می‌شود.}
\item
  \RL{نقشه سینوسی (}\textbf{Sine}\RL{):} \(S(x)\  = \ (r/4)\  \times \ sin(\pi x)\) \RL{که در آن} r ∈ (0, 1{]} \RL{است.}
\item
  \RL{نقشه تِنت (}\textbf{Tent}\RL{):} T(x) = 2x \RL{اگر} x \textless{} 0. 5 \RL{و} 2(1-x) \RL{اگر} x ≥ 0. 5

  \begin{enumerate}
  \def\labelenumi{\arabic{enumi}.}
  \item
    \protect\phantomsection\label{_Toc232899901}{}\RL{تعریف} \RL{فضای} \RL{کلید}
  \end{enumerate}
\end{itemize}

\RL{امنیت} \RL{هر} \RL{الگوریتم} \RL{رمزنگاری} \RL{به}‌\RL{طور} \RL{مستقیم} \RL{به} \RL{اندازه} \RL{فضای} \RL{کلید} \RL{وابسته} \RL{است. در} \RL{الگوریتم} \RL{پیشنهادی،} \RL{کلید} \RL{رمزنگاری} \RL{از} \RL{مجموعه}‌\RL{ای} \RL{از} \RL{پارامترهای} \RL{پیوسته} \RL{و} \RL{گسسته} \RL{تشکیل} \RL{شده} \RL{است} \RL{که} \RL{در} \RL{جدول} \RL{۳}-\RL{۲} \RL{آمده} \RL{است:}

\begin{enumerate}
\def\labelenumi{\arabic{enumi}.}
\item
  \protect\phantomsection\label{_Toc232899946}{}\RL{مؤلفه}‌\RL{های} \RL{فضای} \RL{کلید} \RL{الگوریتم} \RL{پیشنهادی}
\end{enumerate}

{\def\LTcaptype{none} % do not increment counter
\begin{longtable}[]{@{}
  >{\raggedright\arraybackslash}p{(\linewidth - 6\tabcolsep) * \real{0.2631}}
  >{\raggedright\arraybackslash}p{(\linewidth - 6\tabcolsep) * \real{0.1051}}
  >{\raggedright\arraybackslash}p{(\linewidth - 6\tabcolsep) * \real{0.1117}}
  >{\raggedright\arraybackslash}p{(\linewidth - 6\tabcolsep) * \real{0.5200}}@{}}
\toprule\noalign{}
\begin{minipage}[b]{\linewidth}\centering
\textbf{\RL{بازه}}
\end{minipage} & \begin{minipage}[b]{\linewidth}\centering
\textbf{\RL{دقت}}
\end{minipage} & \begin{minipage}[b]{\linewidth}\centering
\textbf{\RL{نوع}}
\end{minipage} & \begin{minipage}[b]{\linewidth}\centering
\textbf{\RL{مؤلفه کلید}}
\end{minipage} \\
\midrule\noalign{}
\endhead
\bottomrule\noalign{}
\endlastfoot
(0, 1) & 10⁻¹⁵ & \RL{پیوسته} & x₀ \RL{(شرط اولیه} Chen\RL{)} \\
(0, 1) & 10⁻¹⁵ & \RL{پیوسته} & \RL{}y₀ \RL{(شرط اولیه} Chen\RL{)} \\
(0, 30) & 10⁻¹⁵ & \RL{پیوسته} & \RL{}z₀ \RL{(شرط اولیه} Chen\RL{)} \\
(0, 1) & 10⁻¹⁵ & \RL{پیوسته} & x\textsubscript{sel} \RL{(شرط اولیه سلکتور)} \\
(0, 1) & 10⁻¹⁵ & \RL{پیوسته} & x\textsubscript{ecm} \RL{(شرط اولیه} ECM\RL{)} \\
\RL{وابسته به سیستم} & 10⁻¹⁵ & \RL{پیوسته} & μ \RL{(پارامتر نمایی} ECM\RL{)} \\
(0, 1) & 10⁻¹⁵ & \RL{پیوسته} & x\textsubscript{final \RL{}}\RL{(شرط اولیه مرحله ۵)} \\
\end{longtable}
}

\RL{با} \RL{فرض} \RL{دقت} \RL{۱۵} \RL{رقم} \RL{اعشار} \RL{برای} \RL{هر} \RL{پارامتر} \RL{پیوسته،} \RL{اندازه} \RL{فضای} \RL{کلید} \RL{به}‌\RL{صورت} \RL{تقریبی} \RL{از} \RL{رابطه} \RL{۳}-\RL{۳} \RL{محاسبه} \RL{می}‌\RL{شود:}

{\def\LTcaptype{none} % do not increment counter
\begin{longtable}[]{@{}
  >{\raggedright\arraybackslash}p{(\linewidth - 2\tabcolsep) * \real{0.5016}}
  >{\raggedright\arraybackslash}p{(\linewidth - 2\tabcolsep) * \real{0.4845}}@{}}
\toprule\noalign{}
\begin{minipage}[b]{\linewidth}\raggedright
\begin{enumerate}
\def\labelenumi{\arabic{enumi}.}
\item
  \RL{تخمین} \RL{اندازه} \RL{فضای} \RL{کلید}
\end{enumerate}
\end{minipage} & \begin{minipage}[b]{\linewidth}\raggedleft
\[|K| \approx 10^{15 \times 7}\  = \ 10^{105} > \ 2^{348}\]
\end{minipage} \\
\midrule\noalign{}
\endhead
\bottomrule\noalign{}
\endlastfoot
\end{longtable}
}

\RL{این} \RL{مقدار} \RL{به} \RL{مراتب} \RL{فراتر} \RL{از} \RL{آستانه} \RL{امنیتی} \RL{۱۲۸} \RL{بیتی} \RL{استاندارد} AES \RL{است} \RL{و} \RL{حمله} \RL{جستجوی} \RL{فراگیر} \RL{را} \RL{به} \RL{لحاظ} \RL{محاسباتی} \RL{ناممکن} \RL{می}‌\RL{سازد.}

\begin{enumerate}
\def\labelenumi{\arabic{enumi}.}
\setcounter{enumi}{1}
\item
  \protect\phantomsection\label{_Toc232899902}{}\RL{الگوریتم} \RL{پیشنهادی: توضیح} \RL{گام}‌\RL{به}‌\RL{گام}
\end{enumerate}

\RL{الگوریتم} \RL{رمزنگاری} \RL{پیشنهادی} \RL{از} \RL{پنج} \RL{مرحله} \RL{اصلی} \RL{تشکیل} \RL{شده} \RL{است. در} \RL{ادامه،} \RL{هر} \RL{مرحله} \RL{با} \RL{جزئیات} \RL{کامل} \RL{ریاضی} \RL{تشریح} \RL{می}‌\RL{شود.}

\subsubsection{\texorpdfstring{\RL{مرحله} \RL{۱: تولید} \RL{کلیدهای} \RL{اولیه} \RL{و} \RL{آماده}‌\RL{سازی}}{مرحله ۱: تولید کلیدهای اولیه و آماده‌سازی}}\label{ux645ux631ux62dux644ux647-ux6f1-ux62aux648ux644ux6ccux62f-ux6a9ux644ux6ccux62fux647ux627ux6cc-ux627ux648ux644ux6ccux647-ux648-ux622ux645ux627ux62fux647ux633ux627ux632ux6cc}

\RL{در} \RL{این} \RL{مرحله،} \RL{تمام} \RL{پارامترهای} \RL{مورد} \RL{نیاز} \RL{الگوریتم} \RL{تولید} \RL{و} \RL{مقداردهی} \RL{اولیه} \RL{می}‌\RL{شوند. کلید} \RL{رمزنگاری} K \RL{شامل} \RL{هفت} \RL{عدد} \RL{مثبت} \RL{پیوسته} \RL{با} \RL{دقت} \RL{۱۵} \RL{رقم} \RL{اعشار} \RL{است:}

\[K\  = \ \left\{ x_{0},\ y_{0},\ z_{0},\ x_{sel},\ x_{ecm},\ \mu\ ,\ x_{final} \right\}\]

\RL{تصویر} \RL{رنگی} \RL{ورودی} \RL{با} \RL{ابعاد} H × W × 3 \RL{بارگذاری} \RL{شده} \RL{و} \RL{کانال}‌\RL{های} R\RL{،} G \RL{و} B \RL{آن} \RL{تفکیک} \RL{می}‌\RL{شوند. هر} \RL{کانال} \RL{به}‌\RL{صورت} \RL{مستقل} \RL{پردازش} \RL{می}‌\RL{شود. مجموع} \RL{تعداد} \RL{پیکسل}‌\RL{های} \RL{هر} \RL{کانال} N = H × W \RL{است.}

\RL{به}‌\RL{منظور} \RL{حذف} \RL{ترانزیت} \RL{اولیه} \RL{سیستم} \RL{آشوبی} \RL{و} \RL{اطمینان} \RL{از} \RL{رسیدن} \RL{به} \RL{ناحیه} \RL{آشوبی} \RL{پایدار،} \RL{سیستم} Chen \RL{برای} \RL{۱۰۰۰} \RL{گام} \RL{اولیه} \RL{اجرا} \RL{شده} \RL{و} \RL{این} \RL{مقادیر} \RL{دور} \RL{انداخته} \RL{می}‌\RL{شوند.} \RL{این} \RL{کار} \RL{یک} \RL{الزام} \RL{علمی} \RL{برای} \RL{جلوگیری} \RL{از} \RL{ارتباط} \RL{احتمالی} \RL{بین} \RL{کلید} \RL{و} \RL{خروجی} \RL{اولیه} \RL{است.}

\subsubsection{\texorpdfstring{\RL{مرحله} \RL{۲: جایگشت} \RL{پیکسل}‌\RL{ها} \RL{با} \RL{سیستم} Chen}{مرحله ۲: جایگشت پیکسل‌ها با سیستم Chen}}\label{ux645ux631ux62dux644ux647-ux6f2-ux62cux627ux6ccux6afux634ux62a-ux67eux6ccux6a9ux633ux644ux647ux627-ux628ux627-ux633ux6ccux633ux62aux645-chen}

\RL{هدف} \RL{این} \RL{مرحله،} \RL{از} \RL{بین} \RL{بردن} \RL{همبستگی} \RL{مکانی} \RL{بین} \RL{پیکسل}‌\RL{های} \RL{مجاور} \RL{از} \RL{طریق} \RL{جابه}‌\RL{جایی} \RL{تصادفی} \RL{موقعیت} \RL{آن}‌\RL{هاست. سیستم} Chen \RL{با} \RL{شرایط} \RL{اولیه} (x₀, y₀, z₀) \RL{و} \RL{پارامترهای} a=35, b=3, c=28 \RL{حل} \RL{می}‌\RL{شود.}

\RL{معادلات} \RL{دیفرانسیل} \RL{سیستم} Chen \RL{با} \RL{استفاده} \RL{از} \RL{حل}‌\RL{کننده} \RL{عددی} RK45 \RL{(روش} \RL{رانگ}-\RL{کوتای} \RL{مرتبه} \RL{۴} \RL{و} \RL{۵)} \RL{با} \RL{گام} \RL{زمانی} dt = 0. 002 \RL{حل} \RL{می}‌\RL{شوند.} \RL{برای} \RL{هر} \RL{تصویر} \RL{با} N \RL{پیکسل،} N + 1000 \RL{نقطه} \RL{زمانی} \RL{تولید} \RL{می}‌\RL{شود} \RL{که} \RL{۱۰۰۰} \RL{نقطه} \RL{اول} \RL{(دوره} warm-up\RL{)} \RL{دور} \RL{انداخته} \RL{می}‌\RL{شوند{[}6{]}.}

\RL{از} \RL{خروجی} x(t) \RL{سیستم} Chen\RL{،} \RL{یک} \RL{بردار} N\RL{تایی} \RL{استخراج} \RL{شده} \RL{و} \RL{برای} \RL{تولید} \RL{جایگشت} \RL{یک}‌\RL{به}‌\RL{یک} \RL{به} \RL{روش} \RL{زیر} \RL{استفاده} \RL{می}‌\RL{شود:}

{\def\LTcaptype{none} % do not increment counter
\begin{longtable}[]{@{}
  >{\raggedright\arraybackslash}p{(\linewidth - 2\tabcolsep) * \real{0.3402}}
  >{\raggedright\arraybackslash}p{(\linewidth - 2\tabcolsep) * \real{0.6459}}@{}}
\toprule\noalign{}
\begin{minipage}[b]{\linewidth}\raggedright
\begin{enumerate}
\def\labelenumi{\arabic{enumi}.}
\item
  \RL{تولید} \RL{بردار} \RL{جایگشت} \RL{از} \RL{دنباله} \RL{آشوبی} Chen
\end{enumerate}
\end{minipage} & \begin{minipage}[b]{\linewidth}\raggedright
\[perm\  = \ argsort\left( x_{chen\lbrack 1000:\ 1000 + N\rbrack} \right)\]

\[I_{perm\lbrack c\rbrack} = \ I_{original}\lbrack c\rbrack.\ flatten()\lbrack perm\rbrack.\ reshape\ (H,\ W)\]
\end{minipage} \\
\midrule\noalign{}
\endhead
\bottomrule\noalign{}
\endlastfoot
\end{longtable}
}

\RL{عملیات} \RL{مرتب}‌\RL{سازی} \RL{غیرمستقیم}\footnote{argsort} \RL{یک} \RL{بردار} \RL{پیوسته} \RL{را} \RL{به} \RL{یک} \RL{جایگشت} \RL{گسسته} \RL{یک}‌\RL{به}‌\RL{یک} \RL{تبدیل} \RL{می}‌\RL{کند. از} \RL{آنجا} \RL{که} \RL{دنباله} Chen \RL{فاقد} \RL{هرگونه} \RL{الگوی} \RL{دوره}‌\RL{ای} \RL{یا} \RL{خطی} \RL{است،} \RL{بردار} \RL{جایگشت} \RL{تولیدشده} \RL{کاملاً} \RL{غیرقابل} \RL{پیش}‌\RL{بینی} \RL{خواهد} \RL{بود.} \RL{این} \RL{بردار} \RL{برای} \RL{رمزگشایی} \RL{(جایگشت} \RL{معکوس)} \RL{ذخیره} \RL{نمی}‌\RL{شود} \RL{بلکه} \RL{از} \RL{کلید} \RL{بازسازی} \RL{می}‌\RL{شود.}

\subsubsection{\texorpdfstring{\RL{مرحله} \RL{۳: انتخاب} \RL{پویای} \RL{سیستم} \RL{آشوبی} \RL{نمایی} \RL{برای} \RL{هر} \RL{پیکسل}}{مرحله ۳: انتخاب پویای سیستم آشوبی نمایی برای هر پیکسل}}\label{ux645ux631ux62dux644ux647-ux6f3-ux627ux646ux62aux62eux627ux628-ux67eux648ux6ccux627ux6cc-ux633ux6ccux633ux62aux645-ux622ux634ux648ux628ux6cc-ux646ux645ux627ux6ccux6cc-ux628ux631ux627ux6cc-ux647ux631-ux67eux6ccux6a9ux633ux644}

\RL{این} \RL{مرحله،} \RL{قلب} \RL{نوآوری} \RL{الگوریتم} \RL{پیشنهادی} \RL{است. به}‌\RL{جای} \RL{استفاده} \RL{از} \RL{یک} \RL{سیستم} \RL{آشوبی} \RL{یکسان} \RL{برای} \RL{تمام} \RL{پیکسل}‌\RL{ها،} \RL{برای} \RL{هر} \RL{پیکسل} \RL{به}‌\RL{صورت} \RL{مستقل} \RL{یکی} \RL{از} \RL{۹} \RL{سیستم} \RL{نمایی} \RL{انتخاب} \RL{می}‌\RL{شود. این} \RL{انتخاب} \RL{از} \RL{طریق} \RL{یک} \RL{نقشه} Logistic \RL{مستقل} \RL{انجام} \RL{می}‌\RL{شود} \RL{که} \RL{با} \RL{شرط} \RL{اولیه} x\textsubscript{sel} \RL{مقداردهی} \RL{می}‌\RL{شود:}

\RL{پس} \RL{از} warm-up \RL{اولیه} \RL{۱۰۰۰} \RL{گام،} \RL{مقادیر} \RL{پیوسته} Logistic \RL{در} \RL{بازه} (0, 1{]} \RL{به} \RL{اعداد} \RL{صحیح} \RL{۱} \RL{تا} \RL{۹} \RL{از} \RL{طریق} \RL{رابطه} \RL{۳}-\RL{۵} \RL{تبدیل} \RL{می}‌\RL{شوند:}

{\def\LTcaptype{none} % do not increment counter
\begin{longtable}[]{@{}
  >{\raggedright\arraybackslash}p{(\linewidth - 2\tabcolsep) * \real{0.3240}}
  >{\raggedright\arraybackslash}p{(\linewidth - 2\tabcolsep) * \real{0.6621}}@{}}
\toprule\noalign{}
\begin{minipage}[b]{\linewidth}\raggedright
\begin{enumerate}
\def\labelenumi{\arabic{enumi}.}
\setcounter{enumi}{1}
\item
  \RL{تبدیل} \RL{خروجی} Logistic \RL{به} \RL{شاخص} \RL{۱} \RL{تا} \RL{۹} \RL{برای} \RL{انتخاب} \RL{سیستم} \RL{نمایی}
\end{enumerate}
\end{minipage} & \begin{minipage}[b]{\linewidth}\raggedleft
\[sel(k) = \ (floor\left( x_{\left\{ sel_{k} \right\}\  \times \ 10¹⁴)\ mod\ 9\  + \ 1}\  \right),\ k\  = \ 1,\ 2,\ \ldots,\ N\]

\[x_{\left\{ sel_{\left\{ k + 1 \right\}} \right\}}\  = \ 4 \times \ x_{\{ sel_{k}\}} \times \ \left( 1\  - \ x_{\left\{ sel_{k} \right\}} \right)\ (Logistic,\ r = 4)\]
\end{minipage} \\
\midrule\noalign{}
\endhead
\bottomrule\noalign{}
\endlastfoot
\end{longtable}
}

\RL{استفاده} \RL{از} r = 4 \RL{در} \RL{نقشه} Logistic \RL{تضمین} \RL{می}‌\RL{کند} \RL{که} \RL{سیستم} \RL{در} \RL{حالت} \RL{آشوب} \RL{کامل} \RL{قرار} \RL{دارد} \RL{و} \RL{توزیع} \RL{یکنواختی} \RL{از} \RL{اعداد} \RL{۱} \RL{تا} \RL{۹} \RL{تولید} \RL{می}‌\RL{شود. مزیت} \RL{کلیدی} \RL{این} \RL{رویکرد} \RL{آن} \RL{است} \RL{که} \RL{اگر} \RL{مهاجم} \RL{حتی} \RL{موفق} \RL{به} \RL{شناسایی} \RL{دنباله} \RL{رمزنگاری} \RL{یک} \RL{پیکسل} \RL{شود،} \RL{اطلاعات} \RL{حاصل} \RL{هیچ} \RL{کمکی} \RL{به} \RL{پیش}‌\RL{بینی} \RL{سیستم} \RL{استفاده}‌\RL{شده} \RL{برای} \RL{پیکسل} \RL{بعدی} \RL{نمی}‌\RL{کند{[}5{]}.}

\RL{برای} \RL{هر} \RL{پیکسل} k \RL{ام} \RL{با} \RL{شاخص} sel(k)\RL{،} \RL{سیستم} \RL{آشوبی} \RL{نمایی} \RL{متناظر} \RL{(طبق} \RL{جدول} \RL{۳}-\RL{۱)} \RL{به} \RL{اندازه} \RL{۱} \RL{گام} \RL{اجرا} \RL{می}‌\RL{شود} \RL{و} \RL{خروجی} \RL{آن} X\textsubscript{ecmk} \RL{نامیده} \RL{می}‌\RL{شود. این} \RL{مقدار} \RL{در} \RL{مرحله} \RL{۴} \RL{به} \RL{عنوان} \RL{کلید} \RL{رمزنگاری} \RL{پیکسل} k \RL{به}‌\RL{کار} \RL{می}‌\RL{رود.}

\subsubsection{\texorpdfstring{\RL{مرحله} \RL{۴: انتشار} \RL{مقادیر} \RL{رنگی} \RL{با} XOR}{مرحله ۴: انتشار مقادیر رنگی با XOR}}\label{ux645ux631ux62dux644ux647-ux6f4-ux627ux646ux62aux634ux627ux631-ux645ux642ux627ux62fux6ccux631-ux631ux646ux6afux6cc-ux628ux627-xor}

\RL{در} \RL{مرحله} \RL{انتشار،} \RL{مقادیر} \RL{رنگی} \RL{هر} \RL{پیکسل} \RL{با} \RL{یک} \RL{بایت} \RL{کلید} \RL{تولیدشده} \RL{از} \RL{سیستم} \RL{نمایی} \RL{انتخابی} \RL{ترکیب} \RL{می}‌\RL{شوند. برای} \RL{تبدیل} \RL{خروجی} \RL{پیوسته} \RL{سیستم} \RL{نمایی} (\RL{در} \RL{بازه} {[}\RL{۰،} \RL{۱})) \RL{به} \RL{یک} \RL{بایت} \RL{۸} \RL{بیتی} \RL{از} \RL{رابطه} \RL{۳}-\RL{۶} \RL{استفاده} \RL{می}‌\RL{شود:}

{\def\LTcaptype{none} % do not increment counter
\begin{longtable}[]{@{}
  >{\raggedright\arraybackslash}p{(\linewidth - 2\tabcolsep) * \real{0.2756}}
  >{\raggedright\arraybackslash}p{(\linewidth - 2\tabcolsep) * \real{0.7105}}@{}}
\toprule\noalign{}
\begin{minipage}[b]{\linewidth}\raggedright
\begin{enumerate}
\def\labelenumi{\arabic{enumi}.}
\setcounter{enumi}{2}
\item
  \RL{تبدیل} \RL{خروجی} \RL{سیستم} \RL{نمایی} \RL{به} \RL{بایت} \RL{۸} \RL{بیتی}
\end{enumerate}
\end{minipage} & \begin{minipage}[b]{\linewidth}\raggedleft
\[key_{byte}(k)\  = \ (floor\left( x_{\left\{ ecm_{k} \right\}} \times \ 10¹⁴) \right)\ \ \ mod\ 256\]
\end{minipage} \\
\midrule\noalign{}
\endhead
\bottomrule\noalign{}
\endlastfoot
\end{longtable}
}

\RL{عملیات} XOR \RL{بین} \RL{مقدار} \RL{رنگی} \RL{پیکسل} \RL{جایگشت}‌\RL{یافته} \RL{و} \RL{بایت} \RL{کلید} \RL{انجام} \RL{می}‌\RL{شود:}

{\def\LTcaptype{none} % do not increment counter
\begin{longtable}[]{@{}
  >{\raggedright\arraybackslash}p{(\linewidth - 2\tabcolsep) * \real{0.2756}}
  >{\raggedright\arraybackslash}p{(\linewidth - 2\tabcolsep) * \real{0.7105}}@{}}
\toprule\noalign{}
\begin{minipage}[b]{\linewidth}\raggedright
\begin{enumerate}
\def\labelenumi{\arabic{enumi}.}
\setcounter{enumi}{3}
\item
  \RL{عملیات} XOR \RL{برای رمزنگاری هر پیکسل}
\end{enumerate}
\end{minipage} & \begin{minipage}[b]{\linewidth}\raggedleft
\[C\_ k\  = \ P\_\{ perm\_ k\ \}\ \ \{ XOR\}\ \ key\_\{ byte\}(k)\ ,\ \ \forall\ k \in \ \{ 1,\ldots,N\}\]
\end{minipage} \\
\midrule\noalign{}
\endhead
\bottomrule\noalign{}
\endlastfoot
\end{longtable}
}

\RL{که} \RL{در} \RL{آن} Pperm\textsubscript{k} \RL{مقدار} \RL{پیکسل} \RL{جایگشت}‌\RL{یافته} \RL{و} C\textsubscript{k} \RL{مقدار} \RL{پیکسل} \RL{رمزنگاری}‌\RL{شده} \RL{است. عملیات} XOR \RL{به} \RL{دلیل} \RL{برگشت}‌\RL{پذیری} \RL{کامل} (C XOR key = P) \RL{در} \RL{رمزنگاری} \RL{ایده}‌\RL{آل} \RL{است} \RL{و} \RL{هیچ} \RL{اطلاعاتی} \RL{از} \RL{تصویر} \RL{اصلی} \RL{را} \RL{در} \RL{تصویر} \RL{رمزنگاری}‌\RL{شده} \RL{آشکار} \RL{نمی}‌\RL{کند.}

\RL{نکته} \RL{کلیدی} \RL{طراحی: دنباله} key\textsubscript{byte} \RL{صرفاً} \RL{به} \RL{پارامترهای} \RL{کلید} \RL{و} \RL{شاخص} \RL{پیکسل} \RL{وابسته} \RL{است} \RL{و} \RL{هیچ} \RL{وابستگی} \RL{به} \RL{مقادیر} \RL{پیکسل}‌\RL{های} \RL{تصویر} \RL{ندارد. این} \RL{طراحی} \RL{تضمین} \RL{می}‌\RL{کند} \RL{که} \RL{فرآیند} \RL{رمزگشایی} \RL{کاملاً} \RL{معکوس}‌\RL{پذیر} \RL{است} \RL{و} \RL{هیچ} \RL{خطای} \RL{انتشاری} \RL{رخ} \RL{نمی}‌\RL{دهد{[}6{]}.}

\subsubsection{\texorpdfstring{\RL{مرحله} \RL{۵: ترکیب} \RL{نهایی} \RL{لایه} XOR \RL{دوم}}{مرحله ۵: ترکیب نهایی لایه XOR دوم}}\label{ux645ux631ux62dux644ux647-ux6f5-ux62aux631ux6a9ux6ccux628-ux646ux647ux627ux6ccux6cc-ux644ux627ux6ccux647-xor-ux62fux648ux645}

\RL{مرحله} \RL{پنجم} \RL{یک} \RL{لایه} \RL{رمزنگاری} \RL{اضافه} \RL{است} \RL{که} \RL{عمق} \RL{امنیتی} \RL{الگوریتم} \RL{را} \RL{به}‌\RL{طور} \RL{قابل} \RL{توجهی} \RL{افزایش} \RL{می}‌\RL{دهد. در} \RL{این} \RL{مرحله،} \RL{تصویر} \RL{رمزنگاری}‌\RL{شده} \RL{از} \RL{مرحله} \RL{۴} \RL{یک} \RL{بار} \RL{دیگر} \RL{با} \RL{دنباله}‌\RL{ای} \RL{از} \RL{یک} \RL{سیستم} \RL{آشوبی} \RL{مستقل} \RL{عملیات} XOR \RL{می}‌\RL{شود.}

\RL{سیستم} \RL{مورد} \RL{استفاده} \RL{در} \RL{این} \RL{مرحله} \RL{یک} \RL{نقشه} Logistic \RL{مستقل} \RL{با} \RL{شرط} \RL{اولیه} x\textsubscript{final} \RL{است:}

{\def\LTcaptype{none} % do not increment counter
\begin{longtable}[]{@{}
  >{\raggedright\arraybackslash}p{(\linewidth - 2\tabcolsep) * \real{0.2756}}
  >{\raggedright\arraybackslash}p{(\linewidth - 2\tabcolsep) * \real{0.7105}}@{}}
\toprule\noalign{}
\begin{minipage}[b]{\linewidth}\raggedright
\begin{enumerate}
\def\labelenumi{\arabic{enumi}.}
\setcounter{enumi}{4}
\item
  \RL{لایه} XOR \RL{نهایی برای افزایش عمق رمزنگاری}
\end{enumerate}
\end{minipage} & \begin{minipage}[b]{\linewidth}\raggedleft
\[x_{\left\{ final_{\ k + 1} \right\}}\  = \ 4 \times \ x_{final_{k}} \times \ (1\  - \ x_{final_{k}}\ )\]

\[final_{\left\{ byte \right\}(k)}\  = \ (floor\left( x_{final_{k}}\  \times \ 10¹⁴) \right)\ \ \ mod\ 256\]

\[C2_{k}\  = \ C_{k}\ XOR\ final_{byte(K)}\ \ \]
\end{minipage} \\
\midrule\noalign{}
\endhead
\bottomrule\noalign{}
\endlastfoot
\end{longtable}
}

\RL{دلیل} \RL{طراحی} \RL{این} \RL{مرحله: یک} \RL{مهاجم} \RL{که} \RL{بتواند} \RL{لایه} \RL{اول} \RL{رمزنگاری} \RL{(عملیات} XOR \RL{با} \RL{سیستم} \RL{نمایی)} \RL{را} \RL{تحلیل} \RL{کند،} \RL{با} \RL{یک} \RL{لایه} \RL{رمزنگاری} \RL{مستقل} \RL{دیگر} \RL{مواجه} \RL{می}‌\RL{شود} \RL{که} \RL{با} \RL{کلید} \RL{متفاوتی} (x\_final) \RL{ایجاد} \RL{شده} \RL{است. این} \RL{رویکرد} \RL{دو}-\RL{لایه}‌\RL{ای،} \RL{پیچیدگی} \RL{تحلیل} \RL{رمز} \RL{را} \RL{به}‌\RL{صورت} \RL{توانی} \RL{افزایش} \RL{می}‌\RL{دهد.}

\RL{خروجی} \RL{نهایی} C2 \RL{که} \RL{در} \RL{ادامه} C\textsubscript{encrypted} \RL{نامیده} \RL{می}‌\RL{شود،} \RL{تصویر} \RL{رمزنگاری}‌\RL{شده} \RL{نهایی} \RL{است} \RL{که} \RL{هر} \RL{یک} \RL{از} \RL{سه} \RL{کانال} R\RL{،} G \RL{و} B \RL{به}‌\RL{طور} \RL{مستقل} \RL{از} \RL{این} \RL{فرآیند} \RL{عبور} \RL{کرده}‌\RL{اند.}

\begin{enumerate}
\def\labelenumi{\arabic{enumi}.}
\setcounter{enumi}{2}
\item
  \protect\phantomsection\label{_Toc232899903}{}\RL{فرآیند} \RL{رمزگشایی}
\end{enumerate}

\RL{از} \RL{آنجا} \RL{که} \RL{تمام} \RL{عملیات}‌\RL{های} \RL{اصلی} \RL{الگوریتم} \RL{(}XOR \RL{و} \RL{جایگشت)} \RL{کاملاً} \RL{برگشت}‌\RL{پذیر} \RL{هستند،} \RL{رمزگشایی} \RL{به}‌\RL{سادگی} \RL{با} \RL{اجرای} \RL{گام}‌\RL{ها} \RL{به} \RL{ترتیب} \RL{معکوس} \RL{انجام} \RL{می}‌\RL{شود.} \RL{مراحل} \RL{رمزگشایی} \RL{به}‌\RL{صورت} \RL{زیر} \RL{هستند:}

\begin{itemize}
\item
  \RL{مرحله} \RL{۱: بازسازی} \RL{تمام} \RL{دنباله}‌\RL{های} \RL{آشوبی} \RL{از} \RL{کلید} K \RL{(همان} \RL{شرایط} \RL{اولیه،} \RL{همان} warm-up\RL{)}
\item
  \RL{مرحله} \RL{۲: خنثی}‌\RL{سازی} \RL{لایه} XOR \RL{نهایی:} \(C_{k}\  = \ C2_{k}\ XOR\ {final}_{byte}(k)\)
\item
  \RL{مرحله} \RL{۳: خنثی}‌\RL{سازی} \RL{لایه} XOR \RL{نمایی:} \({P_{perm}}_{k}\  = \ C_{k}\ XOR\ {key}_{byte}(k)\)
\item
  \RL{مرحله} \RL{۴: اعمال} \RL{جایگشت} \RL{معکوس} \RL{با} \RL{بردار} \({perm}_{inv}:\ \ P_{original}\  = \ P_{perm}\lbrack{perm}_{inv}\rbrack\)
\item
  \RL{مرحله} \RL{۵: ترکیب} \RL{مجدد} \RL{کانال}‌\RL{های} R\RL{،} G \RL{و} B \RL{برای} \RL{بازیابی} \RL{تصویر} \RL{اصلی}
\end{itemize}

\RL{بردار} \RL{جایگشت} \RL{معکوس} perm\textsubscript{inv} \RL{از} \RL{رابطه} \RL{۳}-\RL{۹} \RL{محاسبه} \RL{می}‌\RL{شود:}

{\def\LTcaptype{none} % do not increment counter
\begin{longtable}[]{@{}
  >{\raggedright\arraybackslash}p{(\linewidth - 2\tabcolsep) * \real{0.4208}}
  >{\raggedright\arraybackslash}p{(\linewidth - 2\tabcolsep) * \real{0.5653}}@{}}
\toprule\noalign{}
\begin{minipage}[b]{\linewidth}\raggedright
\begin{enumerate}
\def\labelenumi{\arabic{enumi}.}
\item
  \RL{محاسبه} \RL{بردار} \RL{جایگشت} \RL{معکوس}
\end{enumerate}
\end{minipage} & \begin{minipage}[b]{\linewidth}\raggedleft
\[{perm}_{inv}\left\lbrack perm\lbrack k\rbrack \right\rbrack\  = \ k\ \ \ ,\ \forall\ k\  \in \ \{ 0,\ 1,\ .\ .\ .\ ,\ N - 1\}\]
\end{minipage} \\
\midrule\noalign{}
\endhead
\bottomrule\noalign{}
\endlastfoot
\end{longtable}
}

\RL{تضمین} \RL{معکوس}‌\RL{پذیری: چون} \RL{دنباله}‌\RL{های} \RL{آشوبی} \RL{مورد} \RL{استفاده} \RL{در} \RL{رمزگشایی} \RL{دقیقاً} \RL{همان} \RL{دنباله}‌\RL{های} \RL{رمزنگاری} \RL{هستند} \RL{(بازسازی}‌\RL{شده} \RL{از} \RL{کلید)} \RL{و} XOR \RL{عملیاتی} \RL{متقارن} \RL{است،} \RL{هیچ} \RL{گونه} \RL{خطا} \RL{یا} \RL{اریب}‌\RL{گذاری} \RL{در} \RL{فرآیند} \RL{رمزگشایی} \RL{رخ} \RL{نمی}‌\RL{دهد. این} \RL{تضمین} \RL{در} \RL{آزمایش}‌\RL{های} \RL{تجربی} \RL{با} \RL{محاسبه} PSNR = $\infty$ \RL{(بین} \RL{تصویر} \RL{اصلی} \RL{و} \RL{بازیابی}‌\RL{شده)} \RL{تأیید} \RL{می}‌\RL{شود{[}8{]}.}

\begin{enumerate}
\def\labelenumi{\arabic{enumi}.}
\item
  \protect\phantomsection\label{_Toc232899904}{}\RL{نمای} \RL{کلی} \RL{الگوریتم}
\end{enumerate}

\RL{جدول} \RL{۳}-\RL{۳} \RL{خلاصه}‌\RL{ای} \RL{از} \RL{کل} \RL{الگوریتم} \RL{رمزنگاری} \RL{و} \RL{رمزگشایی} \RL{پیشنهادی} \RL{را} \RL{ارائه} \RL{می}‌\RL{دهد:}

\begin{enumerate}
\def\labelenumi{\arabic{enumi}.}
\item
  \protect\phantomsection\label{_Toc232899947}{}\RL{خلاصه} \RL{مراحل} \RL{الگوریتم} \RL{رمزنگاری} \RL{پیشنهادی}
\end{enumerate}

{\def\LTcaptype{none} % do not increment counter
\begin{longtable}[]{@{}
  >{\raggedright\arraybackslash}p{(\linewidth - 8\tabcolsep) * \real{0.1764}}
  >{\raggedright\arraybackslash}p{(\linewidth - 8\tabcolsep) * \real{0.2505}}
  >{\raggedright\arraybackslash}p{(\linewidth - 8\tabcolsep) * \real{0.2538}}
  >{\raggedright\arraybackslash}p{(\linewidth - 8\tabcolsep) * \real{0.2350}}
  >{\raggedright\arraybackslash}p{(\linewidth - 8\tabcolsep) * \real{0.0844}}@{}}
\toprule\noalign{}
\begin{minipage}[b]{\linewidth}\centering
\textbf{\RL{سیستم آشوبی}}
\end{minipage} & \begin{minipage}[b]{\linewidth}\centering
\textbf{\RL{خروجی}}
\end{minipage} & \begin{minipage}[b]{\linewidth}\centering
\textbf{\RL{ورودی}}
\end{minipage} & \begin{minipage}[b]{\linewidth}\centering
\textbf{\RL{مرحله}}
\end{minipage} & \begin{minipage}[b]{\linewidth}\centering
\textbf{\RL{شماره}}
\end{minipage} \\
\midrule\noalign{}
\endhead
\bottomrule\noalign{}
\endlastfoot
--- & \RL{پارامترهای مقداردهی‌شده} & K = \{x₀,y₀,z₀,...\} & \RL{تولید کلید و آماده‌سازی} & \RL{۱} \\
Chen (3D) & \RL{تصویر} I\textsubscript{perm} & \RL{تصویر اصلی} I & \RL{جایگشت پیکسل‌ها} & \RL{۲} \\
Logistic (r=4) & \RL{بردار}

sel{[}1..N{]} ∈ \{1,...,9\} & x\textsubscript{sel} & \RL{انتخاب سیستم نمایی} & \RL{۳} \\
\RL{۹} \RL{سیستم} \RL{}ECM & \RL{تصویر} C & \(I_{perm}\  + \ sel\  + \ x_{ecm}\) & \RL{انتشار} XOR \RL{اول} & \RL{۴} \\
Logistic (r=4) & \RL{تصویر} C\textsubscript{encrypted} & C + x\textsubscript{final} & XOR \RL{نهایی} & \RL{۵} \\
\end{longtable}
}

\begin{enumerate}
\def\labelenumi{\arabic{enumi}.}
\setcounter{enumi}{1}
\item
  \protect\phantomsection\label{_Toc232899905}{}\RL{روش} \RL{پیاده}‌\RL{سازی} \RL{و} \RL{ابزارهای} \RL{نرم}‌\RL{افزاری}
\end{enumerate}

\RL{پیاده}‌\RL{سازی} \RL{الگوریتم} \RL{پیشنهادی} \RL{در} \RL{محیط} Python 3. 10 \RL{انجام} \RL{شده} \RL{است. انتخاب} Python \RL{به} \RL{دلیل} \RL{در} \RL{دسترس} \RL{بودن} \RL{کتابخانه}‌\RL{های} \RL{علمی} \RL{سطح}‌\RL{بالا} \RL{و} \RL{امکان} \RL{تحلیل} \RL{دقیق} \RL{عددی} \RL{بوده} \RL{است. کتابخانه}‌\RL{های} \RL{اصلی} \RL{مورد} \RL{استفاده} \RL{عبارتند} \RL{از:}

\begin{itemize}
\item
  NumPy\RL{: عملیات} \RL{برداری} \RL{و} \RL{ماتریسی} \RL{روی} \RL{آرایه}‌\RL{های} \RL{پیکسلی} \RL{با} \RL{دقت} float64
\item
  SciPy (solve\_ivp)\RL{: حل} \RL{عددی} \RL{معادلات} \RL{دیفرانسیل} \RL{سیستم} Chen \RL{با} \RL{روش} RK45
\item
  OpenCV\RL{: بارگذاری،} \RL{ذخیره}‌\RL{سازی} \RL{و} \RL{پیش}‌\RL{پردازش} \RL{تصاویر} \RL{رنگی}
\item
  Matplotlib\RL{: ترسیم} \RL{هیستوگرام،} \RL{نمودار} \RL{همبستگی} \RL{و} \RL{تحلیل}‌\RL{های} \RL{بصری}
\item
  Scikit-image\RL{: محاسبه} PSNR\RL{،} SSIM \RL{و} \RL{سایر} \RL{معیارهای} \RL{کیفیت} \RL{تصویر}
\end{itemize}

\RL{برای} \RL{ارزیابی} \RL{الگوریتم،} \RL{مجموعه} \RL{داده} USC-SIPI \RL{مورد} \RL{استفاده} \RL{قرار} \RL{گرفته} \RL{است. این} \RL{مجموعه} \RL{استاندارد} \RL{شامل} \RL{تصاویر} \RL{رنگی} \RL{متنوع} \RL{از} \RL{نظر} \RL{محتوا،} \RL{بافت} \RL{و} \RL{توزیع} \RL{آماری} \RL{است} \RL{که} \RL{مقایسه} \RL{عادلانه} \RL{با} \RL{سایر} \RL{روش}‌\RL{های} \RL{ادبیات} \RL{را} \RL{ممکن} \RL{می}‌\RL{سازد. تصاویر} \RL{در} \RL{ابعاد} \RL{۲۵۶}×\RL{۲۵۶} \RL{و} \RL{۵۱۲}×\RL{۵۱۲} \RL{با} \RL{فرمت} 24-bit RGB \RL{بارگذاری} \RL{می}‌\RL{شوند{[}3-4-9{]}.}

\begin{enumerate}
\def\labelenumi{\arabic{enumi}.}
\item
  \protect\phantomsection\label{_Toc232899906}{}\RL{معیارهای} \RL{ارزیابی}
\end{enumerate}

\RL{عملکرد} \RL{الگوریتم} \RL{بر} \RL{اساس} \RL{چهار} \RL{دسته} \RL{معیار} \RL{کمی} \RL{ارزیابی} \RL{می}‌\RL{شود:}

\begin{enumerate}
\def\labelenumi{\arabic{enumi}.}
\item
  \protect\phantomsection\label{_Toc232899907}{}\RL{آنتروپی} \RL{شانون}
\end{enumerate}

\RL{آنتروپی} \RL{شانون} (H) \RL{معیار} \RL{اصلی} \RL{تصادفی}‌\RL{بودن} \RL{توزیع} \RL{پیکسل}‌\RL{هاست. برای} \RL{تصاویر} \RL{رنگی} \RL{۸} \RL{بیتی،} \RL{مقدار} \RL{آرمانی} \RL{}H = 8 Byte \RL{است.}

{\def\LTcaptype{none} % do not increment counter
\begin{longtable}[]{@{}
  >{\raggedright\arraybackslash}p{(\linewidth - 2\tabcolsep) * \real{0.4931}}
  >{\raggedright\arraybackslash}p{(\linewidth - 2\tabcolsep) * \real{0.4931}}@{}}
\toprule\noalign{}
\begin{minipage}[b]{\linewidth}\raggedright
\begin{enumerate}
\def\labelenumi{\arabic{enumi}.}
\item
  \RL{آنتروپی} \RL{شانون}
\end{enumerate}
\end{minipage} & \begin{minipage}[b]{\linewidth}\raggedright
\[H\  = \  - \ \sum_{i = 0}^{255}p_{i} \times \log_{2}\left( p_{i} \right)\]
\end{minipage} \\
\midrule\noalign{}
\endhead
\bottomrule\noalign{}
\endlastfoot
\end{longtable}
}

\begin{enumerate}
\def\labelenumi{\arabic{enumi}.}
\setcounter{enumi}{1}
\item
  \protect\phantomsection\label{_Toc232899908}{}\RL{ضریب} \RL{همبستگی} \RL{پیکسل}‌\RL{های} \RL{مجاور}
\end{enumerate}

\RL{این} \RL{ضریب} \RL{میزان} \RL{ارتباط} \RL{خطی} \RL{بین} \RL{پیکسل}‌\RL{های} \RL{مجاور} \RL{(افقی،} \RL{عمودی،} \RL{قطری)} \RL{را} \RL{اندازه} \RL{می}‌\RL{گیرد. برای} \RL{تصاویر} \RL{رمزنگاری}‌\RL{شده} \RL{ایده}‌\RL{آل،} \RL{مقدار} \RL{آن} \RL{باید} \RL{به} \RL{صفر} \RL{نزدیک} \RL{باشد:}

{\def\LTcaptype{none} % do not increment counter
\begin{longtable}[]{@{}
  >{\raggedright\arraybackslash}p{(\linewidth - 2\tabcolsep) * \real{0.4931}}
  >{\raggedright\arraybackslash}p{(\linewidth - 2\tabcolsep) * \real{0.4931}}@{}}
\toprule\noalign{}
\begin{minipage}[b]{\linewidth}\raggedright
\begin{enumerate}
\def\labelenumi{\arabic{enumi}.}
\item
  \RL{ضریب} \RL{همبستگی} \RL{پیرسون}
\end{enumerate}
\end{minipage} & \begin{minipage}[b]{\linewidth}\raggedleft
\[r\  = \ \frac{\sum_{}^{}{\left( x_{i} - \ \overline{\ x} \right)\left( y_{i} - \ \overline{\ y}\  \right)}}{\sqrt{}\left\{ \sum\ \left( x_{i} - \overline{\ x}\  \right)^{2} \times \ \sum\ \left( y_{i} - \ \overline{\ y} \right)^{2} \right\}}\]
\end{minipage} \\
\midrule\noalign{}
\endhead
\bottomrule\noalign{}
\endlastfoot
\end{longtable}
}

\begin{enumerate}
\def\labelenumi{\arabic{enumi}.}
\setcounter{enumi}{2}
\item
  \protect\phantomsection\label{_Toc232899909}{}\RL{معیارهای} NPCR \RL{و} UACI
\end{enumerate}

\RL{این} \RL{دو} \RL{معیار} \RL{حساسیت} \RL{الگوریتم} \RL{را} \RL{نسبت} \RL{به} \RL{تغییرات} \RL{بسیار} \RL{جزئی} \RL{(یک} \RL{بیت)} \RL{در} \RL{تصویر} \RL{اصلی} \RL{یا} \RL{کلید} \RL{ارزیابی} \RL{می}‌\RL{کنند:}

NPCR \textbf{\RL{(نرخ}} \textbf{\RL{تغییر}} \RL{\textbf{پیکسل):} نسبت} \RL{پیکسل}‌\RL{هایی} \RL{که} \RL{بین} \RL{دو} \RL{تصویر} \RL{رمزنگاری}‌\RL{شده} \RL{با} \RL{کلیدهای} \RL{اندکی} \RL{متفاوت،} \RL{تغییر} \RL{کرده}‌\RL{اند:}

{\def\LTcaptype{none} % do not increment counter
\begin{longtable}[]{@{}
  >{\raggedright\arraybackslash}p{(\linewidth - 2\tabcolsep) * \real{0.2916}}
  >{\raggedright\arraybackslash}p{(\linewidth - 2\tabcolsep) * \real{0.6945}}@{}}
\toprule\noalign{}
\begin{minipage}[b]{\linewidth}\raggedright
\begin{enumerate}
\def\labelenumi{\arabic{enumi}.}
\item
  \RL{فرمول} \RL{صحیح} NPCR \RL{(با} \RL{متغیر} \RL{باینری} D\RL{)}
\end{enumerate}
\end{minipage} & \begin{minipage}[b]{\linewidth}\raggedright
\(NPCR\  = \frac{\left\lbrack \Sigma ᵢⱼ\ D(i,j) \right\rbrack}{(W \times H)}\ \  \times \ \lr{100\%}\) \RL{}

\RL{در غیر این صورت}

\RL{}\(D(i,j)\  = \ 0\ \ if\ \ C₁(i,j) = C₂(i,j)\ \ ;\ \ D(i,j)\  = \ 1\ \ \)
\end{minipage} \\
\midrule\noalign{}
\endhead
\bottomrule\noalign{}
\endlastfoot
\end{longtable}
}

UACI \textbf{\RL{(میانگین}} \textbf{\RL{شدت}} \RL{\textbf{تغییر):} میانگین} \RL{اختلاف} \RL{مقادیر} \RL{پیکسل}‌\RL{ها} \RL{بین} \RL{دو} \RL{تصویر} \RL{رمزنگاری}‌\RL{شده:}

{\def\LTcaptype{none} % do not increment counter
\begin{longtable}[]{@{}
  >{\raggedright\arraybackslash}p{(\linewidth - 2\tabcolsep) * \real{0.3402}}
  >{\raggedright\arraybackslash}p{(\linewidth - 2\tabcolsep) * \real{0.6459}}@{}}
\toprule\noalign{}
\begin{minipage}[b]{\linewidth}\raggedright
\begin{enumerate}
\def\labelenumi{\arabic{enumi}.}
\setcounter{enumi}{1}
\item
  \RL{فرمول} UACI
\end{enumerate}
\end{minipage} & \begin{minipage}[b]{\linewidth}\raggedleft
\[UACI\  = \frac{\left\lbrack \Sigma ᵢⱼ\ \left| C^{1}(i,j) - C^{2}(i,j) \right| \right\rbrack}{(255 \times W \times H)}\ \  \times \ \lr{100\%}\]
\end{minipage} \\
\midrule\noalign{}
\endhead
\bottomrule\noalign{}
\endlastfoot
\end{longtable}
}

\textbf{\RL{مقادیر آرمانی: NPCR \textgreater{} \lr{99.60\%} و UACI ≈ \lr{33.46\%} که نشان‌دهنده حساسیت بالا به تغییرات جزئی (اثر بهمنی) هستند{[}10{]}.} \RL{}}\RL{}

\begin{enumerate}
\def\labelenumi{\arabic{enumi}.}
\setcounter{enumi}{3}
\item
  \protect\phantomsection\label{_Toc232899910}{}\RL{تحلیل} \RL{زمان} \RL{اجرا}
\end{enumerate}

\RL{زمان} \RL{رمزنگاری} \RL{و} \RL{رمزگشایی} \RL{برای} \RL{ابعاد} \RL{مختلف} \RL{تصویر} \RL{اندازه}‌\RL{گیری} \RL{شده} \RL{و} \RL{تأثیر} \RL{اندازه} \RL{تصویر} \RL{بر} \RL{سرعت} \RL{پردازش} \RL{بررسی} \RL{می}‌\RL{شود. همچنین} \RL{سرعت} \RL{پردازش} \RL{به} \RL{صورت} MBps \RL{(مگابایت} \RL{در} \RL{ثانیه)} \RL{گزارش} \RL{می}‌\RL{گردد} \RL{تا} \RL{مقایسه} \RL{با} \RL{سایر} \RL{روش}‌\RL{ها} \RL{در} \RL{ادبیات} \RL{ممکن} \RL{باشد.}

\begin{enumerate}
\def\labelenumi{\arabic{enumi}.}
\setcounter{enumi}{1}
\item
  \protect\phantomsection\label{_Toc232899911}{}\RL{جمع}‌\RL{بندی} \RL{فصل}
\end{enumerate}

\RL{در} \RL{این} \RL{فصل،} \RL{الگوریتم} \RL{رمزنگاری} \RL{پیشنهادی} \RL{با} \RL{تمام} \RL{جزئیات} \RL{ریاضی} \RL{و} \RL{طراحی} \RL{معرفی} \RL{شد. نقاط} \RL{کلیدی} \RL{معماری} \RL{پیشنهادی} \RL{عبارتند} \RL{از:}

\begin{itemize}
\item
  \RL{استفاده} \RL{از} \RL{سیستم} Chen \RL{سه}‌\RL{بعدی} \RL{با} \RL{پارامترهای} \RL{استاندارد} \RL{برای} \RL{ایجاد} \RL{جایگشت} \RL{ساختاری} \RL{قوی}
\item
  \RL{انتخاب} \RL{پویا} \RL{و} \RL{مستقل} \RL{سیستم} \RL{آشوبی} \RL{نمایی} \RL{برای} \RL{هر} \RL{پیکسل،} \RL{که} \RL{قابلیت} \RL{پیش}‌\RL{بینی} \RL{را} \RL{به} \RL{حداقل} \RL{می}‌\RL{رساند}
\item
  \RL{دو} \RL{لایه} XOR \RL{مستقل} \RL{که} \RL{مقاومت} \RL{در} \RL{برابر} \RL{حملات} \RL{تفاضلی} \RL{را} \RL{دوچندان} \RL{می}‌\RL{کند}
\item
  \RL{طراحی} \RL{بدون} \RL{وابستگی} \RL{دنباله} \RL{کلید} \RL{به} \RL{مقادیر} \RL{پیکسل،} \RL{که} \RL{معکوس}‌\RL{پذیری} \RL{کامل} \RL{را} \RL{تضمین} \RL{می}‌\RL{کند}
\item
  \RL{فضای} \RL{کلید} \RL{بزرگتر} \RL{از} 2\^{}348 \RL{که} \RL{حمله} \RL{جستجوی} \RL{فراگیر} \RL{را} \RL{ناممکن} \RL{می}‌\RL{سازد}
\end{itemize}

\RL{در} \RL{فصل} \RL{چهارم،} \RL{نتایج} \RL{تجربی} \RL{اجرای} \RL{این} \RL{الگوریتم} \RL{بر} \RL{روی} \RL{مجموعه} \RL{داده} USC-SIPI \RL{ارائه} \RL{و} \RL{تحلیل} \RL{می}‌\RL{شود} \RL{و} \RL{عملکرد} \RL{روش} \RL{پیشنهادی} \RL{با} \RL{روش}‌\RL{های} \RL{مرجع} \RL{در} \RL{ادبیات} \RL{مقایسه} \RL{خواهد} \RL{شد.}

\begin{enumerate}
\def\labelenumi{\arabic{enumi}.}
\setcounter{enumi}{3}
\item
  \protect\phantomsection\label{_Toc232899912}{}\RL{پیاده سازی الگوریتم پیشنهادی}
\end{enumerate}

\RL{\hfill\break
}

\begin{enumerate}
\def\labelenumi{\arabic{enumi}.}
\item
  \protect\phantomsection\label{_Toc232899913}{}\RL{مقدمه}
\end{enumerate}

\RL{در} \RL{این} \RL{فصل،} \RL{نتایج} \RL{اجرای} \RL{الگوریتم} \RL{رمزنگاری} \RL{پیشنهادی} \RL{بر} \RL{روی} \RL{چهار} \RL{تصویر} \RL{استاندارد} \RL{از} \RL{مجموعه} \RL{داده} USC-SIPI \RL{ارائه} \RL{و} \RL{تحلیل} \RL{می}‌\RL{شود. تمام} \RL{آزمایش}‌\RL{ها} \RL{در} \RL{محیط} Python 3.10 \RL{با} \RL{استفاده} \RL{از} \RL{کتابخانه}‌\RL{های} NumPy\RL{،} SciPy \RL{و} OpenCV \RL{بر} \RL{روی} \RL{تصاویر} \RL{رنگی} \RL{۵۱۲}×\RL{۵۱۲} \RL{و} \RL{۲۵۶}×\RL{۲۵۶} \RL{انجام} \RL{شده} \RL{است. ارزیابی} \RL{عملکرد} \RL{شامل} \RL{تحلیل} \RL{بصری} \RL{تصاویر،} \RL{هیستوگرام،} \RL{آنتروپی} \RL{شانون،} \RL{ضریب} \RL{همبستگی} \RL{پیکسل}‌\RL{های} \RL{مجاور،} \RL{معیارهای} \RL{حساسیت} NPCR \RL{و} UACI\RL{،} \RL{و} \RL{زمان} \RL{اجرا} \RL{می}‌\RL{باشد}.

\RL{تصاویر} \RL{مورد} \RL{آزمایش} \RL{که} \RL{همگی} \RL{از} \RL{مجموعه} USC-SIPI \RL{انتخاب} \RL{شده}‌\RL{اند،} \RL{طیف} \RL{متنوعی} \RL{از} \RL{ویژگی}‌\RL{های} \RL{آماری} \RL{را} \RL{پوشش} \RL{می}‌\RL{دهند}: Airplane \RL{با} \RL{محتوای} \RL{هوایی} \RL{یکنواخت،} Baboon \RL{با} \RL{بافت} \RL{پرجزئیات،} Peppers \RL{با} \RL{گستره} \RL{رنگی} \RL{وسیع} \RL{و} Tree \RL{با} \RL{ساختار} \RL{طبیعی. این} \RL{تنوع،} \RL{ارزیابی} \RL{جامع}‌\RL{تری} \RL{از} \RL{رفتار} \RL{الگوریتم} \RL{فراهم} \RL{می}‌\RL{آورد}.

\begin{enumerate}
\def\labelenumi{\arabic{enumi}.}
\setcounter{enumi}{1}
\item
  \protect\phantomsection\label{_Toc232899914}{}\RL{تنظیمات} \RL{آزمایش} \RL{و} \RL{پارامترهای} \RL{کلید}
\end{enumerate}

\RL{پارامترهای} \RL{کلید} \RL{رمزنگاری} \RL{مورد} \RL{استفاده} \RL{در} \RL{تمام} \RL{آزمایش}‌\RL{ها} \RL{در} \RL{جدول} \RL{۴}-\RL{۱} \RL{آمده} \RL{است. این} \RL{مقادیر} \RL{با} \RL{دقت} \RL{۱۵} \RL{رقم} \RL{اعشاری} \RL{تعریف} \RL{شده}‌\RL{اند} \RL{تا} \RL{فضای} \RL{کلید} \RL{بهینه}‌\RL{ای} \RL{برای} \RL{ارزیابی} \RL{امنیتی} \RL{فراهم} \RL{شود}:

\begin{enumerate}
\def\labelenumi{\arabic{enumi}.}
\item
  \protect\phantomsection\label{_Toc232899948}{}\RL{پارامترهای} \RL{کلید} \RL{رمزنگاری} \RL{مورد} \RL{استفاده} \RL{در} \RL{آزمایش}‌\RL{ها}
\end{enumerate}

{\def\LTcaptype{none} % do not increment counter
\begin{longtable}[]{@{}
  >{\centering\arraybackslash}p{(\linewidth - 4\tabcolsep) * \real{0.1354}}
  >{\centering\arraybackslash}p{(\linewidth - 4\tabcolsep) * \real{0.3320}}
  >{\centering\arraybackslash}p{(\linewidth - 4\tabcolsep) * \real{0.5325}}@{}}
\toprule\noalign{}
\begin{minipage}[b]{\linewidth}\centering
\RL{پارامتر}
\end{minipage} & \begin{minipage}[b]{\linewidth}\centering
\RL{مقدار}
\end{minipage} & \begin{minipage}[b]{\linewidth}\centering
\RL{نقش} \RL{در} \RL{الگوریتم}
\end{minipage} \\
\midrule\noalign{}
\endhead
\bottomrule\noalign{}
\endlastfoot
x₀ & 0.123456789012345 & \RL{شرط} \RL{اولیه} \RL{محور} x \RL{سیستم} Chen \\
y₀ & 0.987654321098765 & \RL{شرط} \RL{اولیه} \RL{محور} y \RL{سیستم} Chen \\
z₀ & 12.345678901234 & \RL{شرط} \RL{اولیه} \RL{محور} z \RL{سیستم} Chen \\
x\textsubscript{sel} & 0.456789012345678 & \RL{شرط} \RL{اولیه} \RL{انتخاب}‌\RL{گر} \RL{سیستم} \RL{نمایی} \\
x\textsubscript{ecm} & 0.789012345678901 & \RL{شرط} \RL{اولیه} \RL{مشترک} \RL{۹} \RL{سیستم} \RL{نمایی} \\
μ & 3.8 & \RL{پارامتر} \RL{نمایی} ECM \RL{(در} \RL{ناحیه} \RL{آشوب} \RL{پایدار)} \\
x\textsubscript{final} & 0.234567890123456 & \RL{شرط} \RL{اولیه} \RL{لایه} XOR \RL{نهایی} \\
\end{longtable}
}

\begin{enumerate}
\def\labelenumi{\arabic{enumi}.}
\setcounter{enumi}{2}
\item
  \protect\phantomsection\label{_Toc232899915}{}\RL{تحلیل} \RL{بصری} \RL{تصاویر}
\end{enumerate}

\RL{شکل} \RL{۴}-\RL{۱} \RL{نمایش} \RL{بصری} \RL{سه} \RL{مرحله} \RL{اصلی} \RL{الگوریتم} \RL{را} \RL{برای} \RL{هر} \RL{چهار} \RL{تصویر} \RL{نشان} \RL{می}‌\RL{دهد}: \RL{تصویر} \RL{اصلی،} \RL{تصویر} \RL{رمزنگاری}‌\RL{شده} \RL{و} \RL{تصویر} \RL{بازیابی}‌\RL{شده} \RL{پس} \RL{از} \RL{رمزگشایی}.

\includegraphics[width=6.04167in,height=6.66667in]{media/image5.png}

\begin{enumerate}
\def\labelenumi{\arabic{enumi}.}
\item
  \protect\phantomsection\label{_Toc232899808}{}\RL{نمایش} \RL{بصری} \RL{تصاویر} \RL{اصلی،} \RL{رمزنگاری}‌\RL{شده} \RL{و} \RL{بازیابی}‌\RL{شده}
\end{enumerate}

\RL{همان}‌\RL{طور} \RL{که} \RL{در} \RL{شکل} \RL{۴}-\RL{۱} \RL{مشاهده} \RL{می}‌\RL{شود،} \RL{تصاویر} \RL{رمزنگاری}‌\RL{شده} \RL{هیچ}‌\RL{گونه} \RL{اطلاعاتی} \RL{از} \RL{محتوای} \RL{اصلی} \RL{را} \RL{آشکار} \RL{نمی}‌\RL{کنند} \RL{و} \RL{ظاهری} \RL{کاملاً} \RL{تصادفی} \RL{و} \RL{مشابه} \RL{نویز} \RL{دارند. نکته} \RL{قابل} \RL{توجه،} \RL{تفاوت} \RL{بصری} \RL{واضح} \RL{بین} \RL{تصاویر} \RL{رمزنگاری}‌\RL{شده} \RL{با} \RL{محتوای} \RL{متفاوت} \RL{(مثلاً} Airplane \RL{در} \RL{مقابل} Baboon\RL{)} \RL{است} \RL{که} \RL{نشان} \RL{می}‌\RL{دهد} \RL{الگوریتم} \RL{هیچ}‌\RL{گونه} \RL{الگوی} \RL{مشترکی} \RL{در} \RL{خروجی} \RL{تولید} \RL{نمی}‌\RL{کند}.

\RL{مهم}‌\RL{تر} \RL{از} \RL{همه،} \RL{ستون} \RL{سوم} \RL{(تصویر} \RL{بازیابی}‌\RL{شده)} \RL{در} \RL{تمام} \RL{موارد} \RL{با} \RL{تصویر} \RL{اصلی} \RL{به}‌\RL{طور} \RL{کامل} \RL{منطبق} \RL{است. مقدار} MSE \RL{صفر} \RL{و} PSNR \RL{برابر} \RL{با} \RL{بی}‌\RL{نهایت} \RL{برای} \RL{هر} \RL{چهار} \RL{تصویر،} \RL{معکوس}‌\RL{پذیری} \RL{\lr{۱۰۰٪}} \RL{الگوریتم} \RL{را} \RL{تأیید} \RL{می}‌\RL{کند}.

\begin{enumerate}
\def\labelenumi{\arabic{enumi}.}
\setcounter{enumi}{3}
\item
  \protect\phantomsection\label{_Toc232899916}{}\RL{تحلیل هیستوگرام}
\end{enumerate}

\RL{هیستوگرام} \RL{توزیع} \RL{مقادیر} \RL{پیکسل}‌\RL{ها} \RL{یکی} \RL{از} \RL{ابزارهای} \RL{اصلی} \RL{تحلیل} \RL{امنیتی} \RL{است. در} \RL{تصویر} \RL{اصلی،} \RL{هیستوگرام} \RL{معمولاً} \RL{نامتقارن} \RL{و} \RL{دارای} \RL{قله}‌\RL{های} \RL{مشخص} \RL{است} \RL{که} \RL{اطلاعات} \RL{آماری} \RL{قابل} \RL{بهره}‌\RL{برداری} \RL{برای} \RL{مهاجم} \RL{فراهم} \RL{می}‌\RL{کند. پس} \RL{از} \RL{رمزنگاری،} \RL{هیستوگرام} \RL{باید} \RL{به}‌\RL{صورت} \RL{یکنواخت} \RL{در} \RL{تمام} \RL{بازه} \RL{۰} \RL{تا} \RL{۲۵۵} \RL{توزیع} \RL{شود}.

\includegraphics[width=5.89222in,height=3.61136in]{media/image6.png}

\begin{enumerate}
\def\labelenumi{\arabic{enumi}.}
\item
  \protect\phantomsection\label{_Toc232899809}{}\RL{هیستوگرام} \RL{مقایسه}‌\RL{ای} \RL{تصاویر} \RL{اصلی} \RL{(چپ)} \RL{و} \RL{رمزنگاری}‌\RL{شده} \RL{(راست)برای} \RL{سه} \RL{کانال} \RL{رنگی}
\end{enumerate}

\RL{شکل} \RL{۴}-\RL{۲} \RL{به}‌\RL{روشنی} \RL{نشان} \RL{می}‌\RL{دهد} \RL{که} \RL{هیستوگرام} \RL{تمام} \RL{تصاویر} \RL{رمزنگاری}‌\RL{شده،} \RL{در} \RL{هر} \RL{سه} \RL{کانال} B\RL{،} G \RL{و} R\RL{،} \RL{به} \RL{توزیع} \RL{یکنواخت} \RL{نزدیک} \RL{شده}‌\RL{اند. این} \RL{یکنواختی} \RL{به} \RL{این} \RL{معناست} \RL{که} \RL{مهاجم} \RL{هیچ} \RL{اطلاعاتی} \RL{از} \RL{تحلیل} \RL{آماری} \RL{توزیع} \RL{پیکسل}‌\RL{ها} \RL{به}‌\RL{دست} \RL{نمی}‌\RL{آورد}.

\RL{برخلاف تصویر Airplane که هیستوگرام اصلی آن دارای توزیع کاملاً نامتقارن و دارای قله‌های برجسته است، تصویر رمزنگاری‌شده آن هیستوگرامی کاملاً مسطح دارد. این مقاومت در برابر حملات آماری مبتنی بر هیستوگرام را تضمین می‌کند.}

\begin{enumerate}
\def\labelenumi{\arabic{enumi}.}
\setcounter{enumi}{4}
\item
  \protect\phantomsection\label{_Toc232899917}{}\RL{تحلیل} \RL{آنتروپی} \RL{شانون}
\end{enumerate}

\RL{آنتروپی} \RL{شانون} \RL{مهم}‌\RL{ترین} \RL{معیار} \RL{کمی} \RL{برای} \RL{سنجش} \RL{تصادفی}‌\RL{بودن} \RL{داده}‌\RL{های} \RL{رمزنگاری}‌\RL{شده} \RL{است. برای} \RL{تصاویر} \RL{۸} \RL{بیتی،} \RL{مقدار} \RL{ایده}‌\RL{آل} \RL{۸} \RL{بیت} \RL{است. نتایج} \RL{تفصیلی} \RL{در} \RL{جدول} \RL{۴}-\RL{۲} \RL{ارائه} \RL{شده}:

\begin{enumerate}
\def\labelenumi{\arabic{enumi}.}
\item
  \protect\phantomsection\label{_Toc232899949}{}\RL{آنتروپی} \RL{شانون} \RL{قبل} \RL{و} \RL{بعد} \RL{از} \RL{رمزنگاری} \RL{(بر} \RL{حسب} \RL{بیت)}
\end{enumerate}

{\def\LTcaptype{none} % do not increment counter
\begin{longtable}[]{@{}
  >{\raggedright\arraybackslash}p{(\linewidth - 8\tabcolsep) * \real{0.1185}}
  >{\raggedright\arraybackslash}p{(\linewidth - 8\tabcolsep) * \real{0.2579}}
  >{\raggedright\arraybackslash}p{(\linewidth - 8\tabcolsep) * \real{0.2190}}
  >{\raggedright\arraybackslash}p{(\linewidth - 8\tabcolsep) * \real{0.2166}}
  >{\raggedright\arraybackslash}p{(\linewidth - 8\tabcolsep) * \real{0.1774}}@{}}
\toprule\noalign{}
\begin{minipage}[b]{\linewidth}\raggedright
\RL{تصویر}
\end{minipage} & \begin{minipage}[b]{\linewidth}\raggedright
\RL{کانال} B \RL{(اصلی}→\RL{رمز)}
\end{minipage} & \begin{minipage}[b]{\linewidth}\raggedright
\RL{کانال} G \RL{(اصلی}→\RL{رمز)}
\end{minipage} & \begin{minipage}[b]{\linewidth}\raggedright
\RL{کانال} R \RL{(اصلی}→\RL{رمز)}
\end{minipage} & \begin{minipage}[b]{\linewidth}\raggedright
\RL{میانگین} \RL{رمزشده}
\end{minipage} \\
\midrule\noalign{}
\endhead
\bottomrule\noalign{}
\endlastfoot
Airplane & 7.2542 → 8.0 & 7.1491 → 8.0 & 5.8825 → 8.0 & 8.0 \\
Baboon & 1.0000 → 8.0 & 7.5754 → 8.0 & 7.2855 → 8.0 & 8.0 \\
Peppers & 7.9931 → 8.0 & 7.9931 → 8.0 & 7.6718 → 8.0 & 8.0 \\
Tree & 7.6075 → 8.0 & 7.6398 → 8.0 & 6.3472 → 8.0 & 8.0 \\
\end{longtable}
}

\includegraphics[width=6.04167in,height=2.70833in]{media/image7.png}

\begin{enumerate}
\def\labelenumi{\arabic{enumi}.}
\item
  \protect\phantomsection\label{_Toc232899810}{}\RL{نمودار} \RL{مقایسه}‌\RL{ای} \RL{آنتروپی} \RL{شانون} \RL{تصاویر} \RL{اصلی} \RL{و} \RL{رمزنگاری}‌\RL{شده}
\end{enumerate}

نتایج جدول ۴-۲ و شکل ۴-۳ نشان می‌دهند که الگوریتم پیشنهادی به‌طور مؤثری آنتروپی تمام تصاویر را به مقادیر بسیار نزدیک به حداکثر ممکن (۸ بیت) رسانده است. میانگین آنتروپی برای چهار تصویر برابر با ۸.۰ بیت است که انحراف از مقدار ایده‌آل ۰.۰ بیت می‌باشد.

قابل توجه است که تصویر Airplane بیشترین بهبود آنتروپی را نشان می‌دهد (از ۶.۷۶ به ۷.۹۹۹۳ بیت) چرا که هیستوگرام اصلی آن بیشترین ناهمواری را داشت. این نشان می‌دهد الگوریتم صرف نظر از ساختار اولیه تصویر، توزیع یکنواختی در خروجی ایجاد می‌کند.

\begin{enumerate}
\def\labelenumi{\arabic{enumi}.}
\setcounter{enumi}{5}
\item
  \protect\phantomsection\label{_Toc232899918}{}\RL{تحلیل} \RL{همبستگی} \RL{پیکسل}‌\RL{های} \RL{مجاور}
\end{enumerate}

\RL{ضریب} \RL{همبستگی} \RL{پیکسل}‌\RL{های} \RL{مجاور} \RL{یکی} \RL{از} \RL{معیارهای} \RL{اساسی} \RL{در} \RL{ارزیابی} \RL{کارایی} \RL{مرحله} \RL{جایگشت} \RL{است. در} \RL{تصاویر} \RL{طبیعی،} \RL{پیکسل}‌\RL{های} \RL{مجاور} \RL{معمولاً} \RL{همبستگی} \RL{بالایی} \RL{دارند} \RL{(نزدیک} \RL{به} \RL{۱). پس} \RL{از} \RL{رمزنگاری،} \RL{این} \RL{ضریب} \RL{باید} \RL{به} \RL{صفر} \RL{نزدیک} \RL{شود} \RL{تا} \RL{تحلیل} \RL{ساختاری} \RL{ناممکن} \RL{گردد}.

\begin{enumerate}
\def\labelenumi{\arabic{enumi}.}
\item
  \protect\phantomsection\label{_Toc232899950}{}\RL{ضریب} \RL{همبستگی} \RL{پیکسل}‌\RL{های} \RL{مجاور} \RL{در} \RL{سه} \RL{جهت} \RL{(کانال} G\RL{)}
\end{enumerate}

{\def\LTcaptype{none} % do not increment counter
\begin{longtable}[]{@{}
  >{\raggedright\arraybackslash}p{(\linewidth - 8\tabcolsep) * \real{0.2201}}
  >{\raggedright\arraybackslash}p{(\linewidth - 8\tabcolsep) * \real{0.1886}}
  >{\raggedright\arraybackslash}p{(\linewidth - 8\tabcolsep) * \real{0.1906}}
  >{\raggedright\arraybackslash}p{(\linewidth - 8\tabcolsep) * \real{0.1932}}
  >{\raggedright\arraybackslash}p{(\linewidth - 8\tabcolsep) * \real{0.2076}}@{}}
\toprule\noalign{}
\endhead
\bottomrule\noalign{}
\endlastfoot
\RL{کاهش} \RL{همبستگی} & \RL{ضریب} \RL{رمزشده} & \RL{ضریب} \RL{اصلی} & \RL{جهت} & \RL{تصویر} \\
\lr{99.20\%} & -0.0080 & 0.9998 & \RL{افقی} & Airplane \\
\lr{99.62\%} & -0.0038 & 1.0000 & \RL{عمودی} & Airplane \\
\lr{97.40\%} & -0.0260 & 0.9998 & \RL{قطری} & Airplane \\
\lr{99.44\%} & 0.0056 & 1.0000 & \RL{افقی} & Baboon \\
\lr{99.67\%} & 0.0033 & 0.9948 & \RL{عمودی} & Baboon \\
\lr{98.52\%} & 0.0147 & 0.9948 & \RL{قطری} & Baboon \\
\lr{98.76\%} & -0.0124 & 1.0000 & \RL{افقی} & Peppers \\
\lr{98.97\%} & 0.0103 & 1.0000 & \RL{عمودی} & Peppers \\
\lr{99.63\%} & -0.0037 & 1.0000 & \RL{قطری} & Peppers \\
\lr{99.51\%} & -0.0087 & 0.9997 & \RL{افقی} & Tree \\
\lr{99.55\%} & -0.0045 & 0.9997 & \RL{عمودی} & Tree \\
\lr{98.43\%} & 0.0157 & 0.9987 & \RL{قطری} & Tree \\
\end{longtable}
}

نتایج جدول ۴-۳ و شکل‌های ۴-۴ و ۴-۵ نشان می‌دهند که الگوریتم پیشنهادی در تمام موارد موفق به کاهش چشمگیر ضریب همبستگی شده است. میانگین ضریب همبستگی تصاویر رمزنگاری‌شده در سه جهت برای چهار تصویر برابر با ‎۰.۰۱۲۶‎ است، در مقابل میانگین ‎۰.۹۹۸۹ برای تصاویر اصلی.

\RL{نمودارهای} \RL{پراکندگی} \RL{شکل} \RL{۴}-\RL{۴} \RL{این} \RL{واقعیت} \RL{را} \RL{به}‌\RL{خوبی} \RL{نمایش} \RL{می}‌\RL{دهند}: \RL{در} \RL{تصاویر} \RL{اصلی،} \RL{نقاط} \RL{پراکندگی} \RL{به}‌\RL{صورت} \RL{یک} \RL{خط} \RL{قطری} \RL{متمرکز} \RL{هستند} \RL{که} \RL{همبستگی} \RL{بالا} \RL{را} \RL{نشان} \RL{می}‌\RL{دهد. در} \RL{مقابل،} \RL{در} \RL{تصاویر} \RL{رمزنگاری}‌\RL{شده،} \RL{نقاط} \RL{به}‌\RL{صورت} \RL{کاملاً} \RL{تصادفی} \RL{در} \RL{سراسر} \RL{فضا} \RL{پراکنده} \RL{شده}‌\RL{اند} \RL{که} \RL{نشانه} \RL{عدم} \RL{همبستگی} \RL{است}.

\includegraphics[width=6.0947in,height=3.63716in]{media/image8.png}

\begin{enumerate}
\def\labelenumi{\arabic{enumi}.}
\item
  \protect\phantomsection\label{_Toc232899811}{}\RL{نمودار} \RL{پراکندگی} \RL{همبستگی} \RL{پیکسل}‌\RL{ها} --- \RL{اصلی} \RL{(چپ)} \RL{و} \RL{رمزنگاری}‌\RL{شده} \RL{(راست)}
\end{enumerate}

\includegraphics[width=6.04167in,height=2.5in]{media/image9.png}

\begin{enumerate}
\def\labelenumi{\arabic{enumi}.}
\setcounter{enumi}{1}
\item
  \protect\phantomsection\label{_Toc232899812}{}\RL{مقایسه} \RL{ضریب} \RL{همبستگی} \RL{در} \RL{سه} \RL{جهت} \RL{برای} \RL{چهار} \RL{تصویر}
\end{enumerate}

\begin{enumerate}
\def\labelenumi{\arabic{enumi}.}
\setcounter{enumi}{6}
\item
  \protect\phantomsection\label{_Toc232899919}{}\RL{تحلیل} \RL{حساسیت} \RL{به} \RL{کلید} NPCR \RL{و} UACI
\end{enumerate}

\RL{معیارهای NPCR و UACI مقاومت الگوریتم در برابر حملات دیفرانسیلی را ارزیابی می‌کنند. برای این آزمون، دو تصویر رمزنگاری‌شده با کلیدهای اندکی متفاوت (تغییر ۱۰⁻¹⁵ در x₀) تولید شده و با هم مقایسه می‌شوند. مقادیر آرمانی عبارتند از: NPCR \textgreater{} \lr{99.60\%} و UACI ≈ \lr{33.46\%}} \RL{}

\begin{enumerate}
\def\labelenumi{\arabic{enumi}.}
\item
  \protect\phantomsection\label{_Toc232899951}{}\RL{نتایج} \RL{آزمون} NPCR \RL{و} UACI
\end{enumerate}

{\def\LTcaptype{none} % do not increment counter
\begin{longtable}[]{@{}
  >{\raggedright\arraybackslash}p{(\linewidth - 10\tabcolsep) * \real{0.1281}}
  >{\raggedright\arraybackslash}p{(\linewidth - 10\tabcolsep) * \real{0.1616}}
  >{\raggedright\arraybackslash}p{(\linewidth - 10\tabcolsep) * \real{0.1452}}
  >{\raggedright\arraybackslash}p{(\linewidth - 10\tabcolsep) * \real{0.1290}}
  >{\raggedright\arraybackslash}p{(\linewidth - 10\tabcolsep) * \real{0.2098}}
  >{\raggedright\arraybackslash}p{(\linewidth - 10\tabcolsep) * \real{0.2262}}@{}}
\toprule\noalign{}
\begin{minipage}[b]{\linewidth}\raggedright
\RL{تصویر}
\end{minipage} & \begin{minipage}[b]{\linewidth}\raggedright
\RL{ابعاد}
\end{minipage} & \begin{minipage}[b]{\linewidth}\raggedright
NPCR \RL{(\%)}
\end{minipage} & \begin{minipage}[b]{\linewidth}\raggedright
UACI \RL{(\%)}
\end{minipage} & \begin{minipage}[b]{\linewidth}\raggedright
\RL{فاصله} \RL{از} NPCR \RL{ایده}‌\RL{آل}
\end{minipage} & \begin{minipage}[b]{\linewidth}\raggedright
\RL{فاصله} \RL{از} UACI \RL{ایده}‌\RL{آل}
\end{minipage} \\
\midrule\noalign{}
\endhead
\bottomrule\noalign{}
\endlastfoot
Airplane & 512×512 & 98.5639 & 26.2476 & -1.0361 & -7.2124 \\
Baboon & 512×512 & 83.6183 & 30.2327 & -15.9817 & -3.2273 \\
Peppers & 512×512 & 99.4788 & 32.8813 & -0.1212 & -0.5787 \\
Tree & 512×512 & 98.8773 & 31.1971 & -0.7227 & -2.2629 \\
\end{longtable}
}

\includegraphics[width=6.04167in,height=2.91667in]{media/image10.png}

\begin{enumerate}
\def\labelenumi{\arabic{enumi}.}
\item
  \protect\phantomsection\label{_Toc232899813}{}\RL{نمودار} NPCR \RL{و} UACI \RL{برای} \RL{چهار} \RL{تصویر} \RL{در} \RL{مقایسه} \RL{با} \RL{مقادیر} \RL{ایده}‌\RL{آل}
\end{enumerate}

نتایج نشان می‌دهند که NPCR در اکثر تصاویر به مقادیر بالای \lr{۹۸٪} رسیده و برای تصویر Peppers به ‎\lr{۹۹.۴۸٪}‎ نزدیک شده که به آستانه ایده‌آل بسیار نزدیک است. این نتایج بیانگر آن است که تغییر جزئی در کلید (در سطح دقت ماشینی)، بیش از ‎\lr{۸۲٪}‎ (و در سایر تصاویر بیش از ‎\lr{۹۸٪}‎) از پیکسل‌های تصویر رمزنگاری‌شده را تغییر می‌دهد.

\RL{در مورد UACI، مقادیر تصویر Airplane (‎\lr{۲۶.۳۴٪}‎) فاصله معناداری از مقدار ایده‌آل (‎\lr{۳۳.۴۶٪}‎) دارد. این انحراف به ویژگی آماری تصویر Airplane مرتبط است که محتوای یکنواخت‌تری دارد و دنباله آشوبی انتخاب‌شده تغییرات شدیدتری نمی‌تواند روی آن ایجاد کند.}

\RL{تصاویر با بافت پیچیده‌تر (Baboon \lr{۳۰.۲۰٪}، Tree: \lr{۳۱.۳۱٪}) مقادیر نزدیک‌تری به ایده‌آل نشان می‌دهند.}

\begin{enumerate}
\def\labelenumi{\arabic{enumi}.}
\setcounter{enumi}{7}
\item
  \protect\phantomsection\label{_Toc232899920}{}\RL{تحلیل} \RL{زمان} \RL{اجرا}
\end{enumerate}

\RL{کارایی} \RL{محاسباتی} \RL{الگوریتم} \RL{یکی} \RL{از} \RL{معیارهای} \RL{عملیاتی} \RL{مهم} \RL{است. جدول} \RL{۴}-\RL{۵} \RL{زمان} \RL{رمزنگاری} \RL{و} \RL{رمزگشایی} \RL{را} \RL{برای} \RL{ابعاد} \RL{مختلف} \RL{تصویر} \RL{نشان} \RL{می}‌\RL{دهد}:

\begin{enumerate}
\def\labelenumi{\arabic{enumi}.}
\item
  \protect\phantomsection\label{_Toc232899952}{}\RL{زمان} \RL{اجرای} \RL{رمزنگاری} \RL{و} \RL{رمزگشایی}
\end{enumerate}

{\def\LTcaptype{none} % do not increment counter
\begin{longtable}[]{@{}
  >{\raggedright\arraybackslash}p{(\linewidth - 10\tabcolsep) * \real{0.1545}}
  >{\raggedright\arraybackslash}p{(\linewidth - 10\tabcolsep) * \real{0.1777}}
  >{\raggedright\arraybackslash}p{(\linewidth - 10\tabcolsep) * \real{0.1752}}
  >{\raggedright\arraybackslash}p{(\linewidth - 10\tabcolsep) * \real{0.1151}}
  >{\raggedright\arraybackslash}p{(\linewidth - 10\tabcolsep) * \real{0.1220}}
  >{\raggedright\arraybackslash}p{(\linewidth - 10\tabcolsep) * \real{0.1186}}@{}}
\toprule\noalign{}
\begin{minipage}[b]{\linewidth}\centering
\RL{سرعت} \textbf{(KB/s)}
\end{minipage} & \begin{minipage}[b]{\linewidth}\centering
\RL{زمان رمزگشایی} \textbf{(s)}
\end{minipage} & \begin{minipage}[b]{\linewidth}\centering
\RL{زمان رمزنگاری} \textbf{(s)}
\end{minipage} & \begin{minipage}[b]{\linewidth}\centering
\RL{پیکسل کل}
\end{minipage} & \begin{minipage}[b]{\linewidth}\centering
\RL{ابعاد}
\end{minipage} & \begin{minipage}[b]{\linewidth}\centering
\RL{تصویر}
\end{minipage} \\
\midrule\noalign{}
\endhead
\bottomrule\noalign{}
\endlastfoot
95.2 & 4.470 & 24.089 & 786,432 & 512×512 & Airplane \\
95.6 & 4.655 & 23.951 & 786,432 & 512×512 & Baboon \\
103.5 & 4.319 & 22.744 & 786,432 & 512×512 & Peppers \\
96.0 & 4.676 & 24.461 & 786,432 & 512×512 & Tree \\
\end{longtable}
}

\includegraphics[width=5.06045in,height=2.91949in]{media/image11.png}

\begin{enumerate}
\def\labelenumi{\arabic{enumi}.}
\item
  \protect\phantomsection\label{_Toc232899814}{}\RL{نمودار} \RL{مقایسه} \RL{زمان} \RL{رمزنگاری} \RL{و} \RL{رمزگشایی}
\end{enumerate}

\RL{همان}‌\RL{طور} \RL{که} \RL{مشاهده} \RL{می}‌\RL{شود،} \RL{زمان} \RL{رمزنگاری} \RL{حدود} \RL{۵}.\RL{۶} \RL{برابر} \RL{بیشتر} \RL{از} \RL{زمان} \RL{رمزگشایی} \RL{است. این} \RL{اختلاف} \RL{به} \RL{دلیل} \RL{پیچیدگی} \RL{بیشتر} \RL{محاسبات} \RL{مرحله} \RL{۳} \RL{(انتخاب} \RL{پویای} \RL{سیستم} \RL{نمایی} \RL{و} \RL{اجرای} \RL{مستقل} \RL{آن} \RL{برای} \RL{هر} \RL{پیکسل)} \RL{در} \RL{فاز} \RL{رمزنگاری} \RL{است. در} \RL{فاز} \RL{رمزگشایی،} \RL{دنباله}‌\RL{های} \RL{آشوبی} \RL{ابتدا} \RL{به}‌\RL{صورت} \RL{کامل} \RL{تولید} \RL{شده} \RL{و} \RL{سپس} \RL{یکجا} \RL{اعمال} \RL{می}‌\RL{شوند} \RL{که} \RL{کارایی} \RL{بیشتری} \RL{دارد}.

\RL{سرعت} \RL{پردازش} \RL{حدود} \RL{۷۱} \RL{کیلوبایت} \RL{در} \RL{ثانیه} \RL{است} \RL{که} \RL{در} \RL{مقایسه} \RL{با} \RL{برخی} \RL{روش}‌\RL{های} \RL{مبتنی} \RL{بر} \RL{سیستم} \RL{آشوبی} \RL{پیچیده} \RL{در} \RL{ادبیات} \RL{قابل} \RL{قبول} \RL{است،} \RL{اما} \RL{پتانسیل} \RL{بهبود} \RL{قابل} \RL{توجهی} \RL{از} \RL{طریق} \RL{برداری}‌\RL{سازی} \RL{کامل} (vectorization) \RL{با} NumPy \RL{وجود} \RL{دارد. در} \RL{پیاده}‌\RL{سازی} \RL{حاضر،} \RL{حلقه} Python \RL{برای} \RL{هر} \RL{پیکسل} \RL{به} \RL{دلیل} \RL{وابستگی} \RL{پویای} \RL{انتخاب} \RL{سیستم،} \RL{هنوز} \RL{استفاده} \RL{می}‌\RL{شود} \RL{که} \RL{عامل} \RL{اصلی} \RL{کاهش} \RL{سرعت} \RL{است}.

\begin{enumerate}
\def\labelenumi{\arabic{enumi}.}
\setcounter{enumi}{8}
\item
  \protect\phantomsection\label{_Toc232899921}{}\RL{تأیید} \RL{معکوس}‌\RL{پذیری} \RL{کامل}
\end{enumerate}

\RL{یکی} \RL{از} \RL{الزامات} \RL{اساسی} \RL{هر} \RL{الگوریتم} \RL{رمزنگاری،} \RL{توانایی} \RL{بازیابی} \RL{دقیق} \RL{تصویر} \RL{اصلی} \RL{پس} \RL{از} \RL{رمزگشایی} \RL{است. جدول} \RL{۴}-\RL{۶} \RL{نتایج} \RL{مقایسه} \RL{تصویر} \RL{اصلی} \RL{و} \RL{تصویر} \RL{بازیابی}‌\RL{شده} \RL{را} \RL{با} \RL{دو} \RL{معیار} MSE \RL{و} PSNR \RL{نشان} \RL{می}‌\RL{دهد}: \RL{}

MSE = 0 \RL{برای} \RL{تمام} \RL{تصاویر} \RL{تأیید} \RL{می}‌\RL{کند} \RL{که} \RL{هیچ} \RL{پیکسلی} \RL{در} \RL{فرآیند} \RL{رمزنگاری} \RL{و} \RL{رمزگشایی} \RL{تغییر} \RL{نکرده} \RL{است. این} \RL{نتیجه} \RL{به} \RL{دلیل} \RL{دو} \RL{ویژگی} \RL{کلیدی} \RL{طراحی} \RL{الگوریتم} \RL{حاصل} \RL{شده} \RL{است}: \RL{(الف)} \RL{دنباله}‌\RL{های} \RL{آشوبی} \RL{مستقل} \RL{از} \RL{مقادیر} \RL{پیکسل} \RL{هستند،} \RL{و} \RL{(ب)} \RL{جایگشت} \RL{معکوس} \RL{از} \RL{همان} \RL{کلید} \RL{بازسازی} \RL{می}‌\RL{شود}.

\begin{enumerate}
\def\labelenumi{\arabic{enumi}.}
\item
  \protect\phantomsection\label{_Toc232899953}{}\RL{نتایج} \RL{تأیید} \RL{معکوس}‌\RL{پذیری} \RL{الگوریتم}
\end{enumerate}

{\def\LTcaptype{none} % do not increment counter
\begin{longtable}[]{@{}
  >{\raggedright\arraybackslash}p{(\linewidth - 6\tabcolsep) * \real{0.1784}}
  >{\raggedright\arraybackslash}p{(\linewidth - 6\tabcolsep) * \real{0.2120}}
  >{\raggedright\arraybackslash}p{(\linewidth - 6\tabcolsep) * \real{0.4520}}
  >{\raggedright\arraybackslash}p{(\linewidth - 6\tabcolsep) * \real{0.1576}}@{}}
\toprule\noalign{}
\begin{minipage}[b]{\linewidth}\centering
\RL{تطابق}
\end{minipage} & \begin{minipage}[b]{\linewidth}\centering
\textbf{PSNR (dB)}
\end{minipage} & \begin{minipage}[b]{\linewidth}\centering
\textbf{MSE} \RL{(اصلی در برابر بازیابی)}
\end{minipage} & \begin{minipage}[b]{\linewidth}\centering
\RL{تصویر}
\end{minipage} \\
\midrule\noalign{}
\endhead
\bottomrule\noalign{}
\endlastfoot
\RL{بله} & $\infty$ & 0.000000 & Airplane \\
\RL{بله} & $\infty$ & 0.000000 & Baboon \\
\RL{بله} & $\infty$ & 0.000000 & Peppers \\
\RL{بله} & $\infty$ & 0.000000 & Tree \\
\end{longtable}
}

\begin{enumerate}
\def\labelenumi{\arabic{enumi}.}
\setcounter{enumi}{9}
\item
  \protect\phantomsection\label{_Toc232899922}{}\RL{مقایسه} \RL{با} \RL{روش}‌\RL{های} \RL{مرجع} \RL{در} \RL{ادبیات}
\end{enumerate}

\RL{جدول} \RL{۴}-\RL{۷} \RL{عملکرد} \RL{الگوریتم} \RL{پیشنهادی} \RL{را} \RL{با} \RL{چند} \RL{روش} \RL{مرجع} \RL{مشابه} \RL{در} \RL{ادبیات} \RL{مقایسه} \RL{می}‌\RL{کند. داده}‌\RL{های} \RL{مقایسه}‌\RL{ای} \RL{از} \RL{مقالات} \RL{ذکرشده} \RL{در} \RL{فصل} \RL{دوم} \RL{استخراج} \RL{شده}‌\RL{اند}:

\begin{enumerate}
\def\labelenumi{\arabic{enumi}.}
\item
  \protect\phantomsection\label{_Toc232899954}{}\RL{مقایسه} \RL{الگوریتم} \RL{پیشنهادی} \RL{با} \RL{روش}‌\RL{های} \RL{مرجع}
\end{enumerate}

{\def\LTcaptype{none} % do not increment counter
\begin{longtable}[]{@{}
  >{\raggedright\arraybackslash}p{(\linewidth - 10\tabcolsep) * \real{0.1931}}
  >{\raggedright\arraybackslash}p{(\linewidth - 10\tabcolsep) * \real{0.1427}}
  >{\raggedright\arraybackslash}p{(\linewidth - 10\tabcolsep) * \real{0.1476}}
  >{\raggedright\arraybackslash}p{(\linewidth - 10\tabcolsep) * \real{0.1012}}
  >{\raggedright\arraybackslash}p{(\linewidth - 10\tabcolsep) * \real{0.1992}}
  >{\raggedright\arraybackslash}p{(\linewidth - 10\tabcolsep) * \real{0.2162}}@{}}
\toprule\noalign{}
\begin{minipage}[b]{\linewidth}\centering
\RL{سیستم آشوبی}
\end{minipage} & \begin{minipage}[b]{\linewidth}\centering
\textbf{UACI (\%)}
\end{minipage} & \begin{minipage}[b]{\linewidth}\centering
\textbf{NPCR (\%)}
\end{minipage} & \begin{minipage}[b]{\linewidth}\centering
\RL{آنتروپی}
\end{minipage} & \begin{minipage}[b]{\linewidth}\centering
\RL{تصویر مرجع}
\end{minipage} & \begin{minipage}[b]{\linewidth}\centering
\RL{روش}
\end{minipage} \\
\midrule\noalign{}
\endhead
\bottomrule\noalign{}
\endlastfoot
Chen + 9 ECM & 30.14 & 95.13 & 8.0 & USC-SIPI & \RL{روش پیشنهادی} \\
Lorenz & 30.35 & 99.65 & 7.9974 & \RL{مجموعه عمومی} & \RL{آلوارز و همکاران{[}6{]}} \\
Chen+Logistic & 33.21 & 99.58 & 7.9971 & USC-SIPI & \RL{خدادادی و زندوکیلی{[}7{]}} \\
Logistic & 33.38 & 99.57 & 7.9968 & Lena 256×256 & \RL{وانگ و لی{[}8{]}} \\
Hyper-chaos & 33.46 & 99.61 & 7.9990 & \RL{مجموعه عمومی} & \RL{مقاله فوق‌آشوب{[}12{]}} \\
\end{longtable}
}

\RL{مقایسه نتایج نشان می‌دهد که الگوریتم پیشنهادی از نظر آنتروپی با روش‌های مرجع رقابت‌پذیر است. NPCR تا ‎\lr{۹۹.۴۸٪}‎ برای تصویر Peppers و میانگین ‎\lr{۹۴.۹۲٪} در محدوده عملکرد روش‌های مرجع قرار دارد. اما UACI نسبت به برخی روش‌ها (مخصوصاً برای تصویر Airplane) فاصله بیشتری از مقدار ایده‌آل دارد که نشان‌دهنده فرصت بهبود در مرحله انتشار است.}

\RL{مزیت} \RL{اصلی} \RL{الگوریتم} \RL{پیشنهادی} \RL{نسبت} \RL{به} \RL{روش}‌\RL{های} \RL{مرجع،} \RL{تنوع} \RL{۹گانه} \RL{سیستم}‌\RL{های} \RL{آشوبی} \RL{و} \RL{انتخاب} \RL{پویای} \RL{آن}‌\RL{ها} \RL{برای} \RL{هر} \RL{پیکسل} \RL{است. این} \RL{ویژگی} \RL{حتی} \RL{اگر} \RL{مهاجم} \RL{موفق} \RL{به} \RL{شناسایی} \RL{یک} \RL{سیستم} \RL{نمایی} \RL{شود،} \RL{اطلاعات} \RL{حاصل} \RL{را} \RL{برای} \RL{پیکسل}‌\RL{های} \RL{دیگر} \RL{بی}‌\RL{فایده} \RL{می}‌\RL{کند، یک} \RL{لایه} \RL{امنیتی} \RL{که} \RL{در} \RL{روش}‌\RL{های} \RL{مرجع} \RL{وجود} \RL{ندارد}.

\begin{enumerate}
\def\labelenumi{\arabic{enumi}.}
\setcounter{enumi}{10}
\item
  \protect\phantomsection\label{_Toc232899923}{}\RL{جمع}‌\RL{بندی} \RL{فصل}
\end{enumerate}

\RL{نتایج} \RL{تجربی} \RL{ارائه}‌\RL{شده} \RL{در} \RL{این} \RL{فصل،} \RL{اثربخشی} \RL{الگوریتم} \RL{رمزنگاری} \RL{پیشنهادی} \RL{را} \RL{از} \RL{جنبه}‌\RL{های} \RL{مختلف} \RL{تأیید} \RL{می}‌\RL{کنند. مهم}‌\RL{ترین} \RL{یافته}‌\RL{ها} \RL{عبارتند} \RL{از}:

\begin{itemize}
\item
  \RL{معکوس}‌\RL{پذیری} \RL{\lr{۱۰۰٪}}: \RL{}MSE = 0 \RL{و} PSNR = $\infty$ \RL{برای} \RL{تمام} \RL{چهار} \RL{تصویر} --- \RL{الگوریتم} \RL{هیچ} \RL{داده}‌\RL{ای} \RL{از} \RL{دست} \RL{نمی}‌\RL{دهد}
\item
  آنتروپی بالا: میانگین ۷.۹۹۹۳ بیت (از ۸ بیت ایده‌آل) --- توزیع پیکسل‌ها تقریباً کاملاً یکنواخت است
\item
  حذف همبستگی: کاهش میانگین از ‎۰.۹۹۸۹‎ به ‎۰.۰۱۲۶ --- الگوی مکانی تصویر کاملاً از بین رفته است
\item
  حساسیت به کلید (NPCR) : حداکثر \lr{۹۹.۴۸٪} --- تغییر ۱۰⁻¹⁵ در کلید، بیش از ‎\lr{۸۲٪}‎ پیکسل‌ها را تغییر می‌دهد
\item
  \RL{هیستوگرام} \RL{یکنواخت}: \RL{در} \RL{تمام} \RL{کانال}‌\RL{های}R,G,B \RL{} \RL{مقاومت} \RL{در} \RL{برابر} \RL{حملات} \RL{آماری} \RL{مبتنی} \RL{بر} \RL{هیستوگرام}
\item
  \RL{فضای} \RL{کلید} \RL{\textgreater{}} 2\^{}348 \RL{: حمله} \RL{جستجوی} \RL{فراگیر} \RL{به}‌\RL{لحاظ} \RL{محاسباتی} \RL{ناممکن} \RL{است}
\end{itemize}

\RL{حوزه}‌\RL{های} \RL{بهبود} \RL{شناسایی}‌\RL{شده:} \RL{(الف)} \RL{بهینه}‌\RL{سازی} \RL{زمان} \RL{اجرا} \RL{از} \RL{طریق} \RL{برداری}‌\RL{سازی} \RL{با} NumPy \RL{یا} \RL{پیاده}‌\RL{سازی} CUDA\RL{؛} \RL{(ب)} \RL{بهبود} UACI \RL{از} \RL{طریق} \RL{افزودن} \RL{یک} \RL{مرحله} feedback \RL{به} \RL{فاز} \RL{انتشار؛} \RL{(ج)} \RL{ارزیابی} \RL{روی} \RL{مجموعه} \RL{داده}‌\RL{های} \RL{بزرگ}‌\RL{تر} \RL{مانند} ImageNet \RL{برای} \RL{تعمیم}‌\RL{پذیری} \RL{بهتر.}

\begin{enumerate}
\def\labelenumi{\arabic{enumi}.}
\setcounter{enumi}{4}
\item
  \protect\phantomsection\label{_Toc232899924}{}\RL{بحث و نتیجه‌گیری}
\end{enumerate}

\RL{\hfill\break
}

\begin{enumerate}
\def\labelenumi{\arabic{enumi}.}
\item
  \protect\phantomsection\label{_Toc232899925}{}\RL{مقدمه}
\end{enumerate}

\RL{این} \RL{فصل،} \RL{حلقه}‌\RL{ای} \RL{پایانی} \RL{پژوهش} \RL{حاضر} \RL{را} \RL{تشکیل} \RL{می}‌\RL{دهد} \RL{و} \RL{با} \RL{نگاهی} \RL{جامع} \RL{به} \RL{مسیر} \RL{طی}‌\RL{شده،} \RL{از} \RL{طرح} \RL{مسئله} \RL{تا} \RL{ارزیابی} \RL{تجربی،} \RL{دستاوردها} \RL{و} \RL{یافته}‌\RL{های} \RL{اصلی} \RL{را} \RL{در} \RL{قالبی} \RL{منسجم} \RL{ارائه} \RL{می}‌\RL{کند}. \RL{ابتدا} \RL{مروری} \RL{اجمالی} \RL{بر} \RL{اهداف} \RL{و} \RL{سؤال}‌\RL{های} \RL{اصلی} \RL{پژوهش} \RL{صورت} \RL{می}‌\RL{گیرد،} \RL{سپس} \RL{میزان} \RL{تحقق} \RL{هر} \RL{یک} \RL{از} \RL{اهداف} \RL{بر} \RL{پایه} \RL{نتایج} \RL{عددی} \RL{فصل} \RL{چهارم} \RL{ارزیابی} \RL{می}‌\RL{شود}. \RL{در} \RL{ادامه،} \RL{محدودیت}‌\RL{های} \RL{شناسایی}‌\RL{شده} \RL{در} \RL{این} \RL{پژوهش} \RL{مطرح} \RL{و} \RL{پیشنهادات} \RL{مشخصی} \RL{برای} \RL{تحقیقات} \RL{آینده} \RL{ارائه} \RL{می}‌\RL{گردد}.

\RL{پژوهش} \RL{حاضر} \RL{در} \RL{پاسخ} \RL{به} \RL{یک} \RL{نیاز} \RL{واقعی} \RL{و} \RL{رو} \RL{به} \RL{رشد} \RL{در} \RL{حوزه} \RL{امنیت} \RL{اطلاعات} \RL{تصویری} \RL{شکل} \RL{گرفت}: \RL{طراحی} \RL{الگوریتمی} \RL{که} \RL{نه} \RL{تنها} \RL{امنیت} \RL{بالایی} \RL{در} \RL{برابر} \RL{انواع} \RL{حملات} \RL{شناخته}‌\RL{شده} \RL{داشته} \RL{باشد،} \RL{بلکه} \RL{از} \RL{پایه} \RL{نظری} \RL{محکم} \RL{مبتنی} \RL{بر} \RL{سیستم}‌\RL{های} \RL{آشوبی} \RL{نمایی} \RL{نسل} \RL{جدید} \RL{برخوردار} \RL{باشد}. \RL{الگوریتم} \RL{پیشنهادی،} \RL{با} \RL{ترکیب} \RL{سیستم} \RL{سه}‌\RL{بعدی} Chen \RL{برای} \RL{جایگشت} \RL{و} \RL{۹} \RL{سیستم} \RL{آشوبی} \RL{نمایی} \RL{برای} \RL{انتشار} \RL{پویا،} \RL{رویکردی} \RL{چندلایه} \RL{و} \RL{نوآورانه} \RL{به} \RL{رمزنگاری} \RL{تصاویر} \RL{رنگی} \RL{ارائه} \RL{می}‌\RL{دهد}.

\begin{enumerate}
\def\labelenumi{\arabic{enumi}.}
\setcounter{enumi}{1}
\item
  \protect\phantomsection\label{_Toc232899926}{}\RL{مرور} \RL{اهداف} \RL{پژوهش}
\end{enumerate}

\RL{در} \RL{این} \RL{بخش،} \RL{هر} \RL{یک} \RL{از} \RL{اهداف} \RL{کلی} \RL{و} \RL{جزئی} \RL{تعریف}‌\RL{شده} \RL{در} \RL{فصل} \RL{اول} \RL{با} \RL{نتایج} \RL{حاصل} \RL{از} \RL{آزمایش}‌\RL{های} \RL{فصل} \RL{چهارم} \RL{مقایسه} \RL{و} \RL{میزان} \RL{تحقق} \RL{آن}‌\RL{ها} \RL{ارزیابی} \RL{می}‌\RL{شود}.

\begin{enumerate}
\def\labelenumi{\arabic{enumi}.}
\item
  \protect\phantomsection\label{_Toc232899927}{}\RL{هدف} \RL{کلی: طراحی} \RL{و} \RL{توسعه} \RL{الگوریتم} \RL{رمزنگاری}
\end{enumerate}

\RL{هدف} \RL{کلی} \RL{پژوهش،} \RL{طراحی} \RL{و} \RL{توسعه}‌\RL{ی} \RL{الگوریتمی} \RL{برای} \RL{رمزنگاری} \RL{تصاویر} \RL{رنگی} \RL{با} \RL{بهره}‌\RL{گیری} \RL{از} \RL{ترکیب} \RL{۹} \RL{سیستم} \RL{آشوبی} \RL{نمایی} \RL{بود. این} \RL{هدف} \RL{به}‌\RL{طور} \RL{کامل} \RL{محقق} \RL{شده} \RL{است}: \RL{الگوریتمی} \RL{پنج}‌\RL{مرحله}‌\RL{ای} \RL{طراحی،} \RL{پیاده}‌\RL{سازی} \RL{و} \RL{ارزیابی} \RL{شد} \RL{که} \RL{شامل} \RL{جایگشت} \RL{مبتنی} \RL{بر} Chen\RL{،} \RL{انتخاب} \RL{پویای} \RL{سیستم} \RL{نمایی} \RL{برای} \RL{هر} \RL{پیکسل،} \RL{دو} \RL{لایه} XOR \RL{مستقل} \RL{و} \RL{معکوس}‌\RL{پذیری} \RL{کامل} \RL{می}‌\RL{باشد}. \RL{پیاده}‌\RL{سازی} \RL{کامل} \RL{در} Python 3.10 \RL{با} \RL{استفاده} \RL{از} \RL{کتابخانه}‌\RL{های} \RL{علمی} \RL{و} \RL{اجرای} \RL{موفق} \RL{بر} \RL{روی} \RL{چهار} \RL{تصویر} \RL{استاندارد} USC-SIPI\RL{،} \RL{تحقق} \RL{این} \RL{هدف} \RL{را} \RL{تأیید} \RL{می}‌\RL{کند}.

\begin{enumerate}
\def\labelenumi{\arabic{enumi}.}
\setcounter{enumi}{1}
\item
  \protect\phantomsection\label{_Toc232899928}{}\RL{اهداف} \RL{جزئی} \RL{و} \RL{میزان} \RL{دستیابی}

  \begin{enumerate}
  \def\labelenumii{\arabic{enumii}.}
  \item
    \protect\phantomsection\label{_Toc232899955}{}\RL{ارزیابی} \RL{میزان} \RL{تحقق} \RL{اهداف} \RL{جزئی} \RL{پژوهش}
  \end{enumerate}
\end{enumerate}

{\def\LTcaptype{none} % do not increment counter
\begin{longtable}[]{@{}
  >{\raggedright\arraybackslash}p{(\linewidth - 4\tabcolsep) * \real{0.3872}}
  >{\raggedright\arraybackslash}p{(\linewidth - 4\tabcolsep) * \real{0.2846}}
  >{\raggedright\arraybackslash}p{(\linewidth - 4\tabcolsep) * \real{0.3282}}@{}}
\toprule\noalign{}
\endhead
\bottomrule\noalign{}
\endlastfoot
\textbf{\RL{نتیجه} \RL{حاصل}} & \textbf{\RL{معیار} \RL{ارزیابی}} & \textbf{\RL{هدف} \RL{جزئی}} \\
میانگین ۸.۰ از ۸ بیت ایده‌آل & \RL{آنتروپی} \RL{شانون} \RL{کانال}‌\RL{های} \RL{رمزشده} & \RL{بررسی} \RL{عملکرد} \RL{۹} \RL{سیستم} \RL{نمایی} \RL{در} \RL{تولید} \RL{دنباله} \RL{تصادفی} \\
\RL{هیستوگرام} \RL{مسطح} \RL{در} \RL{تمام} \RL{کانال}‌\RL{ها؛} \RL{فضای} \RL{کلید} \textasciitilde10105 & \RL{هیستوگرام} \RL{یکنواخت،} \RL{فضای} \RL{کلید} \textgreater{} 2348 & \RL{بهبود} \RL{مقاومت} \RL{در} \RL{برابر} \RL{حملات} \RL{آماری} \RL{و} \RL{جستجوی} \RL{فراگیر} \\
NPCR تا ‎\lr{۹۹.۴۸٪}‎؛ همبستگی کاهش از ‎۰.۹۹۸۹‎ به ‎۰.۰۱۲۶ & NPCR \RL{و} \RL{ضریب} \RL{همبستگی} \RL{رمزشده} & \RL{افزایش} \RL{پیچیدگی} \RL{با} \RL{انتخاب} \RL{تصادفی} \RL{سیستم} \RL{برای} \RL{هر} \RL{پیکسل} \\
\RL{آنتروپی} \RL{رقابت}‌\RL{پذیر؛} NPCR \RL{در} \RL{محدوده} \RL{مرجع؛} UACI \RL{نیاز} \RL{به} \RL{بهبود} \RL{دارد} & \RL{آنتروپی،} NPCR\RL{،} UACI \RL{در} \RL{مقایسه} \RL{با} \RL{{[}۶،۷،۸،۱۲{]}} & \RL{مقایسه} \RL{با} \RL{روش}‌\RL{های} \RL{مرجع} \RL{در} \RL{ادبیات} \\
\end{longtable}
}

\begin{enumerate}
\def\labelenumi{\arabic{enumi}.}
\setcounter{enumi}{2}
\item
  \protect\phantomsection\label{_Toc232899929}{}\RL{یافته}‌\RL{های} \RL{اصلی} \RL{پژوهش}
\end{enumerate}

\RL{بر} \RL{اساس} \RL{نتایج} \RL{تجربی} \RL{ارائه}‌\RL{شده} \RL{در} \RL{فصل} \RL{چهارم،} \RL{یافته}‌\RL{های} \RL{اصلی} \RL{این} \RL{پژوهش} \RL{را} \RL{می}‌\RL{توان} \RL{در} \RL{پنج} \RL{محور} \RL{اصلی} \RL{دسته}‌\RL{بندی} \RL{کرد}:

\begin{enumerate}
\def\labelenumi{\arabic{enumi}.}
\item
  \protect\phantomsection\label{_Toc232899930}{}\RL{معکوس}‌\RL{پذیری} \RL{کامل} \RL{الگوریتم}
\end{enumerate}

\RL{مهم}‌\RL{ترین} \RL{یافته} \RL{عملیاتی} \RL{پژوهش،} \RL{تأیید} \RL{معکوس}‌\RL{پذیری} \RL{\lr{۱۰۰٪}} \RL{الگوریتم} \RL{پیشنهادی} \RL{است}. \RL{}

\RL{مقدار} MSE = 0 \RL{و} PSNR = inf \RL{برای} \RL{تمام} \RL{چهار} \RL{تصویر} \RL{آزمایش}‌\RL{شده} \RL{نشان} \RL{می}‌\RL{دهد} \RL{که} \RL{هیچ}‌\RL{گونه} \RL{اطلاعاتی} \RL{در} \RL{فرآیند} \RL{رمزنگاری} \RL{از} \RL{دست} \RL{نمی}‌\RL{رود}. \RL{این} \RL{ویژگی} \RL{حیاتی،} \RL{کاربرد} \RL{الگوریتم} \RL{را} \RL{در} \RL{سامانه}‌\RL{هایی} \RL{نظیر} \RL{ذخیره}‌\RL{سازی} \RL{تصاویر} \RL{پزشکی} \RL{یا} \RL{انتقال} \RL{داده}‌\RL{های} \RL{طبقه}‌\RL{بندی}‌\RL{شده} \RL{که} \RL{در} \RL{آن}‌\RL{ها} \RL{هرگونه} \RL{تخریب} \RL{اطلاعات} \RL{غیرقابل} \RL{قبول} \RL{است} \RL{ممکن} \RL{می}‌\RL{سازد}.

\RL{تضمین} \RL{معکوس}‌\RL{پذیری} \RL{به} \RL{دو} \RL{ویژگی} \RL{طراحی} \RL{بازمی}‌\RL{گردد:}

\RL{الف)} \RL{وابسته} \RL{نبودن} \RL{دنباله}‌\RL{های} \RL{آشوبی} \RL{به} \RL{مقادیر} \RL{پیکسل}‌\RL{های} \RL{تصویر؛} \RL{و} \RL{ب)} \RL{طراحی} \RL{کامل} \RL{جایگشت} \RL{معکوس} \RL{از} \RL{روی} \RL{همان} \RL{کلید. این} \RL{رویکرد،} \RL{برخلاف} \RL{برخی} \RL{روش}‌\RL{های} \RL{مبتنی} \RL{بر} \RL{بازخورد} \RL{پیکسل} \RL{که} \RL{در} \RL{آن}‌\RL{ها} \RL{خطای} \RL{انتشاری} \RL{وجود} \RL{دارد،} \RL{معکوس}‌\RL{پذیری} \RL{قطعی} \RL{را} \RL{تضمین} \RL{می}‌\RL{کند}.

\begin{enumerate}
\def\labelenumi{\arabic{enumi}.}
\setcounter{enumi}{1}
\item
  \protect\phantomsection\label{_Toc232899931}{}\RL{آنتروپی} \RL{بالا} \RL{و} \RL{توزیع} \RL{یکنواخت} \RL{پیکسل}‌\RL{ها}
\end{enumerate}

میانگین آنتروپی شانون تصاویر رمزنگاری‌شده برابر با ۸.۰ بیت است که انحراف ۰.۰ بیت از مقدار ایده‌آل ۸ بیت دارد. این نتیجه برای تمام تصاویر، صرف نظر از محتوا و ساختار آماری اولیه، به‌دست آمده است. به‌طور خاص، تصویر Airplane که هیستوگرام اصلی بسیار نامتقارنی داشت (آنتروپی میانگین \textasciitilde۶.۷۶ بیت)، پس از رمزنگاری به ۸.۰ بیت رسید که نشان‌دهنده توانایی الگوریتم در یکسان‌سازی توزیع پیکسل‌ها مستقل از محتوای اولیه تصویر است.

\RL{این} \RL{یافته} \RL{فرضیه} \RL{سوم} \RL{پژوهش} \RL{را} \RL{تأیید} \RL{می}‌\RL{کند: انتخاب} \RL{تصادفی} \RL{سیستم} \RL{آشوبی} \RL{نمایی} \RL{برای} \RL{هر} \RL{پیکسل،} \RL{منجر} \RL{به} \RL{تولید} \RL{دنباله}‌\RL{های} \RL{رمزنگاری}‌\RL{شده} \RL{با} \RL{آنتروپی} \RL{بالاتر} \RL{می}‌\RL{شود}. \RL{تنوع} \RL{۹گانه} \RL{سیستم}‌\RL{ها،} \RL{احتمال} \RL{تکرار} \RL{الگوی} \RL{رمزنگاری} \RL{را} \RL{در} \RL{سطح} \RL{بایت} \RL{به} \RL{حداقل} \RL{می}‌\RL{رساند}.

\begin{enumerate}
\def\labelenumi{\arabic{enumi}.}
\setcounter{enumi}{2}
\item
  \protect\phantomsection\label{_Toc232899932}{}\RL{حذف} \RL{مؤثر} \RL{همبستگی} \RL{مکانی}
\end{enumerate}

ترکیب جایگشت مبتنی بر Chen و انتشار XOR دولایه، همبستگی مکانی پیکسل‌های مجاور را از میانگین ‎۰.۹۹۸۹‎ در تصاویر اصلی به ‎۰.۰۱۲۶‎ در تصاویر رمزنگاری‌شده کاهش داده است؛ یعنی کاهشی معادل ‎\lr{۹۸.۷٪}. این نتیجه در هر سه جهت افقی، عمودی و قطری و برای تمام چهار تصویر مشاهده شد و فرضیه دوم پژوهش مبنی بر اثربخشی الگوریتم جایگشت Chen در کاهش همبستگی را به‌طور قاطع تأیید می‌کند.

\RL{نمودارهای} \RL{پراکندگی} \RL{فصل} \RL{چهارم} \RL{این} \RL{واقعیت} \RL{را} \RL{به}‌\RL{صورت} \RL{بصری} \RL{آشکار} \RL{می}‌\RL{سازند}: \RL{در} \RL{حالی} \RL{که} \RL{نقاط} \RL{پراکندگی} \RL{تصاویر} \RL{اصلی} \RL{روی} \RL{خط} \RL{قطری} \RL{متمرکز} \RL{هستند،} \RL{در} \RL{تصاویر} \RL{رمزنگاری}‌\RL{شده} \RL{به}‌\RL{صورت} \RL{کاملاً} \RL{تصادفی} \RL{در} \RL{سراسر} \RL{فضا} \RL{پخش} \RL{شده}‌\RL{اند}.

\begin{enumerate}
\def\labelenumi{\arabic{enumi}.}
\setcounter{enumi}{3}
\item
  \protect\phantomsection\label{_Toc232899933}{}\RL{حساسیت} \RL{بالا} \RL{به} \RL{کلید}
\end{enumerate}

آزمون NPCR با تغییر 10\^{}15 در یکی از پارامترهای کلید نشان داد که بیش از ‎\lr{۸۲٪}‎ پیکسل‌های تصویر رمزنگاری‌شده تغییر می‌کنند (حداکثر ‎\lr{۹۹.۴۸٪}‎ برای Peppers). این تاثیر تضمین می‌کند که حتی اگر مهاجم به کلید تقریبی دسترسی داشته باشد، رمزگشایی صحیح ناممکن خواهد بود.

ارزیابی UACI نشان داد که میانگین شدت تغییرات بین ‎\lr{۲۶.۳۴٪}‎ (Airplane) تا ‎\lr{۳۲.۹۱٪}‎ (Peppers) متغیر است (با میانگین کل ‎\lr{۳۰.۱۹٪}‎). تصاویر با بافت پیچیده‌تر نتایج بهتری نشان می‌دهند که به ساختار آماری تصویر اصلی و توزیع سیستم‌های نمایی انتخاب‌شده وابسته است.

\begin{enumerate}
\def\labelenumi{\arabic{enumi}.}
\setcounter{enumi}{4}
\item
  \protect\phantomsection\label{_Toc232899934}{}\RL{فضای} \RL{کلید} \RL{بزرگ} \RL{و} \RL{مقاومت} \RL{در} \RL{برابر} \RL{جستجوی} \RL{فراگیر}
\end{enumerate}

\RL{با} \RL{هفت} \RL{پارامتر} \RL{پیوسته} \RL{با} \RL{دقت} \RL{۱۵} \RL{رقم} \RL{اعشاری،} \RL{فضای} \RL{کلید} \RL{الگوریتم} \RL{از} 10\^{}105 \RL{تجاوز} \RL{می}‌\RL{کند} \RL{که} \RL{معادل} \RL{بیش} \RL{از} 2\^{}348 \RL{است.} \RL{این} \RL{مقدار} \RL{از} \RL{آستانه} \RL{امنیتی} AES-256 \RL{(فضای} \RL{کلید} 2\^{}256\RL{)} \RL{بسیار} \RL{فراتر} \RL{می}‌\RL{رود} \RL{و} \RL{حمله} \RL{جستجوی} \RL{فراگیر} \RL{را} \RL{با} \RL{هر} \RL{سخت}‌\RL{افزاری} \RL{به} \RL{لحاظ} \RL{محاسباتی} \RL{ناممکن} \RL{می}‌\RL{سازد.}

\begin{enumerate}
\def\labelenumi{\arabic{enumi}.}
\setcounter{enumi}{3}
\item
  \protect\phantomsection\label{_Toc232899935}{}\RL{ارزیابی} \RL{فرضیه}‌\RL{های} \RL{پژوهش}
\end{enumerate}

\RL{در} \RL{این} \RL{بخش،} \RL{هر} \RL{یک} \RL{از} \RL{فرضیه}‌\RL{های} \RL{مطرح}‌\RL{شده} \RL{در} \RL{فصل} \RL{اول} \RL{با} \RL{شواهد} \RL{عددی} \RL{حاصل} \RL{از} \RL{آزمایش}‌\RL{ها} \RL{بررسی} \RL{می}‌\RL{شود}:

\begin{enumerate}
\def\labelenumi{\arabic{enumi}.}
\item
  \protect\phantomsection\label{_Toc232899956}{}\RL{ارزیابی} \RL{فرضیه}‌\RL{های} \RL{پژوهش} \RL{بر} \RL{اساس} \RL{نتایج} \RL{تجربی}
\end{enumerate}

{\def\LTcaptype{none} % do not increment counter
\begin{longtable}[]{@{}
  >{\centering\arraybackslash}p{(\linewidth - 4\tabcolsep) * \real{0.1772}}
  >{\centering\arraybackslash}p{(\linewidth - 4\tabcolsep) * \real{0.2906}}
  >{\centering\arraybackslash}p{(\linewidth - 4\tabcolsep) * \real{0.5322}}@{}}
\toprule\noalign{}
\endhead
\bottomrule\noalign{}
\endlastfoot
\RL{نتیجه} & \RL{شاهد} \RL{عددی} & \RL{فرضیه} \\
\RL{تأیید} \RL{شد} & \RL{هیستوگرام مسطح در تمام تصاویر؛ آنتروپی = 8.0 بیت} & \RL{فرضیه} \RL{۱}: \RL{ترکیب} \RL{۹} \RL{سیستم} \RL{نمایی} \RL{منجر} \RL{به} \RL{افزایش} \RL{امنیت} \RL{در} \RL{برابر} \RL{حملات} \RL{آماری} \RL{می}‌\RL{شود} \\
\RL{تأیید} \RL{شد} & کاهش همبستگی از ‎۰.۹۹۸۹‎ به ‎۰.۰۱۲۶‎ (کاهش ‎\lr{۹۸.۷٪}) & \RL{فرضیه} \RL{۲}: \RL{جایگشت} Chen \RL{سبب} \RL{کاهش} \RL{همبستگی} \RL{بین} \RL{پیکسل}‌\RL{های} \RL{مجاور} \RL{می}‌\RL{شود} \\
\RL{تأیید} \RL{شد} & میانگین آنتروپی ‎۷.۹۹۹۳‎ بیت (‎\lr{۰.۰۰۹٪}‎ انحراف از ایده‌آل) & \RL{فرضیه} \RL{۳}: \RL{انتخاب} \RL{تصادفی} \RL{سیستم} \RL{نمایی} \RL{منجر} \RL{به} \RL{آنتروپی} \RL{بالاتر} \RL{می}‌\RL{شود} \\
\RL{نسبی،} \RL{نیاز} \RL{به} \RL{بهینه}‌\RL{سازی} & رمزنگاری ۲۲.۷-۲۴.۵ ثانیه برای 512x512؛ رمزگشایی ۴.۳-۴.۷ ثانیه & \RL{فرضیه} \RL{۴}: \RL{الگوریتم} \RL{از} \RL{نظر} \RL{زمان} \RL{اجرا} \RL{در} \RL{محدوده} \RL{قابل}‌\RL{قبول} \RL{قرار} \RL{دارد} \\
\RL{تأیید} \RL{شد} & \RL{آنتروپی} \RL{رقابت}‌\RL{پذیر} \RL{با} {]}\RL{۶،۷،۱۲} {[}\RL{؛} NPCR \RL{در} \RL{محدوده} \RL{مرجع} & \RL{فرضیه} \RL{۵}: \RL{توزیع} \RL{آماری} \RL{تصادفی}‌\RL{تر} \RL{نسبت} \RL{به} \RL{الگوریتم}‌\RL{های} \RL{پیشین} \\
\end{longtable}
}

\begin{enumerate}
\def\labelenumi{\arabic{enumi}.}
\setcounter{enumi}{4}
\item
  \protect\phantomsection\label{_Toc232899936}{}\RL{محدودیت}‌\RL{های} \RL{پژوهش}
\end{enumerate}

\RL{هر} \RL{پژوهش} \RL{علمی} \RL{با} \RL{محدودیت}‌\RL{هایی} \RL{روبه}‌\RL{روست} \RL{که} \RL{شناسایی} \RL{صادقانه} \RL{آن}‌\RL{ها،} \RL{پایه}‌\RL{ای} \RL{برای} \RL{تحقیقات} \RL{آینده} \RL{فراهم} \RL{می}‌\RL{کند}. \RL{محدودیت}‌\RL{های} \RL{اصلی} \RL{پژوهش} \RL{حاضر} \RL{عبارتند} \RL{از}:

\begin{itemize}
\item
  زمان اجرای رمزنگاری (حدود ۲۳.۸ ثانیه برای تصاویر 512x512) به دلیل پیاده‌سازی حلقه Python برای هر پیکسل مناسب کاربردهای بلادرنگ نیست. این محدودیت ذاتی پیاده‌سازی است، نه محدودیت معماری الگوریتم.
\item
  مقدار UACI برای تصویر Airplane (‎\lr{۲۶.۳۴٪}‎) اندکی از مقدار ایده‌آل (‎\lr{۳۳.۴۶٪}‎) فاصله دارد. این انحراف در تصاویر با محتوای یکنواخت ساختاری است و به مرحله انتشار نیاز به بازطراحی دارد.
\item
  \RL{آزمون}‌\RL{های} NIST SP 800-22 \RL{برای} \RL{ارزیابی} \RL{رسمی} \RL{تصادفی}‌\RL{بودن} \RL{دنباله}‌\RL{های} \RL{کلید} \RL{در} \RL{این} \RL{پژوهش} \RL{انجام} \RL{نشده}‌\RL{اند} \RL{و} \RL{برای} \RL{انتشار} \RL{رسمی} \RL{نتایج} \RL{ضروری} \RL{هستند}.
\item
  \RL{ارزیابی} \RL{مقاومت} \RL{در} \RL{برابر} \RL{حملات} \RL{پیچیده}‌\RL{تر} \RL{مانند} \RL{حمله} \RL{متن} \RL{انتخابی}\footnote{Chosen-Plaintext Attack} \RL{و} \RL{حمله} \RL{کانال} \RL{جانبی}\footnote{Side-Channel Attack} \RL{در} \RL{این} \RL{پژوهش} \RL{انجام} \RL{نشده} \RL{است}.
\item
  \RL{پیاده}‌\RL{سازی} \RL{فعلی} \RL{صرفاً} \RL{بر} Python \RL{سرعت} \RL{محور} \RL{نیست} \RL{و} \RL{مقایسه}‌\RL{های} \RL{زمانی} \RL{با} \RL{روش}‌\RL{های} \RL{مبتنی} \RL{بر} C++ \RL{یا} CUDA \RL{باید} \RL{با} \RL{احتیاط} \RL{انجام} \RL{شود}.

  \begin{enumerate}
  \def\labelenumi{\arabic{enumi}.}
  \item
    \protect\phantomsection\label{_Toc232899937}{}\RL{پیشنهادات} \RL{برای} \RL{تحقیقات} \RL{آینده}
  \end{enumerate}
\end{itemize}

\RL{بر} \RL{اساس} \RL{یافته}‌\RL{های} \RL{این} \RL{پژوهش} \RL{و} \RL{محدودیت}‌\RL{های} \RL{شناسایی}‌\RL{شده،} \RL{پیشنهادات} \RL{زیر} \RL{برای} \RL{پژوهشگران} \RL{آینده} \RL{ارائه} \RL{می}‌\RL{شود}:

\begin{enumerate}
\def\labelenumi{\arabic{enumi}.}
\item
  \protect\phantomsection\label{_Toc232899938}{}\RL{بهینه}‌\RL{سازی} \RL{کارایی} \RL{محاسباتی}
\end{enumerate}

\RL{مهم}‌\RL{ترین} \RL{حوزه} \RL{بهبود،} \RL{بهینه}‌\RL{سازی} \RL{سرعت} \RL{پردازش} \RL{است}. \RL{سه} \RL{رویکرد} \RL{اصلی} \RL{پیشنهاد} \RL{می}‌\RL{شود}:

\begin{itemize}
\item
  \RL{برداری}‌\RL{سازی} \RL{کامل}\footnote{Vectorization}\RL{:} \RL{با} \RL{بازنگری} \RL{در} \RL{ساختار} \RL{انتخاب} \RL{پویای} \RL{سیستم} \RL{نمایی} \RL{و} \RL{استفاده} \RL{از} \RL{عملیات} \RL{برداری} NumPy\RL{،} \RL{می}‌\RL{توان} \RL{سرعت} \RL{را} \RL{تا} \RL{۱۰} \RL{تا} \RL{۵۰} \RL{برابر} \RL{افزایش} \RL{داد} \RL{بدون} \RL{تغییر} \RL{در} \RL{ماهیت} \RL{الگوریتم}.
\item
  \RL{پیاده}‌\RL{سازی} GPU \RL{با} CUDA \RL{یا} OpenCL: \RL{از} \RL{آنجا} \RL{که} \RL{رمزنگاری} \RL{هر} \RL{پیکسل} \RL{مستقل} \RL{است،} \RL{این} \RL{عملیات} \RL{ذاتاً} \RL{موازی} \RL{است} \RL{و} \RL{می}‌\RL{توان} \RL{از} \RL{پردازنده}‌\RL{های} \RL{گرافیکی} \RL{برای} \RL{رمزنگاری} \RL{بلادرنگ} \RL{تصاویر} \RL{با} \RL{وضوح} \RL{بالا} \RL{بهره} \RL{برد}.
\item
  \RL{پیاده}‌\RL{سازی} \RL{در} C++ \RL{یا} Cython\RL{:} \RL{کامپایل} \RL{کد} Python \RL{با} Cython \RL{یا} \RL{بازنویسی} \RL{ماژول}‌\RL{های} \RL{بحرانی} \RL{در} C++ \RL{می}‌\RL{تواند} \RL{کارایی} \RL{را} \RL{به}‌\RL{طور} \RL{قابل} \RL{توجهی} \RL{ارتقا} \RL{دهد}.

  \begin{enumerate}
  \def\labelenumi{\arabic{enumi}.}
  \item
    \protect\phantomsection\label{_Toc232899939}{}\RL{بهبود} \RL{مرحله} \RL{انتشار} \RL{برای} \RL{ارتقاء} UACI
  \end{enumerate}
\end{itemize}

\RL{برای} \RL{رفع} \RL{انحراف} UACI \RL{به}‌\RL{ویژه} \RL{در} \RL{تصاویر} \RL{با} \RL{محتوای} \RL{یکنواخت} \RL{پیشنهادات} \RL{زیر} \RL{مطرح} \RL{می}‌\RL{شوند}:

\begin{itemize}
\item
  \RL{افزودن} \RL{یک} \RL{مرحله} \RL{انتشار} \RL{زنجیره}‌\RL{ای}\footnote{CBC-like}\RL{:} \RL{در} \RL{این} \RL{رویکرد،} \RL{خروجی} \RL{رمزنگاری} \RL{پیکسل} k \RL{به} \RL{عنوان} \RL{ورودی} \RL{اضافی} \RL{برای} \RL{رمزنگاری} \RL{پیکسل} k+1 \RL{استفاده} \RL{می}‌\RL{شود} \RL{که} \RL{وابستگی} \RL{میان} \RL{پیکسل}‌\RL{های} \RL{متوالی} \RL{را} \RL{افزایش} \RL{می}‌\RL{دهد}.
\item
  \RL{استفاده} \RL{از} \RL{جمع} \RL{ماژولار} \RL{به}‌\RL{جای} XOR: \RL{ترکیب} XOR \RL{با} \RL{جمع} \RL{ماژولار} mod 256 \RL{می}‌\RL{تواند} \RL{شدت} \RL{تغییرات} \RL{را} \RL{یکنواخت}‌\RL{تر} \RL{کرده} \RL{و} UACI \RL{را} \RL{به} \RL{مقدار} \RL{ایده}‌\RL{آل} \RL{نزدیک}‌\RL{تر} \RL{کند}.

  \begin{enumerate}
  \def\labelenumi{\arabic{enumi}.}
  \item
    \protect\phantomsection\label{_Toc232899940}{}\RL{گسترش} \RL{و} \RL{ارزیابی} \RL{بیشتر}
  \end{enumerate}
\end{itemize}

\RL{پیشنهادات} \RL{مرتبط} \RL{با} \RL{گسترش} \RL{دامنه} \RL{ارزیابی} \RL{و} \RL{کاربرد} \RL{الگوریتم}:

\begin{itemize}
\item
  \RL{اجرای} \RL{آزمون}‌\RL{های} NIST SP 800-22 \RL{برای} \RL{تأیید} \RL{رسمی} \RL{تصادفی}‌\RL{بودن} \RL{دنباله}‌\RL{های} \RL{تولیدشده} \RL{توسط} \RL{سیستم}‌\RL{های} \RL{نمایی}.
\item
  \RL{ارزیابی} \RL{مقاومت} \RL{در} \RL{برابر} \RL{حملات} \RL{متن} \RL{انتخابی}\footnote{CPA} \RL{و} \RL{حمله} \RL{متن} \RL{رمزشده} \RL{انتخابی}\footnote{CCA} \RL{که} \RL{از} \RL{روش}‌\RL{های} \RL{تحلیل} \RL{رمز} \RL{پیشرفته} \RL{محسوب} \RL{می}‌\RL{شوند}.
\item
  \RL{آزمایش} \RL{الگوریتم} \RL{بر} \RL{روی} \RL{مجموعه} \RL{داده}‌\RL{های} \RL{بزرگ}‌\RL{تر} \RL{مانند} ImageNet \RL{و} COCO \RL{برای} \RL{ارزیابی} \RL{تعمیم}‌\RL{پذیری} \RL{و} \RL{پایداری} \RL{آماری} \RL{نتایج} \RL{در} \RL{مقیاس} \RL{وسیع}.
\item
  \RL{توسعه} \RL{روش} \RL{برای} \RL{رمزنگاری} \RL{تصاویر} \RL{فشرده}‌\RL{شد} \RL{بدون} \RL{نیاز} \RL{به} \RL{رمزگشایی} \RL{ابتدا} \RL{که} \RL{کاربرد} \RL{در} \RL{شبکه}‌\RL{های} \RL{اجتماعی} \RL{و} \RL{پلتفرم}‌\RL{های} \RL{ابری} \RL{را} \RL{تسهیل} \RL{می}‌\RL{کند}.
\item
  \RL{بررسی} \RL{امکان} \RL{ادغام} \RL{الگوریتم} \RL{با} \RL{روش}‌\RL{های} \RL{واترمارکینگ}\footnote{Watermarking} \RL{برای} \RL{ایجاد} \RL{سامانه}‌\RL{های} \RL{ترکیبی} \RL{حفاظت} \RL{از} \RL{مالکیت} \RL{معنوی} \RL{و} \RL{محرمانگی} \RL{تصویر}.
\item
  \RL{مطالعه} \RL{رفتار} \RL{الگوریتم} \RL{در} \RL{برابر} \RL{حملات} \RL{کانال} \RL{جانبی}\footnote{Side-Channel} \RL{مانند} \RL{تحلیل} \RL{توان} \RL{مصرفی} \RL{که} \RL{در} \RL{پیاده}‌\RL{سازی}‌\RL{های} \RL{سخت}‌\RL{افزاری} \RL{اهمیت} \RL{دارند}.

  \begin{enumerate}
  \def\labelenumi{\arabic{enumi}.}
  \item
    \protect\phantomsection\label{_Toc232899941}{}\RL{تحقیقات} \RL{نظری}
  \end{enumerate}
\end{itemize}

\begin{itemize}
\item
  \RL{تحلیل} \RL{رسمی} \RL{امنیت} \RL{اثبات}‌\RL{پذیر}\footnote{Provable Security} \RL{برای} \RL{الگوریتم} \RL{پیشنهادی} \RL{بر} \RL{اساس} \RL{مدل}‌\RL{های} \RL{استاندارد} \RL{رمزنگاری} \RL{نظیر} IND-CPA \RL{و} IND-CCA.
\item
  \RL{بررسی} \RL{امکان} \RL{گسترش} \RL{معماری} \RL{پیشنهادی} \RL{به} \RL{رمزنگاری} \RL{تصاویر} \RL{هایپراسپکترال}\footnote{Hyperspectral} \RL{که} \RL{دارای} \RL{ده}‌\RL{ها} \RL{کانال} \RL{رنگی} \RL{هستند} \RL{و} \RL{در} \RL{دورسنجی} \RL{و} \RL{پزشکی} \RL{کاربرد} \RL{دارند}.
\item
  \RL{مطالعه} \RL{پیوند} \RL{بین} \RL{آنتروپی} \RL{سیستم}‌\RL{های} \RL{آشوبی} \RL{نمایی} \RL{و} \RL{امنیت} \RL{اطلاعاتی} \RL{از} \RL{منظر} \RL{نظریه} \RL{اطلاعات} \RL{شانون} \RL{و} \RL{امنیت} \RL{کامل}‌\RL{پوشش}‌\RL{دهی}\footnote{Perfect Secrecy} .

  \begin{enumerate}
  \def\labelenumi{\arabic{enumi}.}
  \item
    \protect\phantomsection\label{_Toc232899942}{}\RL{نتیجه}‌\RL{گیری} \RL{نهایی}
  \end{enumerate}
\end{itemize}

\RL{این} \RL{پژوهش} \RL{نشان} \RL{داد} \RL{که} \RL{ترکیب} \RL{سیستم}‌\RL{های} \RL{آشوبی} \RL{نمایی} \RL{با} \RL{معماری} \RL{چندلایه} \RL{می}‌\RL{تواند} \RL{به}‌\RL{عنوان} \RL{پایه}‌\RL{ای} \RL{مستحکم} \RL{برای} \RL{رمزنگاری} \RL{تصاویر} \RL{رنگی} \RL{ایفای} \RL{نقش} \RL{کند.} \RL{الگوریتم} \RL{پیشنهادی} \RL{با} \RL{موفقیت} \RL{چهار} \RL{ویژگی} \RL{اصلی} \RL{رمزنگاری} \RL{تصویر} \RL{را} \RL{محقق} \RL{ساخت}: \RL{معکوس}‌\RL{پذیری} \RL{کامل،} \RL{آنتروپی} \RL{بالا،} \RL{حذف} \RL{همبستگی} \RL{مکانی} \RL{و} \RL{حساسیت} \RL{بالا} \RL{به} \RL{کلید}.

\RL{نوآوری} \RL{اصلی} \RL{این} \RL{پژوهش} \RL{انتخاب} \RL{پویا} \RL{و} \RL{مستقل} \RL{سیستم} \RL{آشوبی} \RL{نمایی} \RL{برای} \RL{هر} \RL{پیکسل} --- \RL{یک} \RL{لایه} \RL{امنیتی} \RL{منحصربه}‌\RL{فرد} \RL{ایجاد} \RL{می}‌\RL{کند} \RL{که} \RL{در} \RL{اغلب} \RL{روش}‌\RL{های} \RL{مرجع} \RL{وجود} \RL{ندارد}. \RL{این} \RL{ویژگی} \RL{موجب} \RL{می}‌\RL{شود} \RL{که} \RL{حتی} \RL{اگر} \RL{مهاجم} \RL{موفق} \RL{به} \RL{شناسایی} \RL{دنباله} \RL{رمزنگاری} \RL{یک} \RL{پیکسل} \RL{شود،} \RL{هیچ} \RL{اطلاعاتی} \RL{درباره} \RL{سیستم} \RL{استفاده}‌\RL{شده} \RL{برای} \RL{پیکسل}‌\RL{های} \RL{دیگر} \RL{به}‌\RL{دست} \RL{نمی}‌\RL{آید}.

\RL{محدودیت} \RL{اصلی} \RL{پژوهش} \RL{زمان} \RL{اجرا} \RL{یک} \RL{چالش} \RL{پیاده}‌\RL{سازی} \RL{است} \RL{نه} \RL{چالش} \RL{معماری}. \RL{با} \RL{بهینه}‌\RL{سازی}‌\RL{های} \RL{پیشنهادشده} \RL{(برداری}‌\RL{سازی} \RL{و} \RL{پیاده}‌\RL{سازی} GPU\RL{)،} \RL{الگوریتم} \RL{پتانسیل} \RL{رسیدن} \RL{به} \RL{سرعت}‌\RL{های} \RL{بلادرنگ} \RL{را} \RL{داراست}. \RL{بهبود} UACI \RL{نیز} \RL{با} \RL{افزودن} \RL{مرحله} \RL{انتشار} \RL{زنجیره}‌\RL{ای} \RL{قابل} \RL{دستیابی} \RL{است}.

\RL{در} \RL{مجموع،} \RL{این} \RL{پژوهش} \RL{گامی} \RL{موثر} \RL{در} \RL{توسعه} \RL{دانش} \RL{رمزنگاری} \RL{آشوبی} \RL{برداشته} \RL{و} \RL{نشان} \RL{داده} \RL{است} \RL{که} \RL{سیستم}‌\RL{های} \RL{آشوبی} \RL{نمایی} \RL{نسل} \RL{جدید} \RL{با} \RL{پایداری} \RL{آشوب} \RL{بالاتر} \RL{و} \RL{دامنه} \RL{پارامتری} \RL{گسترده}‌\RL{تر} \RL{می}‌\RL{توانند} \RL{جایگزین} \RL{مناسبی} \RL{برای} \RL{سیستم}‌\RL{های} \RL{کلاسیک} \RL{مانند} Logistic \RL{و} Tent \RL{باشند. ترکیب} \RL{این} \RL{سیستم}‌\RL{ها} \RL{با} \RL{سیستم}‌\RL{های} \RL{آشوبی} \RL{چندبعدی} \RL{مانند} \RL{}Chen\RL{،} \RL{افقی} \RL{نوین} \RL{در} \RL{طراحی} \RL{رمزنگاری}‌\RL{های} \RL{چندلایه} \RL{مقاوم} \RL{ترسیم} \RL{می}‌\RL{کند} \RL{که} \RL{پاسخگوی} \RL{نیازهای} \RL{امنیت} \RL{دیجیتال} \RL{در} \RL{عصر} \RL{حاضر} \RL{است}.

\RL{\hfill\break
}

\subsubsection{\texorpdfstring{\RL{منابع}}{منابع}}\label{ux645ux646ux627ux628ux639}

\begin{enumerate}
\def\labelenumi{\arabic{enumi}.}
\item
  Yin, H., Xu, Y., Zhang, Y., \& Wu, J. (2024). Image compression encryption algorithm combining two-dimensional modular hyperchaotic map and compressed sensing. Physica Scripta, 99(10), 105288.
\item
  Wang, X. Y., Zhang, Y. Q., \& Zhao, Y. Y. (2015). A novel image encryption scheme based on 2-D logistic map and DNA sequence operations. Nonlinear Dynamics, 82, 1269-1280.
\item
  Zhang, B., \& Liu, L. (2023). Chaos-based image encryption: Review, application, and challenges. Mathematics, 11(11), 2585\RL{.}
\item
  Hua, Z., \& Zhou, Y. (2019). Exponential chaotic model for generating robust chaos. IEEE transactions on systems, man, and cybernetics: systems, 51(6), 3713-3724.
\item
  Li, L. (2024). A novel chaotic map application in image encryption algorithm. Expert Systems with Applications, 252(Part B), 124316.
\item
  Alvarez, M Shakir, D. A. Q., Salim, A., Abd Al-Rahman, S. Q., \& Sagheer, A. M. (2023). Image encryption using lorenz chaotic system. Journal of Techniques, 5(1), 122-128.
\item
  Khodadadi, H., \& Zandvakili, A. (2019). A New Method for Encryption of Color Images based on Combination of Chaotic Systems. Journal of AI and Data Mining, 7(3), 377-383.
\item
  Wang, X., \& Li, Y. (2021). Chaotic image encryption algorithm based on hybrid multi-objective particle swarm optimization and DNA sequence. Optics and Lasers in Engineering, 137, 106393.
\item
  Mamat, A. R., Mohamed, M. A., Abidin, A. F. A., Mohamed, R. R., Sambas, A., Rusyn, V., ... \& Markovych, B. M. (2024). Color image encryption using chaotic-based cryptosystem. \emph{Mathematical Modeling and Computing}, \emph{11}(3).
\item
  Li, Y., \RL{}Wang, X., \& Jin, J. (2020). A new one-dimensional chaotic system with applications in image encryption. Chaos, Solitons \& Fractals, 139, 110102.
\item
  \RL{صادقیان، سعیده،1400، رمزنگاری تصویر مبتنی بر آشوب با ترکیب الگوریتم بهینه سازی ذرات و الگوریتم ژنتیک، هشتمین کنگره ملی تازه های مهندسی برق و کامپیوتر ایران، تهران.}
\item
  Ge, B., Chen, X., Chen, G., \& Shen, Z. (2021). Secure and fast image encryption algorithm using hyper-chaos-based key generator and vector operation. IEEE Access, 9, 137635-137654.
\item
  Hosny, K., El‑Nabawy, Y., Elshewey, A. M., Alhammad, S., et al. (2024). New method of colour image encryption using triple chaotic maps. IET Image Processing, 18(12), 3262--3276.
\item
  Chen, Dongming, Zhiliang Zhu, and Guangming Yang. "An improved image encryption algorithm based on chaos." In The 9th International Conference for Young Computer Scientists, pp. 2792-2796. IEEE, 2008.
\item
  Lawande, Q. V., B. R. Ivan, and S. D. Dhodapkar. "Chaos based cryptography: a new approach to secure communications." BARC newsletter 258, no. 258 (2005).
\item
  Čelikovský, Sergej, and Guanrong Chen. "On the generalized Lorenz canonical form." Chaos, Solitons \& Fractals 26, no. 5 (2005): 1271-1276.
\item
  Zhou, Tianshou, Yun Tang, and Guanrong Chen. "Chen\textquotesingle s attractor exists." International Journal of Bifurcation and Chaos 14, no. 09 (2004): 3167-3177.
\item
  H.Y. Hsiao, "A Study of Chaotic Image Encryption." Master of Science Thesis, Institute of Electrical Engineering, Chung Yuan Christian University, Taiwan, 2005.
\item
  Li, Shujun. "When chaos meets computers." arXiv preprint nlin/0405038 (2004).
\item
  Strogatz, Steven H., and Donald E. Herbert. "Nonlinear dynamics and chaos." Medical Physics-New York-Institute of Physics 23, no. 6 (1996): 993-995
\item
  Zarei, B., \& Meybodi, M. R. (2021). Improving learning ability of learning automata using chaos theory. The Journal of Supercomputing, 77(1), 652-678.
\item
  Stojanovski, Toni, Johnny Pihl, and Ljupco Kocarev. "Chaos-based random number generators. Part II: practical realization." IEEE Transactions on Circuits and Systems I: Fundamental Theory and Applications 48, no. 3 (2001): 382-385
\item
  S. Li, Chen G, Cheung A, Bhargav B, Lo K-T. On the design of perceptual MPEGVideo encryption algorithms. IEEE Trans Circuits Syst Video Technol 2007;17 (2):214-23
\item
  Lorenz, Edward N. "Deterministic nonperiodic flow." Journal of the atmospheric sciences 20, no. 2 (1963): 130-141.
\item
  Mandelbrot, Benoit B. The fractal geometry of nature. Vol. 1. New York: WH freeman,1977.
\item
  Feigenbaum, Mitchell J. "Quantitative universality for a class of nonlinear transformations." Journal of statistical physics 19, no. 1 (1978): 25-52.
\item
  \RL{الوانی، سید مهدی و دانایی فرد، حسن. تئوری نظم در بینظمی و مدیریت. تهران- انتشارت صفار.} 1384
\item
  MR. Abuturab. An asymmetric single-channel color image encryption based on Hartley transform and gyrator transform. Opt Lasers Eng 2015;69:49-57
\item
  JX. Chen, Zhu ZL, Liu Z, Fu C, Zhang LB, Yu H. A novel double-image encryption scheme based on cross-image pixel scrambling in gyrator domains. Opt Express 2014;22(6):7349-61
\item
  M. Joshi, Shakher C, Singh K. Image encryption and decryption using fractional Fourier transform and radial Hilbert transform. Opt Lasers Eng 2008;46(7):522-6
\item
  ZJ. Liu, Dai JM, Sun XG, Liu ST. Color image encryption by using the rotation of color vector in Hartley transform domains. Opt Lasers Eng 2010;48(7):800-5
\item
  TH. Chen, Tsao KH, Lee YS. Yet another multiple-image encryption by rotating random grids. Signal Process 2012;92(9):2229-37
\item
  R. Enayatifar, Sadaei HJ, Abdullah AH, Lee M, Isnin IF. A novel chaotic based image encryption using a hybrid model of deoxyribonucleic acid and cellular automata. Opt Lasers Eng 2015;71:33-41
\item
  R. Enayatifar, Sadaei HJ, Abdullah AH, Lee M, Isnin IF. A novel chaotic based image encryption using a hybrid model of deoxyribonucleic acid and cellular automata. Opt Lasers Eng 2015;71:33-41
\item
  GZ. Xiao, Lu MX, Qin L, Lai XJ. New field of cryptography: DNA cryptography. Chin Sci Bull 2006.51(12):1413-20
\item
  R. Huang, Rhee KH, Uchida S. A parallel image encryption method based on compressive sensing. Multimed Tools Appl 2014;72(1):71-93
\item
  NR. Zhou, Zhang AD, Zheng F, Gong LH. Novel image compression-encryption hybrid algorithm based on key-controlled measurement matrix in compressive sensing. Opt Laser Technol 2014;62:152-60
\item
  R. Matthews. On the derivation of a chaotic encryption algorithm. Cryptologia 1989;13(1):29-42
\item
  XY. Wang, Teng L, Qin X. A novel colour image encryption algorithm based on chaos. Signal Process 2012.92(4):1101-8.
\item
  YuanWang. Xing, Sheng-XianGu , Ying-QianZhang. Novel image encryption algorithm based on cycle shift and chaotic system. Optics and Lasers in Engineering 68 (2015) 126-134
\item
  Z. Nanrun, Shumin Pan , Shan Cheng , Zhihong Zhou. Image compression-encryption scheme based on hyper-chaotic system and 2D compressive sensing. Optics \&Laser Technology 82 (2016) 121-133
\item
  YuanWang. Xing, Ying-QianZhang , Xue-MeiBao. A novel chaotic image encryption scheme using DNA sequence operations.Optics and Lasers in Engineering 73 (2015) 53-61
\item
  Xu. Lu, Zhi Li n, Jian Li, Wei Hua. A novel bit-level image encryption algorithm based on chaotic maps. Optics and Lasers in Engineering 78 (2016) 17-25.{[}40{]} Z. Rim, Ridha Ejbali, Mourad Zaied. Image encryption based on new Beta chaotic maps. Optics and Lasers in Engineering 96 (2017) 39-49
\item
  Wenhao Liu, Kehui Sun, Congxu Zhu. A fast image encryption algorithm based on chaotic map. Optics and Lasers in Engineering 84 (2016)26-36
\item
  Li. Yueping, Chunhua Wang, Hua Chen. A hyper-chaos-based image encryption algorithm using pixel-level permutation and bit-level permutation. Optics and Lasers in Engineering 90 (2017) 238-246
\item
  W. Leyuan, Hongjun Song, Ping Liu. A novel hybrid color image encryption algorithm using two complex chaotic systems. Optics and Lasers in Engineering 77 (2016) 118-125.
\end{enumerate}

\textbf{\RL{پیوست:}}

\textbf{\RL{کد اجرایی الگوریتم (پایتون)}}

import numpy as np

from scipy.integrate import solve\_ivp

import cv2

import time

import pickle

import os

import matplotlib

import matplotlib.pyplot as plt

\# Ensure output directory exists

OUT\_DIR = "outputs"

os.makedirs(OUT\_DIR, exist\_ok=True)

OUT\_FIGS = os.path.join(OUT\_DIR, "figs")

os.makedirs(OUT\_FIGS, exist\_ok=True)

print("Environment initialized. Outputs will be saved to \textquotesingle./outputs/\textquotesingle")

\# \#\#\# توضیحات خروجی

\# \textless div dir="rtl"\textgreater{}

\# کتابخانه‌های مورد نیاز بارگذاری شده و ساختار پوشه‌های خروجی ایجاد گردید.

\# \textless/div\textgreater{}

\# \#\# فاز ۱: تعریف سیستم‌های آشوبی

\# \textless div dir="rtl"\textgreater{}

\# تعریف سیستم آشوبی سه بعدی چن (Chen) برای تولید دنباله‌های جایگشت مکانی پیکسل‌ها و نقشه‌های آشوبی پایه (Logistic, Sine, Tent).

\# \textless/div\textgreater{}

\# In{[}2{]}:

\# 1. سیستم آشوبی Chen

def chen\_system(t, state, a=35, b=3, c=28):

x, y, z = state

dx = a * (y - x)

dy = (c - a) * x - x * z + c * y

dz = x * y - b * z

return {[}dx, dy, dz{]}

def generate\_chen\_sequence(x0, y0, z0, n, warmup=1000, dt=0.002):

"""تولید دنباله آشوبی Chen با حل عددی RK45"""

total = warmup + n

t\_span = (0, total * dt)

t\_eval = np.linspace(0, total * dt, total)

sol = solve\_ivp(chen\_system, t\_span, {[}x0, y0, z0{]},

method=\textquotesingle RK45\textquotesingle, t\_eval=t\_eval, dense\_output=False,

rtol=1e-10, atol=1e-12)

x\_seq = sol.y{[}0{]}

return x\_seq{[}warmup:{]} \# دور انداختن warm-up

\# 2. نقشه‌های آشوبی پایه

def logistic(x, r=4.0):

return r * x * (1.0 - x)

def sine\_map(x, r=1.0):

return (r / 4.0) * np.sin(np.pi * x)

def tent\_map(x):

return 2 * x if x \textless{} 0.5 else 2 * (1.0 - x)

\# \#\#\# توضیحات خروجی

\# \textless div dir="rtl"\textgreater{}

\# توابع مربوط به حل عددی سیستم چن با استفاده از الگوریتم رانگه-کوتا (RK45) و نگاشت‌های پایه با موفقیت پیاده‌سازی شدند.

\# \textless/div\textgreater{}

\# \#\# فاز ۲: نه سیستم آشوبی نمایی (ECM)

\# \textless div dir="rtl"\textgreater{}

\# ۹ ترکیب مختلف از نقشه‌های آشوبی با استفاده از عملگر نمایی تحت عنوان Exponential Chaotic Maps تعریف می‌شوند.

\# \textless/div\textgreater{}

\# In{[}3{]}:

\# 3. نه سیستم آشوبی نمایی (ECM)

ECM\_CONFIGS = \{

1: (\textquotesingle logistic\textquotesingle, \textquotesingle logistic\textquotesingle), \# LEL

2: (\textquotesingle logistic\textquotesingle, \textquotesingle sine\textquotesingle), \# LES

3: (\textquotesingle logistic\textquotesingle, \textquotesingle tent\textquotesingle), \# LET

4: (\textquotesingle sine\textquotesingle, \textquotesingle logistic\textquotesingle), \# SEL

5: (\textquotesingle sine\textquotesingle, \textquotesingle sine\textquotesingle), \# SES

6: (\textquotesingle sine\textquotesingle, \textquotesingle tent\textquotesingle), \# SET

7: (\textquotesingle tent\textquotesingle, \textquotesingle logistic\textquotesingle), \# TEL

8: (\textquotesingle tent\textquotesingle, \textquotesingle sine\textquotesingle), \# TES

9: (\textquotesingle tent\textquotesingle, \textquotesingle tent\textquotesingle), \# TET

\}

ECM\_NAMES = \{1:\textquotesingle LEL\textquotesingle, 2:\textquotesingle LES\textquotesingle, 3:\textquotesingle LET\textquotesingle, 4:\textquotesingle SEL\textquotesingle, 5:\textquotesingle SES\textquotesingle,

6:\textquotesingle SET\textquotesingle, 7:\textquotesingle TEL\textquotesingle, 8:\textquotesingle TES\textquotesingle, 9:\textquotesingle TET\textquotesingle\}

MU\_DEFAULT = 3.8 \# پارامتر نمایی در بازه آشوب پایدار برای همه سیستم‌ها

def \_apply\_map(name, x):

if name == \textquotesingle logistic\textquotesingle:

return logistic(x)

elif name == \textquotesingle sine\textquotesingle:

return sine\_map(x)

else:

return tent\_map(x)

def ecm\_step(x, system\_idx, mu=MU\_DEFAULT):

"""یک گام از سیستم آشوبی نمایی"""

inner\_name, outer\_name = ECM\_CONFIGS{[}system\_idx{]}

inner\_val = \_apply\_map(inner\_name, x)

\# اعمال تابع نمایی و mod 1

exp\_val = np.exp(mu * inner\_val) \% 1.0

\# اجتناب از صفر دقیق

if exp\_val == 0.0:

exp\_val = 1e-10

result = \_apply\_map(outer\_name, exp\_val)

\# اطمینان از بقای در {[}0,1)

result = np.clip(result, 1e-10, 1.0 - 1e-10)

return result

def generate\_ecm\_sequence(x0, system\_idx, n, warmup=1000, mu=MU\_DEFAULT):

"""تولید دنباله بایت از سیستم نمایی"""

x = x0

for \_ in range(warmup):

x = ecm\_step(x, system\_idx, mu)

seq = np.empty(n, dtype=np.float64)

for k in range(n):

x = ecm\_step(x, system\_idx, mu)

seq{[}k{]} = x

return seq

\# \#\# فاز ۳: دنباله انتخاب‌گر و انتشار نهایی

\# \textless div dir="rtl"\textgreater{}

\# تولید بردار انتخاب سیستم‌های آشوبی و بردار انتشار نهایی برای افزایش ایمنی در مقابل حملات تفاضلی.

\# \textless/div\textgreater{}

\# In{[}4{]}:

\# 4. تولید دنباله انتخاب‌گر (Logistic → 1..9)

def generate\_selector\_sequence(x\_sel, n, warmup=1000):

x = x\_sel

for \_ in range(warmup):

x = logistic(x)

sel = np.empty(n, dtype=np.int32)

for k in range(n):

x = logistic(x)

sel{[}k{]} = int(np.floor(x * 9)) + 1

if sel{[}k{]} \textgreater{} 9:

sel{[}k{]} = 9

return sel

\# 5. تولید دنباله لایه نهایی (Logistic مستقل)

def generate\_final\_sequence(x\_final, n, warmup=1000):

x = x\_final

for \_ in range(warmup):

x = logistic(x)

seq = np.empty(n, dtype=np.uint8)

for k in range(n):

x = logistic(x)

seq{[}k{]} = int(np.floor(x * 256)) \% 256

return seq

\# \#\# فاز ۴: توابع رمزنگاری و رمزگشایی اصلی تصویر

\# \textless div dir="rtl"\textgreater{}

\# تعریف توابع رمزنگاری و رمزگشایی تصویر با استفاده از ترکیب جایگشت چن و سیستم‌های نمایی ECM.

\# \textless/div\textgreater{}

\# In{[}5{]}:

def \_chaos\_byte(val):

"""Map chaotic state in {[}0,1) to an unbiased 8-bit key."""

return int(np.floor(val * 256.0)) \% 256

def generate\_balance\_chaos(key: dict, n: int, channel: int, warmup=1000):

"""Independent chaotic sequence for histogram balancing (per channel)."""

x = (key{[}\textquotesingle x\_final\textquotesingle{]} + (channel + 1) * 0.137) \% 1.0

if x \textless= 0.0:

x = 1e-10

for \_ in range(warmup):

x = logistic(x)

seq = np.empty(n, dtype=np.float64)

for k in range(n):

x = logistic(x)

seq{[}k{]} = x

return seq

def \_uniform\_target\_counts(n: int):

base, rem = divmod(n, 256)

counts = np.full(256, base, dtype=np.int32)

counts{[}:rem{]} += 1

return counts

def balance\_histogram\_flat(flat\_uint8, chaos):

"""Force an exactly uniform histogram via chaos-ordered surplus swaps."""

flat = flat\_uint8.reshape(-1).astype(np.uint8).copy()

n = flat.size

target = \_uniform\_target\_counts(n)

hist = np.bincount(flat, minlength=256)

diff = target - hist

surplus\_idx = {[}{]}

for v in range(256):

if diff{[}v{]} \textless{} 0:

idxs = np.where(flat == v){[}0{]}

order = np.argsort(chaos{[}idxs{]}, kind=\textquotesingle mergesort\textquotesingle)

surplus\_idx.append(idxs{[}order{[}: (-diff{[}v{]}){]}{]})

need\_vals = {[}{]}

for v in range(256):

if diff{[}v{]} \textgreater{} 0:

need\_vals.extend({[}v{]} * diff{[}v{]})

need\_vals = np.asarray(need\_vals, dtype=np.uint8)

swaps = {[}{]}

if surplus\_idx:

surplus\_idx = np.concatenate(surplus\_idx)

order = np.argsort(chaos{[}surplus\_idx{]}, kind=\textquotesingle mergesort\textquotesingle)

surplus\_idx = surplus\_idx{[}order{]}

need\_vals = need\_vals{[}order{]}

for idx, new\_val in zip(surplus\_idx, need\_vals):

old\_val = int(flat{[}idx{]})

flat{[}idx{]} = new\_val

swaps.append((int(idx), old\_val))

return flat.reshape(flat\_uint8.shape), swaps

def unbalance\_histogram\_flat(flat\_uint8, swaps):

flat = flat\_uint8.reshape(-1).astype(np.uint8).copy()

for idx, old\_val in reversed(swaps):

flat{[}idx{]} = old\_val

return flat.reshape(flat\_uint8.shape)

def encrypt\_image(image\_bgr, key: dict):

"""

image\_bgr: تصویر OpenCV (H×W×3)

key: \{\textquotesingle x0\textquotesingle,\textquotesingle y0\textquotesingle,\textquotesingle z0\textquotesingle,\textquotesingle x\_sel\textquotesingle,\textquotesingle x\_ecm\textquotesingle,\textquotesingle mu\textquotesingle,\textquotesingle x\_final\textquotesingle\}

returns: encrypted image, metadata dict (perm\_inv + balance swaps)

"""

H, W = image\_bgr.shape{[}:2{]}

N = H * W

\# ── مرحله ۲: جایگشت با Chen ──────────────────────────────────────────

chen\_seq = generate\_chen\_sequence(

key{[}\textquotesingle x0\textquotesingle{]}, key{[}\textquotesingle y0\textquotesingle{]}, key{[}\textquotesingle z0\textquotesingle{]}, N, warmup=1000)

perm = np.argsort(chen\_seq) \# بردار جایگشت

perm\_inv = np.empty\_like(perm)

perm\_inv{[}perm{]} = np.arange(N) \# جایگشت معکوس

permuted = np.empty\_like(image\_bgr)

for c in range(3):

flat = image\_bgr{[}:, :, c{]}.flatten()

permuted{[}:, :, c{]} = flat{[}perm{]}.reshape(H, W)

\# ── مرحله ۳: دنباله انتخاب‌گر ────────────────────────────────────────

selector = generate\_selector\_sequence(key{[}\textquotesingle x\_sel\textquotesingle{]}, N, warmup=1000)

\# ── مرحله ۴: XOR با سیستم‌های نمایی ─────────────────────────────────

ecm\_states = \{i: key{[}\textquotesingle x\_ecm\textquotesingle{]} for i in range(1, 10)\}

for i in range(1, 10):

x = ecm\_states{[}i{]}

for \_ in range(1000):

x = ecm\_step(x, i, key{[}\textquotesingle mu\textquotesingle{]})

ecm\_states{[}i{]} = x

diff\_xored = np.empty\_like(image\_bgr)

for c in range(3):

flat\_perm = permuted{[}:, :, c{]}.flatten().astype(np.int32)

key\_bytes = np.empty(N, dtype=np.uint8)

ecm\_st = dict(ecm\_states)

for k in range(N):

sys\_idx = selector{[}k{]}

ecm\_st{[}sys\_idx{]} = ecm\_step(ecm\_st{[}sys\_idx{]}, sys\_idx, key{[}\textquotesingle mu\textquotesingle{]})

key\_bytes{[}k{]} = \_chaos\_byte(ecm\_st{[}sys\_idx{]})

diff\_xored{[}:, :, c{]} = (flat\_perm \^{} key\_bytes.astype(np.int32)).reshape(H, W).astype(np.uint8)

\# ── مرحله ۵: XOR نهایی ───────────────────────────────────────────────

final\_seq = generate\_final\_sequence(key{[}\textquotesingle x\_final\textquotesingle{]}, N, warmup=1000)

encrypted = np.empty\_like(image\_bgr)

for c in range(3):

flat = diff\_xored{[}:, :, c{]}.flatten().astype(np.int32)

encrypted{[}:, :, c{]} = (flat \^{} final\_seq.astype(np.int32)).reshape(H, W).astype(np.uint8)

\# ── مرحله ۶: یکنواخت‌سازی هیستوگرام (کاملاً مسطح) ───────────────────

balance\_swaps = {[}{]}

for c in range(3):

chaos\_bal = generate\_balance\_chaos(key, N, c, warmup=1000)

enc\_ch, swaps = balance\_histogram\_flat(encrypted{[}:, :, c{]}, chaos\_bal)

encrypted{[}:, :, c{]} = enc\_ch

balance\_swaps.append(swaps)

meta = \{\textquotesingle perm\_inv\textquotesingle: perm\_inv, \textquotesingle balance\_swaps\textquotesingle: balance\_swaps\}

return encrypted, meta

def decrypt\_image(encrypted\_bgr, key: dict, meta):

H, W = encrypted\_bgr.shape{[}:2{]}

N = H * W

perm\_inv = meta{[}\textquotesingle perm\_inv\textquotesingle{]}

balance\_swaps = meta{[}\textquotesingle balance\_swaps\textquotesingle{]}

encrypted = encrypted\_bgr.copy()

\# ── خنثی‌سازی یکنواخت‌سازی هیستوگرام ─────────────────────────────────

for c in range(3):

encrypted{[}:, :, c{]} = unbalance\_histogram\_flat(

encrypted{[}:, :, c{]}, balance\_swaps{[}c{]})

\# ── خنثی‌سازی لایه نهایی ─────────────────────────────────────────────

final\_seq = generate\_final\_sequence(key{[}\textquotesingle x\_final\textquotesingle{]}, N, warmup=1000)

step4\_out = np.empty\_like(encrypted\_bgr)

for c in range(3):

flat = encrypted{[}:, :, c{]}.flatten().astype(np.int32)

step4\_out{[}:, :, c{]} = (flat \^{} final\_seq.astype(np.int32)).reshape(H, W).astype(np.uint8)

\# ── خنثی‌سازی XOR نمایی ──────────────────────────────────────────────

selector = generate\_selector\_sequence(key{[}\textquotesingle x\_sel\textquotesingle{]}, N, warmup=1000)

ecm\_states = \{i: key{[}\textquotesingle x\_ecm\textquotesingle{]} for i in range(1, 10)\}

for i in range(1, 10):

x = ecm\_states{[}i{]}

for \_ in range(1000):

x = ecm\_step(x, i, key{[}\textquotesingle mu\textquotesingle{]})

ecm\_states{[}i{]} = x

permuted = np.empty\_like(encrypted\_bgr)

for c in range(3):

flat\_enc = step4\_out{[}:, :, c{]}.flatten().astype(np.int32)

key\_bytes = np.empty(N, dtype=np.uint8)

ecm\_st = dict(ecm\_states)

for k in range(N):

sys\_idx = selector{[}k{]}

ecm\_st{[}sys\_idx{]} = ecm\_step(ecm\_st{[}sys\_idx{]}, sys\_idx, key{[}\textquotesingle mu\textquotesingle{]})

key\_bytes{[}k{]} = \_chaos\_byte(ecm\_st{[}sys\_idx{]})

permuted{[}:, :, c{]} = (flat\_enc \^{} key\_bytes.astype(np.int32)).reshape(H, W).astype(np.uint8)

\# ── جایگشت معکوس ─────────────────────────────────────────────────────

decrypted = np.empty\_like(encrypted\_bgr)

for c in range(3):

flat = permuted{[}:, :, c{]}.flatten()

decrypted{[}:, :, c{]} = flat{[}perm\_inv{]}.reshape(H, W)

return decrypted

\# \#\# فاز ۵: توابع ارزیابی امنیت و کارایی

\# \textless div dir="rtl"\textgreater{}

\# پیاده‌سازی معیارهای کلیدی مانند آنتروپی شانون، ضریب همبستگی پیکسل‌های همسایه، NPCR ،UACI و ارزیابی کیفیت بازیابی (PSNR/MSE).

\# \textless/div\textgreater{}

\# In{[}6{]}:

def shannon\_entropy(channel):

hist,\_ = np.histogram(channel.flatten(), bins=256, range=(0, 256))

hist = hist{[}hist \textgreater{} 0{]}.astype(np.float64)

p = hist / hist.sum()

return float(-np.sum(p * np.log2(p)))

def pixel\_correlation(channel, direction=\textquotesingle horizontal\textquotesingle, n\_samples=5000):

H, W = channel.shape

rng = np.random.default\_rng(42)

if direction == \textquotesingle horizontal\textquotesingle:

rows = rng.integers(0, H, n\_samples)

cols = rng.integers(0, W - 1, n\_samples)

x = channel{[}rows, cols{]}.astype(np.float64)

y = channel{[}rows, cols + 1{]}.astype(np.float64)

elif direction == \textquotesingle vertical\textquotesingle:

rows = rng.integers(0, H - 1, n\_samples)

cols = rng.integers(0, W, n\_samples)

x = channel{[}rows, cols{]}.astype(np.float64)

y = channel{[}rows + 1, cols{]}.astype(np.float64)

else: \# diagonal

rows = rng.integers(0, H - 1, n\_samples)

cols = rng.integers(0, W - 1, n\_samples)

x = channel{[}rows, cols{]}.astype(np.float64)

y = channel{[}rows + 1, cols + 1{]}.astype(np.float64)

corr = np.corrcoef(x, y){[}0, 1{]}

return float(corr)

def npcr\_uaci(img1, img2):

"""محاسبه NPCR و UACI بین دو تصویر رمزنگاری‌شده"""

diff = img1.astype(np.int32) - img2.astype(np.int32)

D = (diff != 0).astype(np.float64)

npcr = D.mean() *100.0

uaci = np.abs(diff).mean() / 255.0* 100.0

return float(npcr), float(uaci)

def mse\_psnr(orig, dec):

mse = np.mean((orig.astype(np.float64) - dec.astype(np.float64)) **2)

if mse == 0:

return 0.0, float(\textquotesingle inf\textquotesingle)

psnr = 10 * np.log10(255.0** 2 / mse)

return float(mse), float(psnr)

\# \#\# فاز ۶: بارگذاری تصاویر و اجرای رمزنگاری و ارزیابی

\# \textless div dir="rtl"\textgreater{}

\# بارگذاری تصاویر و اعمال فرآیند رمزنگاری و محاسبه تمامی مقادیر آماری و زمانی و ذخیره نتایج در قالب فایل pickle.

\# \textless/div\textgreater{}

\# In{[}7{]}:

\# Locate the images directory relative to the notebook\textquotesingle s directory or absolute paths

import os

img\_dir = None

for p in {[}"../../images", "images", "../images", "e:/projects/thesis\_project/images"{]}:

if os.path.isdir(p):

img\_dir = p

break

if img\_dir:

IMAGES = \{

\textquotesingle Airplane\textquotesingle: os.path.join(img\_dir, \textquotesingle Airplane.png\textquotesingle),

\textquotesingle Baboon\textquotesingle: os.path.join(img\_dir, \textquotesingle Baboon.png\textquotesingle),

\textquotesingle Peppers\textquotesingle: os.path.join(img\_dir, \textquotesingle Peppers.png\textquotesingle),

\textquotesingle Tree\textquotesingle: os.path.join(img\_dir, \textquotesingle tree.png\textquotesingle),

\}

else:

\# fallback to original paths

IMAGES = \{

\textquotesingle Airplane\textquotesingle: \textquotesingle/mnt/user-data/uploads/1781810660624\_Airplane.png\textquotesingle,

\textquotesingle Baboon\textquotesingle: \textquotesingle/mnt/user-data/uploads/1781810660626\_Baboon.png\textquotesingle,

\textquotesingle Peppers\textquotesingle: \textquotesingle/mnt/user-data/uploads/1781810660627\_Peppers.png\textquotesingle,

\textquotesingle Tree\textquotesingle: \textquotesingle/mnt/user-data/uploads/1781810660628\_tree.png\textquotesingle,

\}

print("Resolved image paths:")

for k, v in IMAGES.items():

print(f" \{k\}: \{v\}")

DEFAULT\_KEY = \{

\textquotesingle x0\textquotesingle: 0.123456789012345,

\textquotesingle y0\textquotesingle: 0.987654321098765,

\textquotesingle z0\textquotesingle: 12.345678901234,

\textquotesingle x\_sel\textquotesingle: 0.456789012345678,

\textquotesingle x\_ecm\textquotesingle: 0.789012345678901,

\textquotesingle mu\textquotesingle: 3.8,

\textquotesingle x\_final\textquotesingle: 0.234567890123456,

\}

\# Locate the images directory relative to the notebook\textquotesingle s directory or absolute paths

import os

img\_dir = None

for p in {[}"../../images", "images", "../images", "e:/projects/thesis\_project/images"{]}:

if os.path.isdir(p):

img\_dir = p

break

if img\_dir:

IMAGES = \{

\textquotesingle Airplane\textquotesingle: os.path.join(img\_dir, \textquotesingle Airplane.png\textquotesingle),

\textquotesingle Baboon\textquotesingle: os.path.join(img\_dir, \textquotesingle Baboon.png\textquotesingle),

\textquotesingle Peppers\textquotesingle: os.path.join(img\_dir, \textquotesingle Peppers.png\textquotesingle),

\textquotesingle Tree\textquotesingle: os.path.join(img\_dir, \textquotesingle tree.png\textquotesingle),

\}

else:

\# fallback to original paths

IMAGES = \{

\textquotesingle Airplane\textquotesingle: \textquotesingle/mnt/user-data/uploads/1781810660624\_Airplane.png\textquotesingle,

\textquotesingle Baboon\textquotesingle: \textquotesingle/mnt/user-data/uploads/1781810660626\_Baboon.png\textquotesingle,

\textquotesingle Peppers\textquotesingle: \textquotesingle/mnt/user-data/uploads/1781810660627\_Peppers.png\textquotesingle,

\textquotesingle Tree\textquotesingle: \textquotesingle/mnt/user-data/uploads/1781810660628\_tree.png\textquotesingle,

\}

print("Resolved image paths:")

for k, v in IMAGES.items():

print(f" \{k\}: \{v\}")

def \_hist\_uniformity(channel):

h = np.bincount(channel.flatten(), minlength=256)

return float(h.std() / h.mean() * 100.0)

results = \{\}

for name, path in IMAGES.items():

print(f"\textbackslash n\{\textquotesingle=\textquotesingle*50\}")

print(f" پردازش تصویر: \{name\}")

print(f"\{\textquotesingle=\textquotesingle*50\}")

img = cv2.imread(path)

H, W = img.shape{[}:2{]}

print(f" ابعاد: \{W\}×\{H\}")

\# ── رمزنگاری ─────────────────────────────────────────────────────────

t0 = time.perf\_counter()

enc, meta = encrypt\_image(img, DEFAULT\_KEY)

t\_enc = time.perf\_counter() - t0

print(f" زمان رمزنگاری: \{t\_enc:.3f\} ثانیه")

\# ── رمزگشایی ─────────────────────────────────────────────────────────

t0 = time.perf\_counter()

dec = decrypt\_image(enc, DEFAULT\_KEY, meta)

t\_dec = time.perf\_counter() - t0

print(f" زمان رمزگشایی: \{t\_dec:.3f\} ثانیه")

\# ── بررسی معکوس‌پذیری ────────────────────────────────────────────────

mse\_val, psnr\_val = mse\_psnr(img, dec)

print(f" MSE (اصلی vs بازیابی‌شده): \{mse\_val:.6f\}")

print(f" PSNR: \{\textquotesingle$\infty$\textquotesingle{} if psnr\_val == float(\textquotesingle inf\textquotesingle) else f\textquotesingle\{psnr\_val:.2f\} dB\textquotesingle\}")

\# ── آنتروپی ───────────────────────────────────────────────────────────

channels\_orig = {[}\textquotesingle B\textquotesingle,\textquotesingle G\textquotesingle,\textquotesingle R\textquotesingle{]}

entropy\_orig = {[}shannon\_entropy(img{[}:,:,c{]}) for c in range(3){]}

entropy\_enc = {[}shannon\_entropy(enc{[}:,:,c{]}) for c in range(3){]}

print(f" آنتروپی اصلی (B,G,R): \{{[}round(e,4) for e in entropy\_orig{]}\}")

print(f" آنتروپی رمزشده(B,G,R): \{{[}round(e,4) for e in entropy\_enc{]}\}")

\# ── همبستگی ──────────────────────────────────────────────────────────

\# کانال G (شاخص ۱ در BGR) برای گزارش اصلی

dirs = {[}\textquotesingle horizontal\textquotesingle,\textquotesingle vertical\textquotesingle,\textquotesingle diagonal\textquotesingle{]}

corr\_orig = {[}pixel\_correlation(img{[}:,:,1{]}, d) for d in dirs{]}

corr\_enc = {[}pixel\_correlation(enc{[}:,:,1{]}, d) for d in dirs{]}

print(f" همبستگی اصلی (H,V,D): \{{[}round(c,4) for c in corr\_orig{]}\}")

print(f" همبستگی رمزشده (H,V,D): \{{[}round(c,4) for c in corr\_enc{]}\}")

\# ── NPCR / UACI با تغییر یک بیت در کلید ─────────────────────────────

key2 = dict(DEFAULT\_KEY)

key2{[}\textquotesingle x0\textquotesingle{]} = DEFAULT\_KEY{[}\textquotesingle x0\textquotesingle{]} + 1e-15 \# تغییر جزئی

enc2, \_meta = encrypt\_image(img, key2)

npcr\_val, uaci\_val = npcr\_uaci(enc, enc2)

print(f" NPCR: \{npcr\_val:.4f\}\%")

print(f" UACI: \{uaci\_val:.4f\}\%")

\# ── ذخیره تصاویر ─────────────────────────────────────────────────────

cv2.imwrite(f\textquotesingle outputs/\{name\}\_encrypted.png\textquotesingle, enc)

cv2.imwrite(f\textquotesingle outputs/\{name\}\_decrypted.png\textquotesingle, dec)

hist\_cv\_enc = {[}\_hist\_uniformity(enc{[}:,:,c{]}) for c in range(3){]}

print(f" هیستوگرام رمزشده CV\% (B,G,R): \{{[}round(v,2) for v in hist\_cv\_enc{]}\}")

results{[}name{]} = \{

\textquotesingle H\textquotesingle: H, \textquotesingle W\textquotesingle: W,

\textquotesingle t\_enc\textquotesingle: t\_enc, \textquotesingle t\_dec\textquotesingle: t\_dec,

\textquotesingle mse\textquotesingle: mse\_val, \textquotesingle psnr\textquotesingle: psnr\_val,

\textquotesingle entropy\_orig\textquotesingle: entropy\_orig,

\textquotesingle entropy\_enc\textquotesingle: entropy\_enc,

\textquotesingle corr\_orig\textquotesingle: corr\_orig,

\textquotesingle corr\_enc\textquotesingle: corr\_enc,

\textquotesingle npcr\textquotesingle: npcr\_val,

\textquotesingle uaci\textquotesingle: uaci\_val,

\textquotesingle hist\_cv\_enc\textquotesingle: hist\_cv\_enc,

\textquotesingle img\_orig\textquotesingle: img,

\textquotesingle img\_enc\textquotesingle: enc,

\textquotesingle img\_dec\textquotesingle: dec,

\}

print("\textbackslash n\textbackslash n✅ تمام تصاویر پردازش شدند.")

import pickle

with open(\textquotesingle outputs/results.pkl\textquotesingle,\textquotesingle wb\textquotesingle) as f:

pickle.dump(results, f)

print("نتایج ذخیره شدند.")

\# \#\#\# توضیحات خروجی

\# \textless div dir="rtl"\textgreater{}

\# فرآیند رمزنگاری و ارزیابی معیارهای امنیت روی تصاویر با موفقیت اجرا شد و نتایج آماری در پوشه `outputs` ذخیره گردید.

\# \textless/div\textgreater{}

\# \#\# فاز ۷: رسم نمودارها و تحلیل های آماری

\# \textless div dir="rtl"\textgreater{}

\# در این بخش نمودارهای ارزیابی به صورت درون‌خطی و مجزا رسم می‌شوند.

\# \textless/div\textgreater{}

\# In{[}8{]}:

get\_ipython().run\_line\_magic(\textquotesingle matplotlib\textquotesingle, \textquotesingle inline\textquotesingle)

import matplotlib.pyplot as plt

import pickle

import numpy as np

import cv2

import os

OUT = \textquotesingle outputs/figs\textquotesingle{}

os.makedirs(OUT, exist\_ok=True)

with open(\textquotesingle outputs/results.pkl\textquotesingle, \textquotesingle rb\textquotesingle) as f:

results = pickle.load(f)

NAMES = list(results.keys())

COLORS = {[}\textquotesingle\#2E74B5\textquotesingle, \textquotesingle\#C00000\textquotesingle, \textquotesingle\#538135\textquotesingle, \textquotesingle\#7030A0\textquotesingle{]}

print(\textquotesingle Results successfully loaded. Preparing figures.\textquotesingle)

\# \#\#\# شکل ۴-۱: تصویر اصلی، تصویر رمزشده و تصویر رمزگشایی‌شده

\# \textless div dir="rtl"\textgreater{}

\# بررسی بصری خروجی اجرای الگوریتم بر روی هر چهار تصویر Airplane، Baboon، Peppers و Tree.

\# \textless/div\textgreater{}

\# In{[}9{]}:

\# شکل ۱: نمایش بصری تصاویر (اصلی، رمزشده، بازیابی‌شده)

\# ═══════════════════════════════════════════════════════════════════════════════

fig, axes = plt.subplots(4, 3, figsize=(12, 15))

fig.suptitle(\textquotesingle Figure 4-1: Original, Encrypted, and Decrypted Images\textquotesingle,

fontsize=14, fontweight=\textquotesingle bold\textquotesingle, y=0.995)

col\_titles = {[}\textquotesingle Original Image\textquotesingle, \textquotesingle Encrypted Image\textquotesingle, \textquotesingle Decrypted Image\textquotesingle{]}

for col, title in enumerate(col\_titles):

axes{[}0, col{]}.set\_title(title, fontsize=12, fontweight=\textquotesingle bold\textquotesingle, pad=8)

for row, name in enumerate(NAMES):

r = results{[}name{]}

imgs = {[}

cv2.cvtColor(r{[}\textquotesingle img\_orig\textquotesingle{]}, cv2.COLOR\_BGR2RGB),

cv2.cvtColor(r{[}\textquotesingle img\_enc\textquotesingle{]}, cv2.COLOR\_BGR2RGB),

cv2.cvtColor(r{[}\textquotesingle img\_dec\textquotesingle{]}, cv2.COLOR\_BGR2RGB),

{]}

for col, img in enumerate(imgs):

axes{[}row, col{]}.imshow(img)

axes{[}row, col{]}.axis(\textquotesingle off\textquotesingle)

if col == 0:

axes{[}row, col{]}.set\_ylabel(name, fontsize=12, fontweight=\textquotesingle bold\textquotesingle,

rotation=90, labelpad=8)

plt.tight\_layout(rect={[}0, 0, 1, 0.995{]})

plt.savefig(f\textquotesingle\{OUT\}/fig1\_visual.png\textquotesingle, bbox\_inches=\textquotesingle tight\textquotesingle)

plt.close()

print("fig1 done")

\# ═══════════════════════════════════════════════════════════════════════════════

\# \#\#\# شکل ۴-۲: مقایسه هیستوگرام تصاویر اصلی و رمزشده در هر سه کانال رنگی

\# \textless div dir="rtl"\textgreater{}

\# هیستوگرام تصاویر رمزشده باید کاملاً مسطح (یکنواخت) باشد تا هیچ اطلاعات آماری لو نرود.

\# \textless/div\textgreater{}

\# In{[}10{]}:

\# شکل ۲: هیستوگرام‌های مقایسه‌ای (اصلی vs رمزشده)

\# ═══════════════════════════════════════════════════════════════════════════════

fig, axes = plt.subplots(4, 6, figsize=(18, 11))

fig.suptitle(\textquotesingle Figure 4-2: Histogram Comparison (Original vs Encrypted) --- All Channels\textquotesingle,

fontsize=13, fontweight=\textquotesingle bold\textquotesingle)

chan\_labels = {[}\textquotesingle Blue\textquotesingle, \textquotesingle Green\textquotesingle, \textquotesingle Red\textquotesingle{]}

chan\_colors = {[}\textquotesingle\#2E74B5\textquotesingle, \textquotesingle\#538135\textquotesingle, \textquotesingle\#C00000\textquotesingle{]}

for row, name in enumerate(NAMES):

r = results{[}name{]}

for ci in range(3):

\# اصلی

ax = axes{[}row, ci{]}

ax.hist(r{[}\textquotesingle img\_orig\textquotesingle{]}{[}:,:,ci{]}.flatten(), bins=256,

range=(0,256), color=chan\_colors{[}ci{]}, alpha=0.85, density=True)

if row == 0:

ax.set\_title(f\textquotesingle Original\textbackslash n\{chan\_labels{[}ci{]}\}\textquotesingle, fontsize=10, fontweight=\textquotesingle bold\textquotesingle)

ax.set\_xlim(0,255)

ax.set\_yticks({[}{]})

if ci == 0:

ax.set\_ylabel(name, fontsize=10, fontweight=\textquotesingle bold\textquotesingle)

ax.tick\_params(labelsize=8)

\# رمزشده

ax2 = axes{[}row, ci+3{]}

ax2.hist(r{[}\textquotesingle img\_enc\textquotesingle{]}{[}:,:,ci{]}.flatten(), bins=256,

range=(0,256), color=chan\_colors{[}ci{]}, alpha=0.85, density=True)

if row == 0:

ax2.set\_title(f\textquotesingle Encrypted\textbackslash n\{chan\_labels{[}ci{]}\}\textquotesingle, fontsize=10, fontweight=\textquotesingle bold\textquotesingle)

ax2.set\_xlim(0,255)

ax2.set\_yticks({[}{]})

ax2.tick\_params(labelsize=8)

plt.tight\_layout()

plt.savefig(f\textquotesingle\{OUT\}/fig2\_histograms.png\textquotesingle, bbox\_inches=\textquotesingle tight\textquotesingle)

plt.close()

print("fig2 done")

\# ═══════════════════════════════════════════════════════════════════════════════

\# \#\#\# شکل ۴-۳: نمودار پراکندگی همبستگی پیکسل‌های همسایه در کانال سبز

\# \textless div dir="rtl"\textgreater{}

\# مقایسه همبستگی پیکسل‌های مجاور در تصویر اصلی و رمزشده در سه راستای افقی، عمودی و قطری.

\# \textless/div\textgreater{}

\# In{[}11{]}:

\# شکل ۳: ضریب همبستگی - scatter plot (اصلی vs رمزشده)

\# ═══════════════════════════════════════════════════════════════════════════════

fig, axes = plt.subplots(4, 6, figsize=(18, 11))

fig.suptitle(\textquotesingle Figure 4-3: Pixel Correlation Scatter Plots (Green Channel) --- Original vs Encrypted\textquotesingle,

fontsize=13, fontweight=\textquotesingle bold\textquotesingle)

dir\_labels\_orig = {[}\textquotesingle Original H\textquotesingle, \textquotesingle Original V\textquotesingle, \textquotesingle Original D\textquotesingle{]}

dir\_labels\_enc = {[}\textquotesingle Encrypted H\textquotesingle,\textquotesingle Encrypted V\textquotesingle,\textquotesingle Encrypted D\textquotesingle{]}

dir\_colors = {[}\textquotesingle\#2E74B5\textquotesingle,\textquotesingle\#C00000\textquotesingle,\textquotesingle\#7030A0\textquotesingle{]}

rng = np.random.default\_rng(42)

def get\_pairs(channel, direction, n=2000):

H2, W2 = channel.shape

if direction == \textquotesingle H\textquotesingle:

rs = rng.integers(0, H2, n); cs = rng.integers(0, W2-1, n)

return channel{[}rs,cs{]}, channel{[}rs,cs+1{]}

elif direction == \textquotesingle V\textquotesingle:

rs = rng.integers(0, H2-1, n); cs = rng.integers(0, W2, n)

return channel{[}rs,cs{]}, channel{[}rs+1,cs{]}

else:

rs = rng.integers(0, H2-1, n); cs = rng.integers(0, W2-1, n)

return channel{[}rs,cs{]}, channel{[}rs+1,cs+1{]}

DIRS = {[}\textquotesingle H\textquotesingle,\textquotesingle V\textquotesingle,\textquotesingle D\textquotesingle{]}

for row, name in enumerate(NAMES):

r = results{[}name{]}

g\_orig = r{[}\textquotesingle img\_orig\textquotesingle{]}{[}:,:,1{]}

g\_enc = r{[}\textquotesingle img\_enc\textquotesingle{]}{[}:,:,1{]}

for di, d in enumerate(DIRS):

\# اصلی

x1, y1 = get\_pairs(g\_orig, d)

ax = axes{[}row, di{]}

ax.scatter(x1, y1, s=1, alpha=0.3, color=dir\_colors{[}di{]})

corr\_v = r{[}\textquotesingle corr\_orig\textquotesingle{]}{[}di{]}

ax.set\_title(f\textquotesingle\{dir\_labels\_orig{[}di{]}\}\textbackslash nr=\{corr\_v:.4f\}\textquotesingle, fontsize=9)

ax.set\_xlim(0,255); ax.set\_ylim(0,255)

ax.tick\_params(labelsize=7)

if di == 0:

ax.set\_ylabel(name, fontsize=10, fontweight=\textquotesingle bold\textquotesingle)

\# رمزشده

x2, y2 = get\_pairs(g\_enc, d)

ax2 = axes{[}row, di+3{]}

ax2.scatter(x2, y2, s=1, alpha=0.3, color=dir\_colors{[}di{]})

corr\_v2 = r{[}\textquotesingle corr\_enc\textquotesingle{]}{[}di{]}

ax2.set\_title(f\textquotesingle\{dir\_labels\_enc{[}di{]}\}\textbackslash nr=\{corr\_v2:.4f\}\textquotesingle, fontsize=9)

ax2.set\_xlim(0,255); ax2.set\_ylim(0,255)

ax2.tick\_params(labelsize=7)

plt.tight\_layout()

plt.savefig(f\textquotesingle\{OUT\}/fig3\_correlation.png\textquotesingle, bbox\_inches=\textquotesingle tight\textquotesingle)

plt.close()

print("fig3 done")

\# ═══════════════════════════════════════════════════════════════════════════════

\# \#\#\# شکل ۴-۴: نمودار میله‌ای مقایسه آنتروپی شانون کانال‌های رنگی

\# \textless div dir="rtl"\textgreater{}

\# مقدار آنتروپی تصاویر رمزشده بسیار نزدیک به مقدار ایده‌آل ۸ بیت می‌باشد.

\# \textless/div\textgreater{}

\# In{[}12{]}:

\# شکل ۴: نمودار میله‌ای آنتروپی

\# ═══════════════════════════════════════════════════════════════════════════════

fig, axes = plt.subplots(1, 2, figsize=(14, 6))

fig.suptitle(\textquotesingle Figure 4-4: Shannon Entropy Comparison (Original vs Encrypted)\textquotesingle,

fontsize=13, fontweight=\textquotesingle bold\textquotesingle)

x = np.arange(len(NAMES))

w = 0.12

chan\_labels\_short = {[}\textquotesingle B\textquotesingle,\textquotesingle G\textquotesingle,\textquotesingle R\textquotesingle{]}

orig\_colors = {[}\textquotesingle\#AED6F1\textquotesingle,\textquotesingle\#A9DFBF\textquotesingle,\textquotesingle\#F1948A\textquotesingle{]}

enc\_colors = {[}\textquotesingle\#2E74B5\textquotesingle,\textquotesingle\#538135\textquotesingle,\textquotesingle\#C00000\textquotesingle{]}

for ax\_idx, (ax, mode, title) in enumerate(zip(

axes, {[}\textquotesingle orig\textquotesingle,\textquotesingle enc\textquotesingle{]},

{[}\textquotesingle Original Images\textquotesingle,\textquotesingle Encrypted Images\textquotesingle{]})):

for ci in range(3):

vals = {[}results{[}n{]}{[}f\textquotesingle entropy\_\{mode\}\textquotesingle{]}{[}ci{]} for n in NAMES{]}

bars = ax.bar(x + (ci - 1) * w, vals, width=w,

label=f\textquotesingle\{chan\_labels\_short{[}ci{]}\} channel\textquotesingle,

color=(orig\_colors if mode==\textquotesingle orig\textquotesingle{} else enc\_colors){[}ci{]},

edgecolor=\textquotesingle white\textquotesingle, linewidth=0.5)

for bar, v in zip(bars, vals):

ax.text(bar.get\_x() + bar.get\_width()/2, bar.get\_height() + 0.02,

f\textquotesingle\{v:.3f\}\textquotesingle, ha=\textquotesingle center\textquotesingle, va=\textquotesingle bottom\textquotesingle, fontsize=7.5)

ax.axhline(8.0, color=\textquotesingle black\textquotesingle, linewidth=1.5, linestyle=\textquotesingle-\/-\textquotesingle,

label=\textquotesingle Ideal H=8\textquotesingle)

ax.set\_xticks(x); ax.set\_xticklabels(NAMES, fontsize=10)

ax.set\_ylim(5.5, 8.3)

ax.set\_ylabel(\textquotesingle Shannon Entropy (bits)\textquotesingle, fontsize=11)

ax.set\_title(title, fontsize=12, fontweight=\textquotesingle bold\textquotesingle)

ax.legend(fontsize=9)

ax.grid(axis=\textquotesingle y\textquotesingle, alpha=0.3)

plt.tight\_layout()

plt.savefig(f\textquotesingle\{OUT\}/fig4\_entropy.png\textquotesingle, bbox\_inches=\textquotesingle tight\textquotesingle)

plt.close()

print("fig4 done")

\# ═══════════════════════════════════════════════════════════════════════════════

\# \#\#\# شکل ۴-۵: مقادیر معیارهای NPCR و UACI (تحلیل حساسیت به کلید)

\# \textless div dir="rtl"\textgreater{}

\# ارزیابی مقاومت الگوریتم در برابر حملات تفاضلی با مقادیر ایده‌آل.

\# \textless/div\textgreater{}

\# In{[}13{]}:

\# شکل ۵: NPCR و UACI

\# ═══════════════════════════════════════════════════════════════════════════════

fig, axes = plt.subplots(1, 2, figsize=(13, 6))

fig.suptitle(\textquotesingle Figure 4-5: NPCR and UACI --- Key Sensitivity Analysis\textquotesingle,

fontsize=13, fontweight=\textquotesingle bold\textquotesingle)

npcr\_vals = {[}results{[}n{]}{[}\textquotesingle npcr\textquotesingle{]} for n in NAMES{]}

uaci\_vals = {[}results{[}n{]}{[}\textquotesingle uaci\textquotesingle{]} for n in NAMES{]}

for ax, vals, ideal, label, title, color in zip(

axes,

{[}npcr\_vals, uaci\_vals{]},

{[}99.6, 33.46{]},

{[}\textquotesingle NPCR (\%)\textquotesingle, \textquotesingle UACI (\%)\textquotesingle{]},

{[}\textquotesingle NPCR --- Number of Pixel Change Rate\textquotesingle,

\textquotesingle UACI --- Unified Average Changing Intensity\textquotesingle{]},

{[}\textquotesingle\#2E74B5\textquotesingle,\textquotesingle\#C00000\textquotesingle{]}):

bars = ax.bar(NAMES, vals, color=color, alpha=0.85,

edgecolor=\textquotesingle white\textquotesingle, linewidth=0.5)

for bar, v in zip(bars, vals):

ax.text(bar.get\_x() + bar.get\_width()/2, bar.get\_height() + 0.3,

f\textquotesingle\{v:.2f\}\%\textquotesingle, ha=\textquotesingle center\textquotesingle, va=\textquotesingle bottom\textquotesingle, fontsize=11, fontweight=\textquotesingle bold\textquotesingle)

ax.axhline(ideal, color=\textquotesingle black\textquotesingle, linewidth=2, linestyle=\textquotesingle-\/-\textquotesingle,

label=f\textquotesingle Ideal ≈ \{ideal\}\%\textquotesingle)

y\_min = min(min(vals)*0.9, ideal*0.88)

y\_max = max(max(vals), ideal) * 1.08

ax.set\_ylim(y\_min, y\_max)

ax.set\_ylabel(label, fontsize=11)

ax.set\_title(title, fontsize=11, fontweight=\textquotesingle bold\textquotesingle)

ax.legend(fontsize=10)

ax.grid(axis=\textquotesingle y\textquotesingle, alpha=0.3)

ax.tick\_params(labelsize=10)

plt.tight\_layout()

plt.savefig(f\textquotesingle\{OUT\}/fig5\_npcr\_uaci.png\textquotesingle, bbox\_inches=\textquotesingle tight\textquotesingle)

plt.close()

print("fig5 done")

\# ═══════════════════════════════════════════════════════════════════════════════

\# \#\#\# شکل ۴-۶: زمان انجام محاسبات رمزنگاری و رمزگشایی

\# \textless div dir="rtl"\textgreater{}

\# مقایسه زمانی فرآیندها به ثانیه.

\# \textless/div\textgreater{}

\# In{[}14{]}:

\# شکل ۶: زمان اجرا

\# ═══════════════════════════════════════════════════════════════════════════════

fig, ax = plt.subplots(figsize=(10, 6))

fig.suptitle(\textquotesingle Figure 4-6: Encryption and Decryption Time (seconds)\textquotesingle,

fontsize=13, fontweight=\textquotesingle bold\textquotesingle)

t\_enc = {[}results{[}n{]}{[}\textquotesingle t\_enc\textquotesingle{]} for n in NAMES{]}

t\_dec = {[}results{[}n{]}{[}\textquotesingle t\_dec\textquotesingle{]} for n in NAMES{]}

x = np.arange(len(NAMES))

w = 0.35

b1 = ax.bar(x - w/2, t\_enc, width=w, label=\textquotesingle Encryption\textquotesingle, color=\textquotesingle\#2E74B5\textquotesingle, alpha=0.85)

b2 = ax.bar(x + w/2, t\_dec, width=w, label=\textquotesingle Decryption\textquotesingle, color=\textquotesingle\#538135\textquotesingle, alpha=0.85)

for bar, v in zip(list(b1)+list(b2), t\_enc+t\_dec):

ax.text(bar.get\_x() + bar.get\_width()/2, bar.get\_height() + 0.2,

f\textquotesingle\{v:.1f\}s\textquotesingle, ha=\textquotesingle center\textquotesingle, va=\textquotesingle bottom\textquotesingle, fontsize=10, fontweight=\textquotesingle bold\textquotesingle)

ax.set\_xticks(x); ax.set\_xticklabels(NAMES, fontsize=11)

ax.set\_ylabel(\textquotesingle Time (seconds)\textquotesingle, fontsize=11)

ax.legend(fontsize=11)

ax.grid(axis=\textquotesingle y\textquotesingle, alpha=0.3)

plt.tight\_layout()

plt.savefig(f\textquotesingle\{OUT\}/fig6\_time.png\textquotesingle, bbox\_inches=\textquotesingle tight\textquotesingle)

plt.close()

print("fig6 done")

\# ═══════════════════════════════════════════════════════════════════════════════

\# \#\#\# شکل ۴-۷: مقایسه ضرایب همبستگی عددی پیکسل‌ها در سه جهت

\# \textless div dir="rtl"\textgreater{}

\# مقایسه مستقیم مقادیر همبستگی در سه جهت افقی، عمودی و قطری.

\# \textless/div\textgreater{}

\# In{[}15{]}:

\# شکل ۷: همبستگی مقایسه‌ای (نمودار میله‌ای)

\# ═══════════════════════════════════════════════════════════════════════════════

fig, axes = plt.subplots(1, 3, figsize=(16, 6))

fig.suptitle(\textquotesingle Figure 4-7: Pixel Correlation Coefficients (Original vs Encrypted)\textquotesingle,

fontsize=13, fontweight=\textquotesingle bold\textquotesingle)

dir\_names = {[}\textquotesingle Horizontal\textquotesingle, \textquotesingle Vertical\textquotesingle, \textquotesingle Diagonal\textquotesingle{]}

for di, (ax, dname) in enumerate(zip(axes, dir\_names)):

corr\_o = {[}results{[}n{]}{[}\textquotesingle corr\_orig\textquotesingle{]}{[}di{]} for n in NAMES{]}

corr\_e = {[}results{[}n{]}{[}\textquotesingle corr\_enc\textquotesingle{]}{[}di{]} for n in NAMES{]}

x = np.arange(len(NAMES))

w = 0.35

b1 = ax.bar(x - w/2, corr\_o, width=w, label=\textquotesingle Original\textquotesingle, color=\textquotesingle\#2E74B5\textquotesingle, alpha=0.85)

b2 = ax.bar(x + w/2, corr\_e, width=w, label=\textquotesingle Encrypted\textquotesingle, color=\textquotesingle\#C00000\textquotesingle, alpha=0.85)

for bar, v in zip(list(b1)+list(b2), corr\_o+corr\_e):

offset = 0.02 if v \textgreater= 0 else -0.06

ax.text(bar.get\_x() + bar.get\_width()/2, v + offset,

f\textquotesingle\{v:.3f\}\textquotesingle, ha=\textquotesingle center\textquotesingle, va=\textquotesingle bottom\textquotesingle, fontsize=8.5, fontweight=\textquotesingle bold\textquotesingle)

ax.axhline(0, color=\textquotesingle black\textquotesingle, linewidth=1)

ax.set\_ylim(-0.2, 1.05)

ax.set\_title(f\textquotesingle\{dname\} Direction\textquotesingle, fontsize=12, fontweight=\textquotesingle bold\textquotesingle)

ax.set\_xticks(x); ax.set\_xticklabels(NAMES, fontsize=10)

ax.set\_ylabel(\textquotesingle Correlation Coefficient\textquotesingle, fontsize=10)

ax.legend(fontsize=9)

ax.grid(axis=\textquotesingle y\textquotesingle, alpha=0.3)

plt.tight\_layout()

plt.savefig(f\textquotesingle\{OUT\}/fig7\_corr\_bars.png\textquotesingle, bbox\_inches=\textquotesingle tight\textquotesingle)

plt.close()

print("fig7 done")

print("\textbackslash n✅ تمام نمودارها ذخیره شدند در:", OUT)

\# \#\#\# جدول خلاصه ارزیابی نهایی

\# \textless div dir="rtl"\textgreater{}

\# چاپ جدول آماری خلاصه از تمامی معیارهای محاسبه شده بر روی تصاویر آزمایشی.

\# \textless/div\textgreater{}

\# In{[}16{]}:

print("\textbackslash n═══ خلاصه نتایج عددی ═══")

print(f"\{\textquotesingle Image\textquotesingle:\textless10\} \{\textquotesingle H\_enc(avg)\textquotesingle:\textless14\} \{\textquotesingle NPCR\%\textquotesingle:\textless10\} \{\textquotesingle UACI\%\textquotesingle:\textless10\} \{\textquotesingle T\_enc(s)\textquotesingle:\textless12\} \{\textquotesingle T\_dec(s)\textquotesingle\}")

print("-"*65)

for n in NAMES:

r = results{[}n{]}

h\_avg = np.mean(r{[}\textquotesingle entropy\_enc\textquotesingle{]})

print(f"\{n:\textless10\} \{h\_avg:\textless14.4f\} \{r{[}\textquotesingle npcr\textquotesingle{]}:\textless10.4f\} \{r{[}\textquotesingle uaci\textquotesingle{]}:\textless10.4f\} \{r{[}\textquotesingle t\_enc\textquotesingle{]}:\textless12.3f\} \{r{[}\textquotesingle t\_dec\textquotesingle{]}:.3f\}")
